\chapter{开放量子系统的约化动力学}
\label{chap:phys-oqs-open-quantum-systems}

\section{精确约化动力学}
\label{sec:phys-oqs-exact}

上一章讨论的 Hamiltonian 在一个封闭 Hilbert 空间上生成幺正演化；一旦只保留其中一部分自由度，真正需要描述的对象便从态矢量转为约化密度算符，环境自由度通过记忆核和相关函数进入动力学。开放系统理论的起点因此仍是联合 Hilbert 空间上的精确幺正动力学，而不是某一个腔模、二能级原子或 Lindblad 方程。先把“系统”定义为直接观测的自由度，把其余自由度统称为环境；只在精确约化映射建立以后，才引入弱耦合、短记忆和旋波等条件。这样可以清楚区分哪些性质来自量子力学本身，哪些性质来自后续近似。

接下来考察一般系统--环境结构与精确约化映射。

把所有未直接观测的自由度统称为环境 $B$，其自由 Hamiltonian 为 $H_B$。系统--环境相互作用写成一般双线性形式
\begin{equation}
\boxed{
H_{SB}=\lambda H_I,
\qquad
H_I=\sum_\mu S_\mu\otimes B_\mu,}
\label{eq:phys-oqs-general-system-bath}
\end{equation}
其中 $S_\mu$ 与 $B_\mu$ 分别只作用于系统和环境 Hilbert 空间，$\lambda$ 取为无量纲的 bookkeeping 参数，只标记系统--环境耦合的微扰阶数；$H_I$ 本身保留能量量纲。总 Hamiltonian 为
\begin{equation}
\boxed{H_{\rm tot}=H_S+H_B+\lambda H_I.}
\label{eq:phys-oqs-total-h}
\end{equation}
对任意给定初始总态 $\rho_{\rm tot}(0)$，严格总演化为
\[
\rho_{\rm tot}(t)
=U_{\rm tot}(t,0)\rho_{\rm tot}(0)U_{\rm tot}^\dagger(t,0),
\]
因此系统约化态定义为
\begin{equation}
\boxed{
\rho(t)=\Tr_B\!\left[
U_{\rm tot}(t,0)\rho_{\rm tot}(0)U_{\rm tot}^\dagger(t,0)
\right].}
\label{eq:phys-oqs-exact-reduced-state}
\end{equation}
\eqnrefs{eq:phys-oqs-exact-reduced-state} 是开放系统理论的严格起点，不要求弱耦合、Markov、secular 或热平衡。

若进一步固定环境初态 $\rho_B$ 并采用无初始相关的制备
\begin{equation}
\rho_{\rm tot}(0)=\rho(0)\otimes\rho_B,
\label{eq:phys-oqs-factorized-initial-state}
\end{equation}
则可把约化演化定义为对任意系统初态都适用的映射
\begin{equation}
\boxed{
\Phi_t[\rho(0)]
=\Tr_B\!\left[
U_{\rm tot}(t,0)(\rho(0)\otimes\rho_B)U_{\rm tot}^\dagger(t,0)
\right].}
\label{eq:phys-oqs-exact-reduced-map}
\end{equation}
该映射自动保持迹、Hermiticity 和完全正性。更一般地，若制备过程只允许系统初态落在集合 $\mathfrak D_{\mathcal A}$ 中，可以引入 assignment map
\begin{equation}
\boxed{
\mathcal A:\mathfrak D_{\mathcal A}\to\mathcal S(\Hil_S\otimes\Hil_B),
\qquad
\Tr_B\mathcal A[\rho]=\rho,
}
\label{eq:phys-oqs-assignment-map}
\end{equation}
并定义
\begin{equation}
\boxed{
\Phi_{t,0}^{\mathcal A}[\rho]
=\Tr_B\!\left[U_{\rm tot}(t,0)\mathcal A[\rho]U_{\rm tot}^\dagger(t,0)\right],
\qquad \rho\in\mathfrak D_{\mathcal A}.
}
\label{eq:phys-oqs-assignment-reduced-map}
\end{equation}
乘积赋值 $\mathcal A_0[\rho]=\rho\otimes\rho_B$ 定义在整个系统态空间上，并直接给出\eqnrefs{eq:phys-oqs-exact-reduced-map} 的 CPTP 映射。若初始总态含系统--环境相关，\eqnrefs{eq:phys-oqs-exact-reduced-state} 对该给定初态仍然精确，但“只给系统初态 \(\rho(0)\) 就唯一决定总初态”的赋值通常只在一个兼容子集上定义。此时可以在这个 compatibility domain 上讨论约化映射，却一般不能把它延拓成定义在整个系统态空间上的同一个 CPTP 映射。因而是否采用因子化初态首先是一个制备与映射定义域问题，而不是 Born 弱耦合截断本身。后面的 Born 展开将在已经选定参考环境态与初始兼容条件以后，独立地作为耦合阶数近似引入。

\eqnrefs{eq:phys-oqs-exact-reduced-map} 通常不能仅用有限个系统算符写成闭合的时间局域微分方程。\secref{sec:phys-oqs-gksl}将从它对应的相互作用绘景方程出发，依次引入 Born、Markov 与 secular 条件，构造可计算的 GKSL 生成元。


接下来考察系统谱与 Bohr 频率记号。

后续所有频率分解都相对于实际采用的系统 Hamiltonian $H_S$ 进行。为建立离散 Bohr 频率的微观 GKSL 公式，以下先假设相关系统能区具有纯点谱；有限维系统自动满足这一条件。若系统含连续谱或阈值连续部分，应把下面的谱投影求和改写成投影值测度积分，并另行控制长时间极限与连续频率共振，不能直接把离散 secular 求和逐项照搬。设
\begin{equation}
H_S=\sum_{\epsilon}\epsilon\,\Pi(\epsilon),
\qquad
\Pi(\epsilon)\Pi(\epsilon')=\delta_{\epsilon\epsilon'}\Pi(\epsilon),
\label{eq:phys-oqs-system-spectral-resolution}
\end{equation}
其中 $\Pi(\epsilon)$ 是能量 $\epsilon$ 的谱投影。任意两能量块之间的带符号 Bohr 频率定义为
\[
\omega_{\epsilon'\epsilon}=\frac{\epsilon'-\epsilon}{\hbar}.
\]
若系统谱非简并并以 $|i\rangle$ 表示本征态，则记
\begin{equation}
\boxed{
\omega_{ij}=\frac{E_i-E_j}{\hbar},
\qquad
\omega_{ji}=-\omega_{ij}.
}
\label{eq:phys-oqs-bohr-signed}
\end{equation}
向下跃迁 $i\to j$ 满足 $E_i>E_j$，所以 $\omega_{ij}>0$。这一定义只固定频率符号，不要求系统是二能级或多能级原子；真正的 Bohr 算符分解将在微观生成元需要时再引入。

普通 $\otimes$ 始终表示物理 Hilbert 空间的张量积；密度算符向量化以后使用的 $\Kron$ 表示矩阵 Kronecker 积，两者不混用。

\section{从精确记忆核到 GKSL 生成元}
\label{sec:phys-oqs-gksl}

接下来考察精确 Nakajima--Zwanzig 与 Born 展开。

本节继续把系统 Hamiltonian 保持为一般的 $H_S$。只有在最后的具体模型中，才会把它替换为经过明确近似得到的有效 Hamiltonian。系统--环境总 Hamiltonian 统一写为
\[
H_{\rm tot}=H_S+H_B+\lambda H_I,
\qquad
H_I=\sum_\mu S_\mu\otimes B_\mu,
\]
其中 $\lambda$ 只用于标记系统--环境耦合的微扰阶数。定义自由部分
\[
H_0=H_S+H_B
\]
以及相互作用绘景
\[
\widetilde X(t)
=\ee^{+\ii H_0t/\hbar}X\ee^{-\ii H_0t/\hbar}.
\]
对总态，
\[
\widetilde\rho_{\rm tot}(t)
=\ee^{+\ii H_0t/\hbar}
\rho_{\rm tot}(t)
\ee^{-\ii H_0t/\hbar}.
\]
设
\[
U_S(t)=\ee^{-\ii H_St/\hbar},
\qquad
U_B(t)=\ee^{-\ii H_Bt/\hbar}.
\]
利用偏迹对环境幺正变换的不变性，
\[
\begin{aligned}
\widetilde\rho(t)
&=\Tr_B\widetilde\rho_{\rm tot}(t)\\
&=U_S^\dagger(t)
\Tr_B\!\left[U_B^\dagger(t)\rho_{\rm tot}(t)U_B(t)\right]
U_S(t)\\
&=U_S^\dagger(t)\rho(t)U_S(t).
\end{aligned}
\]
因此两种绘景下的约化态满足
\begin{equation}
\boxed{
\widetilde\rho(t)
=\ee^{+\ii H_St/\hbar}\rho(t)
 \ee^{-\ii H_St/\hbar},
\qquad
\rho(t)
=\ee^{-\ii H_St/\hbar}\widetilde\rho(t)
 \ee^{+\ii H_St/\hbar}.
}
\label{eq:phys-oqs-picture-dictionary}
\end{equation}
这个关系将在本节末端逐项用于把相互作用绘景生成元变回 Schr\"odinger 绘景。

相互作用绘景下总态严格满足
\begin{equation}
\boxed{
\dot{\widetilde\rho}_{\rm tot}(t)
=-\frac{\ii\lambda}{\hbar}
[\widetilde H_I(t),\widetilde\rho_{\rm tot}(t)].
}
\label{eq:phys-oqs-interaction-vn}
\end{equation}
在施加任何弱耦合近似之前，先把\eqnrefs{eq:phys-oqs-interaction-vn}写成精确的投影算符记忆方程。固定一个平稳参考环境态
\[
[H_B,\rho_B]=0
\]
并定义 Nakajima--Zwanzig 投影
\begin{equation}\label{eq:phys-oqs-nz-projectors}
\boxed{
\mathcal P X
=
\Tr_B(X)\otimes\rho_B,
\qquad
\mathcal Q=1-\mathcal P .
}
\end{equation}
同时把相互作用绘景 Liouvillian 定义为
\begin{equation}\label{eq:phys-oqs-interaction-liouvillian}
\mathcal L(t)X
=
-\frac{\ii\lambda}{\hbar}
[\widetilde H_I(t),X].
\end{equation}
于是精确总态方程就是
\[
\dot{\widetilde\rho}_{\rm tot}(t)
=
\mathcal L(t)\widetilde\rho_{\rm tot}(t).
\]

\begin{derivation}[精确 Nakajima--Zwanzig 记忆方程]
分别左乘 \(\mathcal P\) 和 \(\mathcal Q\)，得到两个耦合方程
\[
\frac{\dd}{\dd t}\mathcal P\widetilde\rho_{\rm tot}
=
\mathcal P\mathcal L(t)\mathcal P\widetilde\rho_{\rm tot}
+
\mathcal P\mathcal L(t)\mathcal Q\widetilde\rho_{\rm tot},
\]
\[
\frac{\dd}{\dd t}\mathcal Q\widetilde\rho_{\rm tot}
=
\mathcal Q\mathcal L(t)\mathcal P\widetilde\rho_{\rm tot}
+
\mathcal Q\mathcal L(t)\mathcal Q\widetilde\rho_{\rm tot}.
\]
定义 \(\mathcal Q\) 子空间上的传播子
\begin{equation}\label{eq:phys-oqs-q-propagator}
\mathcal G_Q(t,s)
=
\mathcal T
\exp\!\left[
\int_s^t
\mathcal Q\mathcal L(u)\mathcal Q\,\dd u
\right].
\end{equation}
第二个方程由常数变易法精确积分为
\[
\begin{aligned}
\mathcal Q\widetilde\rho_{\rm tot}(t)
={}&
\mathcal G_Q(t,0)
\mathcal Q\widetilde\rho_{\rm tot}(0)
\\
&+
\int_0^t\dd s\,
\mathcal G_Q(t,s)
\mathcal Q\mathcal L(s)\mathcal P
\widetilde\rho_{\rm tot}(s).
\end{aligned}
\]
代回 \(\mathcal P\) 方程，得到
\begin{align}
\frac{\dd}{\dd t}\mathcal P\widetilde\rho_{\rm tot}(t)
={}&
\mathcal P\mathcal L(t)\mathcal P
\widetilde\rho_{\rm tot}(t)
\nonumber\\
&+
\mathcal P\mathcal L(t)\mathcal G_Q(t,0)
\mathcal Q\widetilde\rho_{\rm tot}(0)
\nonumber\\
&+
\int_0^t\dd s\,
\mathcal K(t,s)
\mathcal P\widetilde\rho_{\rm tot}(s),
\label{eq:phys-oqs-exact-nz}
\end{align}
其中精确记忆核为
\begin{equation}\label{eq:phys-oqs-exact-memory-kernel}
\boxed{
\mathcal K(t,s)
=
\mathcal P\mathcal L(t)
\mathcal G_Q(t,s)
\mathcal Q\mathcal L(s)\mathcal P.
}
\end{equation}
\eqnrefs{eq:phys-oqs-exact-nz}没有做 Born 或 Markov 近似。第二行是初始系统--环境相关产生的 inhomogeneous term；第三行则明确显示约化动力学一般依赖全部历史状态。

\textbf{物理意义。}Nakajima--Zwanzig 方程是开放量子系统精确的约化动力学方程：不做任何近似，但把环境的影响编码进记忆核 $\mathcal K(t,s)$ 和初始相关项。记忆核的非零性意味着系统当前演化依赖全部历史——这是非 Markov 效应。在 Born 近似（弱耦合 $\lambda\ll1$）下，inhomogeneous 项消失、记忆核只依赖 $t-s$；再取 Markov 极限（环境关联时间 $\tau_c\ll$ 系统演化时间）即得 Lindblad 方程。实例：原子在非 Markov 光子浴中（如光子带隙材料内），记忆核振荡不衰减，自发辐射可被完全抑制。
\end{derivation}

若采用因子化制备
\[
\rho_{\rm tot}(0)
=
\rho(0)\otimes\rho_B,
\]
则
\[
\mathcal Q\rho_{\rm tot}(0)=0
\]
使初始相关项消失。再把环境算符中心化，使
\[
\Tr_B(B_\mu\rho_B)=0,
\]
便有
\[
\mathcal P\mathcal L(t)\mathcal P=0.
\]
此时\eqnrefs{eq:phys-oqs-exact-nz}只剩记忆积分。

由于\eqnrefs{eq:phys-oqs-interaction-liouvillian}中
\(\mathcal L=\mathcal O(\lambda)\)，
\[
\mathcal G_Q(t,s)
=
1+\mathcal O(\lambda).
\]
所以在只保留 \(\lambda^2\) 的 Born 阶数时，
\[
\mathcal K(t,s)
=
\mathcal P\mathcal L(t)
\mathcal Q\mathcal L(s)\mathcal P
+
\mathcal O(\lambda^3).
\]
在已经中心化的情况下
\(\mathcal P\mathcal L(s)\mathcal P=0\)，于是
\(\mathcal Q\mathcal L(s)\mathcal P
=\mathcal L(s)\mathcal P\)，最终得到
\begin{equation}\label{eq:phys-oqs-born-memory-from-nz}
\boxed{
\dot{\widetilde\rho}(t)
=
-\frac{\lambda^2}{\hbar^2}
\int_0^t\dd s\,
\Tr_B
\left[
\widetilde H_I(t),
\left[
\widetilde H_I(s),
\widetilde\rho(s)\otimes\rho_B
\right]
\right]
+
\mathcal O(\lambda^3).
}
\end{equation}
因此 Born 近似在这一语言中不是把精确历史依赖直接删掉，而是先把精确
\(\mathcal Q\)-子空间传播子按耦合常数展开到所需阶数；Markov 近似则是之后对
\eqnrefs{eq:phys-oqs-born-memory-from-nz}中的历史状态和积分上限作时间尺度近似。


同一精确约化动力学也可以重排为 time-local 的 TCL 展开；它与 Nakajima--Zwanzig 记忆核是两种不同的组织方式，而不是两个不同的物理模型。

在进入 Markov 极限之前，还可以从同一个二阶展开得到 time-convolutionless 形式。中心化条件使约化动力学的一阶项消失，因此
\[
\dot{\widetilde\rho}(t)=O(\lambda^2),
\qquad
\widetilde\rho(s)-\widetilde\rho(t)=O(\lambda^2)
\]
在固定有限时间区间上成立。把这个差异放回已经带 \(\lambda^2\) 的 Born 核，只改变 \(O(\lambda^4)\) 项，所以二阶精度下可把历史态换成当前态，但仍保留有限上限 \(t\)：
\begin{equation}\label{eq:phys-oqs-tcl2}
\boxed{
\dot{\widetilde\rho}(t)
=\mathcal K_{\rm TCL}^{(2)}(t)\widetilde\rho(t)+O(\lambda^3),
}
\end{equation}
其中
\begin{equation}\label{eq:phys-oqs-tcl2-generator}
\boxed{
\mathcal K_{\rm TCL}^{(2)}(t)\rho
=-\frac{\lambda^2}{\hbar^2}
\int_0^t\dd s\,
\Tr_B
\left[
\widetilde H_I(t),
\left[
\widetilde H_I(s),
\rho\otimes\rho_B
\right]
\right].
}
\end{equation}

\begin{derivation}[Born 记忆方程与 TCL2 在二阶的关系]
由式~\eqnrefs{eq:phys-oqs-born-memory-from-nz}，Born 记忆方程形如
\begin{equation*}
 \dot{\widetilde\rho}(t)
 =-\int_0^t\dd s\,\mathcal K_{\rm NZ}(t,s)\,\widetilde\rho(s),
\end{equation*}
其中 Born 核 \(\mathcal K_{\rm NZ}\) 在耦合 \(\lambda\) 的微扰展开中首阶为 \(\lambda^2\)。把历史态 \(\widetilde\rho(s)\) 在当前时刻 \(t\) 展开：
\begin{equation*}
 \widetilde\rho(s)
 =\widetilde\rho(t)-\int_s^t\dot{\widetilde\rho}(u)\dd u.
\end{equation*}
代回记忆方程，
\begin{equation*}
 \dot{\widetilde\rho}(t)
 =-\int_0^t\dd s\,\mathcal K_{\rm NZ}(t,s)\,\widetilde\rho(t)
 +\int_0^t\dd s\,\mathcal K_{\rm NZ}(t,s)\int_s^t\dd u\,\dot{\widetilde\rho}(u).
\end{equation*}
由中心化环境下一阶生成元为零，约化动力学首阶为 \(\lambda^2\)，故
\begin{equation*}
 \dot{\widetilde\rho}(u)=O(\lambda^2).
\end{equation*}
于是上式右端第二项的被积函数阶为 \(\lambda^2\cdot O(\lambda^2)=O(\lambda^4)\)，整体
\begin{equation*}
 \int_0^t\dd s\,\mathcal K_{\rm NZ}(t,s)\int_s^t\dd u\,\dot{\widetilde\rho}(u)
 =O(\lambda^4).
\end{equation*}
因此历史态可替换为
\begin{equation*}
 \widetilde\rho(s)
 =\widetilde\rho(t)+O(\lambda^2),
\end{equation*}
且 Born 核前已有显式 \(\lambda^2\)，故历史态替换产生的修正为
\begin{equation*}
 \lambda^2\,[\widetilde\rho(s)-\widetilde\rho(t)]
 =\lambda^2\cdot O(\lambda^2)=O(\lambda^4).
\end{equation*}
于是截断到二阶时，记忆方程化为
\begin{equation*}
 \dot{\widetilde\rho}(t)
 =-\int_0^t\dd s\,\mathcal K_{\rm NZ}(t,s)\,\widetilde\rho(t)+O(\lambda^4),
\end{equation*}
与式~\eqnrefs{eq:phys-oqs-tcl2-generator} 的 TCL2 生成元 \(\mathcal K_{\rm TCL}^{(2)}(t)=-\int_0^t\dd s\,\mathcal K_{\rm NZ}(t,s)\) 给出相同的约化动力学。积分上限仍为 \(t\)，因此 \(\mathcal K_{\rm TCL}^{(2)}(t)\) 一般显含初始时刻，环境相关函数的有限记忆仍保留在系数中。time-local 只是把历史依赖重新组织进时间依赖生成元，并没有自动完成 Markov 极限。

\textbf{物理意义。}TCL2 (time-convolutionless 二阶) 方程把记忆积分重新表达为时间局域方程 $\dot\rho=\mathcal K(t)\rho$，生成元 $\mathcal K(t)$ 显含时间但不含积分。这在数值上比记忆方程更高效（无需存储历史）。两种表述在二阶 Born 近似下等价，但在高阶可能不同。实例：自旋--玻色子模型在强耦合下记忆核长程振荡，TCL 展开可能发散而 NZ 记忆方程仍收敛——选择哪种表述取决于具体物理。
\end{derivation}

二阶等价并不意味着 TCL 与 NZ 在高阶逐项相同。设因子化制备下的精确约化传播映射在有限时间内可逆，并写成
\begin{equation}
\widetilde\Phi_t
=\id()+\lambda\Phi_1(t)+\lambda^2\Phi_2(t)
+\lambda^3\Phi_3(t)+\lambda^4\Phi_4(t)+\cdots .
\label{eq:phys-oqs-map-perturbative-expansion}
\end{equation}
精确 time-local 生成元定义为
\begin{equation}
\boxed{
\mathcal K_{\rm TCL}(t)=\dot{\widetilde\Phi}_t\widetilde\Phi_t^{-1},
}
\label{eq:phys-oqs-exact-tcl-generator}
\end{equation}
它只要求当前时刻的传播映射可逆，不要求动力学 Markov。中心化条件给 $\Phi_1=0$，于是逐阶相乘得到
\begin{equation}
\boxed{
\begin{aligned}
\mathcal K_{\rm TCL}(t)
={}&\lambda^2\dot\Phi_2(t)
+\lambda^3\dot\Phi_3(t)\\
&+\lambda^4\left[\dot\Phi_4(t)-\dot\Phi_2(t)\Phi_2(t)\right]
+\mathcal O(\lambda^5).
\end{aligned}}
\label{eq:phys-oqs-tcl-through-fourth-order}
\end{equation}
若环境是高斯态且耦合对环境变量为线性，奇数阶累积量消失，从而 $\Phi_3=0$；此时 TCL2 的下一项确实是 TCL4。一般非高斯环境即使已经中心化，也不能自动删去 $O(\lambda^3)$。这里的阶数是固定有限时间上的微扰阶数：截断到 $\lambda^n$ 后，局部余项为 $O(\lambda^{n+1})$，但这一估计通常不对 $t\sim O(\lambda^{-2})$ 的弛豫时间一致。若要在这样的长时间尺度上使用低阶 TCL，必须再结合弱耦合/van Hove 极限或相应重求和。另一个独立失效机制是 $\widetilde\Phi_t$ 失去可逆性；此时精确约化映射仍可良好定义，而\eqnrefs{eq:phys-oqs-exact-tcl-generator} 的 time-local 生成元可能出现奇点。

\begin{derivation}[TCL4 与 NZ 的逐阶对应]
由 $\Phi_1=0$（中心化条件使一阶累积量消失），传播映射的逆到四阶为
\[
\widetilde\Phi_t^{-1}
=\id()-\lambda^2\Phi_2-\lambda^3\Phi_3
+\lambda^4(\Phi_2^2-\Phi_4)+\mathcal O(\lambda^5).
\]
这里 \(\Phi_n\) 是 \(n\) 阶累积量传播映射，\(\id()\) 为恒等映射。该展开由 Neumann 级数 \((\id()+X)^{-1}=\id()-X+X^2-\cdots\) 逐阶得到，其中 \(X=\lambda^2\Phi_2+\lambda^3\Phi_3+\lambda^4\Phi_4+\cdots\)。具体地，\((\id()+X)^{-1}=\id()-X+X^2-X^3+\cdots\)，保留到 \(\lambda^4\) 阶给出 \(\id()-\lambda^2\Phi_2-\lambda^3\Phi_3+\lambda^4(\Phi_2^2-\Phi_4)\)。其中 \(\lambda^4\) 阶的 \(\Phi_2^2\) 来自 \(X^2\) 中 \((\lambda^2\Phi_2)^2\)，\(-\Phi_4\) 来自 \(-X\) 中 \(-\lambda^4\Phi_4\)；\(\lambda^3\) 阶的 \(-\Phi_3\) 来自 \(-X\) 中 \(-\lambda^3\Phi_3\)。再乘上
\[
\dot{\widetilde\Phi}_t
=\lambda^2\dot\Phi_2+\lambda^3\dot\Phi_3
+\lambda^4\dot\Phi_4+\mathcal O(\lambda^5),
\]
得到\eqnrefs{eq:phys-oqs-tcl-through-fourth-order}。具体地，\(\lambda^4\) 阶给出 \(\dot\Phi_4-\Phi_2\dot\Phi_2-\dot\Phi_2\Phi_2+\Phi_2^2\Phi_2\)，其中后三项在 time-local 生成元中重组为可逆性修正。展开逐阶写出：
\[
 \dot{\widetilde\Phi}_t\widetilde\Phi_t^{-1}
 =\lambda^2\dot\Phi_2
 +\lambda^3\dot\Phi_3
 +\lambda^4(\dot\Phi_4-\dot\Phi_2\Phi_2-\Phi_2\dot\Phi_2+\Phi_2^2\Phi_2)
 +\mathcal O(\lambda^5),
\]
其中 \(\lambda^2\) 阶给出 TCL2 生成元，\(\lambda^4\) 阶给出 TCL4 修正。各项物理含义：\(\dot\Phi_4\) 是四阶累积量的直接贡献；\(-\dot\Phi_2\Phi_2-\Phi_2\dot\Phi_2\) 来自传播映射逆展开中二阶项与 \(\dot{\widetilde\Phi}_t\) 中二阶项的交叉乘积；\(\Phi_2^2\Phi_2\) 来自逆展开中 \(\Phi_2^2\) 与 \(\dot\Phi_2\) 的乘积。另一方面，把 NZ 方程写成传播映射形式
\[
\dot{\widetilde\Phi}_t
=\int_0^t\dd s\,\Sigma(t,s)\widetilde\Phi_s,
\qquad
\Sigma=\lambda^2\Sigma_2+\lambda^3\Sigma_3+\lambda^4\Sigma_4+\cdots,
\]
则二阶满足
\[
\dot\Phi_2(t)=\int_0^t\Sigma_2(t,s)\dd s,
\]
该式表明二阶传播映射 \(\Phi_2\) 由二阶记忆核 \(\Sigma_2\) 的时间积分给出。四阶满足
\[
\dot\Phi_4(t)
=\int_0^t\Sigma_4(t,s)\dd s
+\int_0^t\Sigma_2(t,s)\Phi_2(s)\dd s.
\]
第二项含 \(\Phi_2(s)\)，体现了二阶过程对四阶的反馈。
代入 TCL4 后，第二项再减去 $\dot\Phi_2(t)\Phi_2(t)$，正好去掉把两个二阶过程简单串联得到的可约部分。具体地，NZ4 中 \(\int_0^t\Sigma_2(t,s)\Phi_2(s)\dd s\) 含可约部分 \(\dot\Phi_2(t)\Phi_2(t)\)（两个二阶过程独立串联），减去后剩余项为不可约四阶贡献。因而 NZ 把高阶信息组织成 memory kernel，TCL 则把同一微扰信息组织成当前时刻生成元；两者在同一截断阶给出相同传播映射，但核与生成元本身不能逐项等同。这一区别在数值实现中至关重要：TCL 是 time-local 的，每步只需当前态；NZ 则需保留历史积分，但生成元结构更简单。具体地，TCL4 生成元含可逆性修正项 \(-\Phi_2\dot\Phi_2-\dot\Phi_2\Phi_2+\Phi_2^2\Phi_2\)，这些项在 time-local 框架内完全由当前时刻的 \(\Phi_2(t)\) 及其导数决定；而 NZ4 核 \(\Sigma_4(t,s)\) 则显式依赖两个时刻 \(t,s\)，体现了记忆效应。两种方法的计算复杂度对比：
\[
 \text{TCL}: O(N_{\rm step}),\qquad
 \text{NZ}: O(N_{\rm step}^2),
\]
因为 NZ 需要在每步对历史积分。作为应用举例，在 Born--Markov 极限下，TCL2 退化为 Lindblad 量子主方程 \(\dot\rho=-\ii[H_{\rm LS},\rho]/\hbar+\sum_\alpha\gamma_\alpha(L_\alpha\rho L_\alpha^\dagger-\frac12\{L_\alpha^\dagger L_\alpha,\rho\})\)，其中 Lamb 移位 \(H_{\rm LS}\) 与耗散率 \(\gamma_\alpha\) 均由二阶核 \(\Sigma_2\) 的 Fourier 变换给出；TCL4 则提供有限时间下的非 Markov 修正，其量级由 \((\lambda/\Delta)^2\) 控制，\(\Delta\) 为 Bohr 频率间隔。Lindblad 方程的完全正定性由 Kraus 表示保证：传播映射 \(\Phi_t(\rho)=\sum_k K_k(t)\rho K_k^\dagger(t)\) 自动保持正定性与迹守恒。更一般地，TCL 与 NZ 的等价性在任意有限阶都成立：TCL\(n\) 与 NZ\(n\) 给出相同的传播映射到 \(\lambda^n\) 阶，但 TCL 生成元可能含 \(\widetilde\Phi_t^{-1}\) 的奇点（当传播映射失去可逆性时），而 NZ 核则始终良定义；这一差异在描述耗散相变等传播映射趋近不可逆的临界点时尤为显著。
\end{derivation}

接下来考察Markov 极限、时间尺度与 Redfield 方程。

若原始环境算符尚未中心化，可作
\[
B_\mu\longrightarrow
B_\mu-\Tr_B(B_\mu\rho_B)\id(B),
\]
并把产生的系统平均场项吸收到 \(H_S\) 中。因此后文统一采用
\begin{equation}
\boxed{\Tr_B(B_\mu\rho_B)=0.}
\label{eq:phys-oqs-bath-centered}
\end{equation}
在二阶 Born 展开中，\eqnrefs{eq:phys-oqs-born-memory-from-nz}等价于先把精确 $\mathcal Q$-子空间传播子按耦合常数展开，并在所保留阶数内用参考环境态
\(\rho_B\) 闭合环境相关函数；它并不声称环境在精确动力学中“永远不变”。因此 Born 的控制首先是\emph{耦合阶数控制}，而不是时间尺度近似。对中心化耦合，固定有限时间上的二阶截断写成
\[
\dot{\widetilde\rho}=\lambda^2\mathcal K_2[\widetilde\rho]
+\mathcal O(\lambda^3),
\]
若环境为高斯态且环境耦合算符线性，奇阶累积量消失，下一项才提升为 $\mathcal O(\lambda^4)$。这一固定时间的微扰余项并不自动对弛豫时间 $t\sim\tau_R=O(\lambda^{-2})$ 一致；长时间一致性需要后面的弱耦合/van Hove 极限或其他重求和。

Born 二阶闭合还要求系统对参考环境态 $\rho_B$ 的反作用在所保留阶数内可以忽略：由系统诱导的环境耗尽、相干位移或更高阶相关必须进入被截去的高阶项，而不能与二阶项同阶。有限小环境、单个窄带振子或明显回波往往不满足这一条件；这类慢自由度应并入扩大的系统，而不是继续用固定参考态 $\rho_B$ 闭合。只有在进入 Markov 步骤时，环境相关时间 $\tau_B$ 与系统弛豫时间 $\tau_R$ 的分离才成为新的独立小参数：
\begin{equation}
\boxed{
\epsilon_{\rm M}\equiv\frac{\tau_B}{\tau_R}\ll1.
}
\label{eq:phys-oqs-markov-memory-parameter}
\end{equation}

把\eqnrefs{eq:phys-oqs-born-memory-from-nz}中的高阶余项暂时省略，并令
\(\tau=t-s\)，则
\[
\dot{\widetilde\rho}(t)
=-\frac{\lambda^2}{\hbar^2}
\int_0^t\dd\tau\,
\Tr_B
[\widetilde H_I(t),
[\widetilde H_I(t-\tau),
\widetilde\rho(t-\tau)\otimes\rho_B]].
\]
若环境相关时间为 $\tau_B$，而系统约化态在这段时间内变化很小，即系统弛豫时间 $\tau_R$ 满足 $\tau_B\ll\tau_R$，就可以选择一个粗粒化时间 $\Delta t$ 使
\begin{equation}
\boxed{\tau_B\ll\Delta t\ll\tau_R.}
\label{eq:phys-oqs-markov-coarse-graining}
\end{equation}
在这样的时间窗口上，Markov 近似取
\[
\widetilde\rho(t-\tau)\simeq\widetilde\rho(t)
\quad(\tau\lesssim\tau_B),
\qquad
\int_0^t\dd\tau\longrightarrow\int_0^\infty\dd\tau.
\]
第一步丢弃系统在环境记忆时间内的变化，第二步丢弃有限初始时刻造成的边界记忆；二者都由\eqnrefs{eq:phys-oqs-markov-coarse-graining}控制，而不是独立的经验假设。
由此得到二阶 Born--Markov 核
\begin{equation}
\boxed{
\dot{\widetilde\rho}(t)
=-\frac{\lambda^2}{\hbar^2}
\int_0^\infty\dd\tau\,
\Tr_B
[\widetilde H_I(t),
[\widetilde H_I(t-\tau),
\widetilde\rho(t)\otimes\rho_B]].
}
\label{eq:phys-oqs-born-markov-kernel}
\end{equation}
若用相关函数包络 $c(\tau)$ 定义
\[
\tau_B=\frac{\int_0^\infty \tau c(\tau)\dd\tau}{\int_0^\infty c(\tau)\dd\tau},
\qquad
\epsilon_{\rm tail}(t)=
\frac{\int_t^\infty c(\tau)\dd\tau}{\int_0^\infty c(\tau)\dd\tau},
\]
则历史态替换的相对误差量级为 $O(\tau_B/\tau_R)$，把上限延长到无穷的边界误差由 $\epsilon_{\rm tail}(t)$ 控制。弱耦合下 $\tau_R^{-1}=O(\lambda^2)$，所以在固定 $\tau_B$ 的 van Hove 标度中，历史态替换对二阶生成元的首个反馈是 $O(\lambda^4)$；但 Born 截断本身在一般非高斯环境中仍可能有 $O(\lambda^3)$ 修正。环境长尾、有限模回波或 $t\lesssim\tau_B$ 时，$\epsilon_{\rm tail}$ 不小，不能使用这一 Markov 极限。

在进入 Markov 极限以后，环境相关函数最自然地由其频谱表述。

把
\[
\widetilde H_I(t)
=\sum_\mu\widetilde S_\mu(t)\otimes\widetilde B_\mu(t)
\]
代入\eqnrefs{eq:phys-oqs-born-markov-kernel}，并定义平稳环境二时相关函数
\begin{equation}
\boxed{
C_{\mu\nu}(\tau)
=\Tr_B[\widetilde B_\mu(\tau)B_\nu\rho_B].
}
\label{eq:phys-oqs-bath-correlation}
\end{equation}
由 $[H_B,\rho_B]=0$，相关函数只依赖时间差。将
\[
[A,[B,R]]=ABR-ARB-BRA+RBA
\]
逐项代入。例如第一项为
\[
\begin{aligned}
&\Tr_B\!\left[
\widetilde H_I(t)\widetilde H_I(t-\tau)
\widetilde\rho(t)\otimes\rho_B
\right]\\
&=\sum_{\mu\nu}
\widetilde S_\mu(t)\widetilde S_\nu(t-\tau)
\widetilde\rho(t)\,
C_{\mu\nu}(\tau).
\end{aligned}
\]
对另外三项作相同展开并使用偏迹循环性，得到
\begin{equation}
\boxed{
\begin{aligned}
\dot{\widetilde\rho}
=\frac{\lambda^2}{\hbar^2}
\sum_{\mu\nu}\int_0^\infty\dd\tau\Big\{&
C_{\mu\nu}(\tau)
[\widetilde S_\nu(t-\tau)\widetilde\rho\widetilde S_\mu(t)
-\widetilde S_\mu(t)\widetilde S_\nu(t-\tau)\widetilde\rho]\\
+&C_{\nu\mu}(-\tau)
[\widetilde S_\mu(t)\widetilde\rho\widetilde S_\nu(t-\tau)
-\widetilde\rho\widetilde S_\nu(t-\tau)\widetilde S_\mu(t)]
\Big\}.
\end{aligned}}
\label{eq:phys-oqs-bm-correlation-form}
\end{equation}
这一步把系统算符结构与环境统计完全分离。

定义双边谱矩阵
\begin{equation}
\boxed{
\widetilde C_{\mu\nu}(\omega)
=\int_{-\infty}^{+\infty}\dd\tau\,
\ee^{+\ii\omega\tau}C_{\mu\nu}(\tau),
}
\label{eq:phys-oqs-two-sided-spectrum}
\end{equation}
以及对应的耗散谱矩阵
\begin{equation}
\boxed{
\Gamma_{\mu\nu}(\omega)
=\frac{\lambda^2}{\hbar^2}
\widetilde C_{\mu\nu}(\omega).
}
\label{eq:phys-oqs-gamma-from-spectrum}
\end{equation}
这里 $\mu,\nu$ 是系统--环境耦合通道指标，而不是能级指标。$\Gamma(\omega)$ 可以先作为连续频率的环境响应定义；系统最终只在自己的 Bohr 频率处读取它。

下面证明\eqnrefs{eq:phys-oqs-gamma-from-spectrum} 的半正定性。对 Hermitian 的 $H_I$，总可以把耦合通道重新组合成 Hermitian 的 $S_\mu$ 与 $B_\mu$。在这一基底中，取同时对角化 $H_B$ 与平稳态 $\rho_B$ 的基底
\[
H_B|n\rangle=E_n|n\rangle,
\qquad
\rho_B=\sum_np_n|n\rangle\langle n|,
\qquad
p_n\ge0,
\]
并定义
\[
\omega_{mn}=\frac{E_m-E_n}{\hbar}.
\]
在\eqnrefs{eq:phys-oqs-bath-correlation} 中插入完备关系，得到
\[
C_{\mu\nu}(\tau)
=\sum_{nm}p_n
\ee^{-\ii\omega_{mn}\tau}
(B_\mu)_{nm}(B_\nu)_{nm}^*.
\]
Fourier 变换后
\begin{equation}
\boxed{
\Gamma_{\mu\nu}(\omega)
=\frac{2\pi\lambda^2}{\hbar^2}
\sum_{nm}p_n
(B_\mu)_{nm}(B_\nu)_{nm}^*
\delta(\omega-\omega_{mn}).
}
\label{eq:phys-oqs-gamma-spectral-representation}
\end{equation}
因此 $\Gamma_{\nu\mu}^*=\Gamma_{\mu\nu}$。再对任意复向量 $\mathbf c$ 计算
\[
\begin{aligned}
\mathbf c^\dagger\Gamma(\omega)\mathbf c
&=\frac{2\pi\lambda^2}{\hbar^2}
\sum_{nm}p_n\delta(\omega-\omega_{mn})
\left|\sum_\mu c_\mu^*(B_\mu)_{nm}\right|^2\\
&\ge0.
\end{aligned}
\]
故
\begin{equation}
\boxed{
\Gamma(\omega)=\Gamma^\dagger(\omega),
\qquad
\Gamma(\omega)\succeq0.
}
\label{eq:phys-oqs-gamma-psd}
\end{equation}
这保证后面每一个频率块都可以化为非负速率的标准 GKSL 通道。

同时把系统耦合算符按 Bohr 频率分解，才能把环境频谱与具体跃迁通道一一对应。

对当前采用的系统 Hamiltonian 作谱分解
\[
H_S=\sum_\epsilon\epsilon\Pi(\epsilon).
\]
任意系统耦合算符按 Bohr 频率分解为
\begin{equation}
\boxed{
S_\mu(\omega)
=\sum_{\epsilon'-\epsilon=\hbar\omega}
\Pi(\epsilon)S_\mu\Pi(\epsilon').
}
\label{eq:phys-oqs-bohr-component}
\end{equation}
由 $H_S\Pi(\epsilon)=\epsilon\Pi(\epsilon)$，逐项计算
\[
\begin{aligned}
[H_S,S_\mu(\omega)]
&=\sum_{\epsilon'-\epsilon=\hbar\omega}
(\epsilon-\epsilon')
\Pi(\epsilon)S_\mu\Pi(\epsilon')\\
&=-\hbar\omega S_\mu(\omega).
\end{aligned}
\]
于是
\begin{equation}
\boxed{
\widetilde S_\mu(t)
=\sum_\omega
\ee^{-\ii\omega t}S_\mu(\omega).
}
\label{eq:phys-oqs-s-interaction-bohr}
\end{equation}
在这里固定的频率约定下，$\omega>0$ 的 $S_\mu(\omega)$ 把系统从能量 $\epsilon'$ 降到较低能量 $\epsilon$；相应的 $S_\mu^\dagger(\omega)$ 描述反向跃迁。

$\omega=0$ 是 Bohr 分解的一部分，而不是另加的噪声模型。它连接等能量子空间：若存在简并，$S_\mu(0)$ 可以在简并子空间内混合状态；只有当相应零频算符在所选能量基中对角时，零频通道才表现为通常意义的纯退相干。具体地，若
\[
[Q,H_S]=0,
\qquad
Q=\sum_iq_i|i\rangle\langle i|,
\]
且该通道的零频速率记为 $\gamma_Q(0)$，则
\[
\Diss{\sqrt{\gamma_Q(0)}Q}\rho
=\gamma_Q(0)\left(Q\rho Q-\frac12\{Q^2,\rho\}\right).
\]
投影到 $|i\rangle\langle j|$，
\[
\langle i|Q\rho Q|j\rangle=q_iq_j\rho_{ij},
\qquad
\frac12\langle i|\{Q^2,\rho\}|j\rangle
=\frac{q_i^2+q_j^2}{2}\rho_{ij},
\]
所以
\begin{equation}
\boxed{
\dot\rho_{ij}\big|_\varphi
=-\frac{\gamma_Q(0)}2(q_i-q_j)^2\rho_{ij},
\qquad
\dot\rho_{ii}\big|_\varphi=0.
}
\label{eq:phys-oqs-zero-frequency-dephasing}
\end{equation}
\eqnrefs{eq:phys-oqs-zero-frequency-dephasing} 说明：对角零频噪声不改变 population，而是随机化不同能级之间的相对相位；后文的 $\gamma_{ij}^{\varphi}$ 正是这类零频通道在给定相干元上的总贡献。

把\eqnrefs{eq:phys-oqs-s-interaction-bohr} 代入\eqnrefs{eq:phys-oqs-bm-correlation-form}。例如其中一项变为
\[
\begin{aligned}
&\int_0^\infty\dd\tau\,
C_{\mu\nu}(\tau)
\widetilde S_\nu(t-\tau)\widetilde\rho\widetilde S_\mu(t)\\
&=\sum_{\omega,\omega'}
\ee^{\ii(\omega'-\omega)t}
\left[
\int_0^\infty\dd\tau\,
\ee^{+\ii\omega\tau}C_{\mu\nu}(\tau)
\right]
S_\nu(\omega)\widetilde\rho S_\mu^\dagger(\omega').
\end{aligned}
\]
因此所有记忆积分都归结为单边 Fourier 变换
\begin{equation}
\boxed{
G_{\mu\nu}(\omega)
=\int_0^\infty\dd\tau\,
\ee^{+\ii\omega\tau}C_{\mu\nu}(\tau).
}
\label{eq:phys-oqs-one-sided-response}
\end{equation}
利用分布恒等式
\[
\int_0^\infty\dd\tau\,\ee^{\ii x\tau}
=\pi\delta(x)+\ii\,\mathcal P\frac1x,
\]
得到逐矩阵元的 Plemelj 关系
\begin{equation}
\boxed{
G_{\mu\nu}(\omega)
=\frac12\widetilde C_{\mu\nu}(\omega)
+\frac{\ii}{2\pi}\mathcal P
\int_{-\infty}^{+\infty}\dd\omega'\,
\frac{\widetilde C_{\mu\nu}(\omega')}{\omega-\omega'}.
}
\label{eq:phys-oqs-plemelj}
\end{equation}
因此耗散与 Lamb shift 来自同一个复响应。一般地定义
\begin{equation}
\boxed{
\Gamma_{\mu\nu}(\omega)
=\frac{\lambda^2}{\hbar^2}
\left[G_{\mu\nu}(\omega)+G_{\nu\mu}^*(\omega)\right],
}
\label{eq:phys-oqs-gamma-from-g}
\end{equation}
和
\begin{equation}
\boxed{
\Delta_{\mu\nu}^{\rm LS}(\omega)
=\frac{\lambda^2}{2\ii\hbar}
\left[G_{\mu\nu}(\omega)-G_{\nu\mu}^*(\omega)\right].
}
\label{eq:phys-oqs-delta-ls}
\end{equation}
在前面采用的 Hermitian 环境算符基底中，
\[
C_{\nu\mu}(-\tau)=C_{\mu\nu}^*(\tau),
\]
于是\eqnrefs{eq:phys-oqs-gamma-from-g} 与\eqnrefs{eq:phys-oqs-gamma-from-spectrum} 完全一致，而
\begin{equation}
\boxed{
\Delta_{\mu\nu}^{\rm LS}(\omega)
=\frac{\lambda^2}{2\pi\hbar}
\mathcal P\int_{-\infty}^{+\infty}\dd\omega'\,
\frac{\widetilde C_{\mu\nu}(\omega')}{\omega-\omega'}.
}
\label{eq:phys-oqs-delta-hilbert}
\end{equation}
$\Gamma$ 的量纲为频率，$\Delta^{\rm LS}$ 的量纲为能量。前者只读取环境谱在系统 Bohr 频率处的在壳值；后者是整个频谱的 Hilbert 变换。

接下来考察有限粗粒化、partial secular 与 Kossakowski 矩阵。

把 Bohr 分解代入\eqnrefs{eq:phys-oqs-bm-correlation-form} 并保留所有 $\omega,\omega'$ 交叉项，就得到通常所称的二阶 Redfield 生成元。它已经是 time-local 且在 Markov 时间尺度上系数近似常数，但一般不具有 GKSL 形式；特别是交叉频率块对应的系数矩阵不必整体半正定，所以 Redfield 演化不能仅凭“二阶、time-local”就宣称完全正。完全正性要么来自受控的有限粗粒化构造，要么来自进一步的 secular 极限。

Bohr 分解代入 Born--Markov 核后，尚未 secular 化的项都带有
\[
\ee^{\ii(\omega'-\omega)t}
\]
这样的相对相位。粗粒化时间还必须足以平均掉不同 Bohr 频率之间的相对相位。若希望删除一对 $\omega\neq\omega'$ 的交叉项，需要同时存在
\[
\tau_B\ll\Delta t\ll\tau_R,
\qquad
|\omega-\omega'|\Delta t\gg1.
\]
这样的 $\Delta t$ 存在的充分尺度条件是
\begin{equation}
|\omega-\omega'|\,\tau_R\gg1.
\label{eq:phys-oqs-secular-condition}
\end{equation}
因此 secular 并不是 Born--Markov 的自动结果，而是对同一粗粒化窗口追加的频率分辨条件。完全 secular 近似于是保留 $\omega=\omega'$。若两条跃迁严格简并，或其频率差小到不满足\eqnrefs{eq:phys-oqs-secular-condition}，则它们必须保留在同一频率块内；此时 $\Gamma_{\mu\nu}(\omega)$ 的非对角元描述环境不能分辨这些跃迁所产生的干涉耗散。

这个条件还可以从有限粗粒化平均直接看出。令
\(\Delta\omega=\omega'-\omega\)，则
\begin{align}
\frac1{\Delta t}
\int_{t-\Delta t/2}^{t+\Delta t/2}
\ee^{\ii\Delta\omega s}\dd s
&=
\ee^{\ii\Delta\omega t}
\frac{
\sin(\Delta\omega\Delta t/2)
}{
\Delta\omega\Delta t/2
}
\nonumber\\
&=
\ee^{\ii\Delta\omega t}
\operatorname{sinc}
\!\left(
\frac{\Delta\omega\Delta t}{2}
\right).
\label{eq:phys-oqs-finite-coarse-grain-sinc}
\end{align}
因此完全 secular 是
\(|\Delta\omega|\Delta t\gg1\) 的极限；当一组近简并频率满足
\[
|\omega-\omega'|\Delta t\lesssim1,
\]
它们应作为同一个 coarse-grained frequency block 保留，这就是 partial-secular 的自然来源。

为在保留近简并交叉项的同时保证完全正性，不能从 Redfield 方程中任意删项，而应直接在有限窗口 $\Delta t$ 上构造二阶 coarse-grained Kossakowski 矩阵。在前面采用的 Hermitian 环境算符基底中，把复合指标记为 $a=(\mu,\omega)$，定义
\begin{equation}
\boxed{
\Gamma_{\mu\omega,\nu\omega'}^{(\Delta t)}
=\frac{\lambda^2}{\hbar^2\Delta t}
\int_0^{\Delta t}\dd s\int_0^{\Delta t}\dd u\,
\ee^{+\ii\omega s-\ii\omega'u}
C_{\mu\nu}(s-u).
}
\label{eq:phys-oqs-coarse-grained-kossakowski}
\end{equation}
对任意复系数 $z_{\mu\omega}$，令
\[
X_z=\sum_{\nu,\omega'}z_{\nu\omega'}
\int_0^{\Delta t}\ee^{-\ii\omega'u}\widetilde B_\nu(u)\dd u.
\]
平稳性给
\begin{align}
\sum_{\mu\omega,\nu\omega'}
 z_{\mu\omega}^*
\Gamma_{\mu\omega,\nu\omega'}^{(\Delta t)}
 z_{\nu\omega'}
&=\frac{\lambda^2}{\hbar^2\Delta t}
\Tr_B(X_z^\dagger X_z\rho_B)\nonumber\\
&\ge0.
\label{eq:phys-oqs-coarse-grained-gram}
\end{align}
所以
\begin{equation}
\boxed{\Gamma^{(\Delta t)}\succeq0}
\label{eq:phys-oqs-coarse-grained-psd}
\end{equation}
是一个直接的 Gram 结论，而不是额外假设。令复合指标 $a=(\mu,\omega)$ 并记 $S_a=S_\mu(\omega)$，则有限粗粒化生成元的耗散部分可直接写为
\begin{equation}
\boxed{
\mathcal D_{\Delta t}[\rho]
=\sum_{a,b}\Gamma_{ab}^{(\Delta t)}
\left(
S_b\rho S_a^\dagger
-\frac12\{S_a^\dagger S_b,\rho\}
\right).
}
\label{eq:phys-oqs-coarse-grained-dissipator}
\end{equation}
与双时间积分的 anti-Hermitian 部分产生的 $H_{\rm LS}^{(\Delta t)}$ 合并后，
\[
\mathcal L_{\Delta t}\rho
=-\frac{\ii}{\hbar}[H_S+H_{\rm LS}^{(\Delta t)},\rho]
+\mathcal D_{\Delta t}[\rho]
\]
已经是 GKSL 形式；因此这里的完全正性来自\eqnrefs{eq:phys-oqs-coarse-grained-psd}，而不是来自随后再做 complete secular。等价地，引入归一化有限时间滤波器
\[
F_{\Delta t}(x)=\frac1{\sqrt{\Delta t}}
\int_0^{\Delta t}\ee^{\ii xs}\dd s
=\sqrt{\Delta t}\,\ee^{\ii x\Delta t/2}
\operatorname{sinc}\!\left(\frac{x\Delta t}{2}\right),
\]
可把\eqnrefs{eq:phys-oqs-coarse-grained-kossakowski}写成
\begin{equation}
\Gamma_{\mu\omega,\nu\omega'}^{(\Delta t)}
=\frac{\lambda^2}{\hbar^2}
\int_{-\infty}^{+\infty}\frac{\dd\varpi}{2\pi}
\widetilde C_{\mu\nu}(\varpi)
F_{\Delta t}(\omega-\varpi)
F_{\Delta t}^*(\omega'-\varpi).
\label{eq:phys-oqs-coarse-grained-spectral-filter}
\end{equation}
因此 partial secular 的含义是：只要 $|\omega-\omega'|\Delta t\lesssim1$，两个滤波窗口仍有显著重叠，交叉项就应作为同一正半定 Kossakowski 块保留；当 $|\omega-\omega'|\Delta t\gg1$ 时，其重叠按 sinc 尾部衰减，才进入完全 secular 极限。一个保持完全正性的离散化办法，是把近简并复合指标分成若干频率簇，完整保留每个簇对应的 principal block，只把不同簇之间的整块耦合置零；所得块对角矩阵仍为半正定。反之，任意“凭频率看起来接近就删掉若干单独矩阵元”的做法都可能破坏\eqnrefs{eq:phys-oqs-coarse-grained-psd}。

\begin{note}[四个容易混淆的逻辑层级]
time-local 只表示方程可写成 $\dot\rho(t)=\mathcal K(t)\rho(t)$，其系数仍可携带有限初始时刻和环境记忆，因此 time-local 不等于 Markov。Redfield 方程来自 Born--Markov 的二阶闭合，但未必具有半正定 Kossakowski 矩阵，因此 Redfield 不等于 GKSL。有限 $\Delta t$ 的 coarse graining 保留满足 $|\omega-\omega'|\Delta t\lesssim1$ 的近简并块，是 partial secular；只有在所有待区分频率满足 $|\omega-\omega'|\Delta t\gg1$ 时，才可进入 complete secular。最后，partial secular 本身也不是“任意保留一些交叉项”：只有像\eqnrefs{eq:phys-oqs-coarse-grained-kossakowski}这样从有限时间 Gram 矩阵整体构造，才能保证完全正性。

在 $\tau_B\ll\Delta t\ll\tau_R$ 内，Markov 历史态替换的相对误差为 $O(\epsilon_{\rm M})=O(\tau_B/\tau_R)$，有限频率平均的残余交叉项由 $\operatorname{sinc}[(\omega-\omega')\Delta t/2]$ 控制；若 $|\omega-\omega'|\Delta t\gg1$，其包络至多为 $O((|\omega-\omega'|\Delta t)^{-1})$。这些误差与 Born/TCL 的耦合阶数截断是不同来源，不能用一个“小耦合”条件同时替代。
\end{note}

若不同 Bohr 频率之间的拍频在粗粒化时间内完全平均掉，就进一步得到完全 secular 的 GKSL 极限。

把\eqnrefs{eq:phys-oqs-one-sided-response} 的 Hermitian 与 anti-Hermitian 部分分别代回同一频率块。由\eqnrefs{eq:phys-oqs-gamma-from-g}--\eqnrefs{eq:phys-oqs-delta-ls}，可写成
\[
G(\omega)
=\frac{\hbar^2}{2\lambda^2}\Gamma(\omega)
+\frac{\ii\hbar}{\lambda^2}\Delta^{\rm LS}(\omega).
\]
其中 Hermitian 部分组合为耗散子，anti-Hermitian 部分组合为对易子。定义
\begin{equation}
\boxed{
\HLS
=\sum_\omega\sum_{\mu\nu}
\Delta_{\mu\nu}^{\rm LS}(\omega)
S_\mu^\dagger(\omega)S_\nu(\omega).
}
\label{eq:phys-oqs-hls-general}
\end{equation}
由
\[
[H_S,S_\nu(\omega)]=-\hbar\omega S_\nu(\omega),
\qquad
[H_S,S_\mu^\dagger(\omega)]=+\hbar\omega S_\mu^\dagger(\omega),
\]
逐项得到
\[
\begin{aligned}
[H_S,S_\mu^\dagger(\omega)S_\nu(\omega)]
={}&[H_S,S_\mu^\dagger(\omega)]S_\nu(\omega)\\
&+S_\mu^\dagger(\omega)[H_S,S_\nu(\omega)]\\
={}&0.
\end{aligned}
\]
故
\begin{equation}
\boxed{[H_S,\HLS]=0.}
\label{eq:phys-oqs-hls-commutes}
\end{equation}
secular 后相互作用绘景中的生成元为
\begin{equation}
\boxed{
\begin{aligned}
\dot{\widetilde\rho}
={}&-\frac{\ii}{\hbar}[\HLS,\widetilde\rho]\\
&+\sum_\omega\sum_{\mu\nu}
\Gamma_{\mu\nu}(\omega)
\left[
S_\nu(\omega)\widetilde\rho S_\mu^\dagger(\omega)
-\frac12\{S_\mu^\dagger(\omega)S_\nu(\omega),\widetilde\rho\}
\right].
\end{aligned}}
\label{eq:phys-oqs-gksl-interaction-picture}
\end{equation}

现在使用\eqnrefs{eq:phys-oqs-picture-dictionary} 返回 Schr\"odinger 绘景。由
\[
\rho=U_S\widetilde\rho U_S^\dagger
\]
对时间求导，
\[
\begin{aligned}
\dot\rho
&=\dot U_S\widetilde\rho U_S^\dagger
+U_S\dot{\widetilde\rho}U_S^\dagger
+U_S\widetilde\rho\dot U_S^\dagger\\
&=-\frac{\ii}{\hbar}[H_S,\rho]
+U_S\dot{\widetilde\rho}U_S^\dagger.
\end{aligned}
\]
另一方面，由\eqnrefs{eq:phys-oqs-s-interaction-bohr} 等价地得到
\[
U_SS_\nu(\omega)U_S^\dagger
=\ee^{+\ii\omega t}S_\nu(\omega),
\qquad
U_SS_\mu^\dagger(\omega)U_S^\dagger
=\ee^{-\ii\omega t}S_\mu^\dagger(\omega).
\]
因此 recycle 结构中的相位逐项抵消：
\[
\begin{aligned}
&U_S
[S_\nu(\omega)\widetilde\rho S_\mu^\dagger(\omega)]
U_S^\dagger\\
&=\ee^{+\ii\omega t}S_\nu(\omega)
\rho
\ee^{-\ii\omega t}S_\mu^\dagger(\omega)\\
&=S_\nu(\omega)\rho S_\mu^\dagger(\omega).
\end{aligned}
\]
同样，$S_\mu^\dagger S_\nu$ 中的相位也抵消；再由\eqnrefs{eq:phys-oqs-hls-commutes}，
\[
U_S\HLS U_S^\dagger=\HLS.
\]
于是得到 Schr\"odinger 绘景中的微观 GKSL 方程
\begin{equation}
\boxed{
\begin{aligned}
\dot\rho
={}&-\frac{\ii}{\hbar}[H_S+\HLS,\rho]\\
&+\sum_\omega\sum_{\mu\nu}
\Gamma_{\mu\nu}(\omega)
\left[
S_\nu(\omega)\rho S_\mu^\dagger(\omega)
-\frac12\{S_\mu^\dagger(\omega)S_\nu(\omega),\rho\}
\right].
\end{aligned}}
\label{eq:phys-oqs-gksl-microscopic}
\end{equation}

由\eqnrefs{eq:phys-oqs-gamma-psd}，每个 $\omega$ 块都可以酉对角化：
\[
\Gamma(\omega)
=U^\dagger(\omega)
\operatorname{diag}[\gamma_1(\omega),\gamma_2(\omega),\ldots]
U(\omega),
\qquad
\gamma_q(\omega)\ge0.
\]
定义含速率的 collapse operator
\begin{equation}
\boxed{
C_q(\omega)
=\sqrt{\gamma_q(\omega)}
\sum_\mu U_{q\mu}(\omega)S_\mu(\omega).
}
\label{eq:phys-oqs-collapse-from-gamma}
\end{equation}
于是\eqnrefs{eq:phys-oqs-gksl-microscopic} 化为标准形式
\begin{equation}
\boxed{
\dot\rho
=-\frac{\ii}{\hbar}[H,\rho]
+\sum_{\omega,q}\Diss{C_q(\omega)}\rho,
\qquad
H\equiv H_S+\HLS,
}
\label{eq:phys-oqs-gksl-standard}
\end{equation}
其中
\begin{equation}
\boxed{
\Diss{C}\rho
=C\rho C^\dagger-\frac12\{C^\dagger C,\rho\}.
}
\label{eq:phys-oqs-dissipator}
\end{equation}
正半定性\eqnrefs{eq:phys-oqs-gamma-psd} 在这里承担了关键作用：它保证所有 $\gamma_q(\omega)$ 非负，因此耗散部分可以写成一组真实非负速率的 Lindblad 通道。反对易项使 $\Tr\rho$ 的一阶变化精确抵消 recycle 项，而非负速率又使生成元满足 GKSL 定理的完全正性条件。后文\eqnrefs{eq:phys-oqs-infinitesimal-completeness} 会进一步把这个结论写成显式的无穷小 Kraus 完备关系；因此“环境谱矩阵半正定 $\Rightarrow$ 生成元完全正”不是额外假设，而是从平稳环境相关函数到 CPTP 半群的闭合链条。

接下来考察热平衡环境：KMS 关系、detailed balance 与 Gibbs 稳态。

若环境不是任意平稳态，而是温度 $T=(k_B\beta)^{-1}$ 的 Gibbs 态
\begin{equation}
\rho_B^{(\beta)}=\frac{\ee^{-\beta H_B}}{Z_B},
\qquad
Z_B=\Tr_B\ee^{-\beta H_B},
\label{eq:phys-oqs-bath-gibbs-state}
\end{equation}
则平稳性之外还多出 Kubo--Martin--Schwinger（KMS）约束。沿用\eqnrefs{eq:phys-oqs-gamma-spectral-representation}的 Hermitian 耦合基底，在能量本征态上有
\[
p_n=\frac{\ee^{-\beta E_n}}{Z_B}.
\]
把 $n,m$ 互换，并利用 $\delta(\omega-\omega_{mn})$ 强制 $E_m-E_n=\hbar\omega$，可得谱矩阵的精确 detailed-balance 关系
\begin{equation}
\boxed{
\Gamma_{\mu\nu}(-\omega)
=\ee^{-\beta\hbar\omega}\,
\Gamma_{\nu\mu}(\omega).
}
\label{eq:phys-oqs-kms-detailed-balance}
\end{equation}
因此正频率的向下跃迁与负频率的向上跃迁不是两组独立速率；热因子 $\ee^{-\beta\hbar\omega}$ 精确固定了它们的比值。

系统自身的 Gibbs 态记为
\begin{equation}
\rho_{S,\beta}=\frac{\ee^{-\beta H_S}}{Z_S}.
\label{eq:phys-oqs-system-gibbs-state}
\end{equation}
由 $[H_S,S_\mu(\omega)]=-\hbar\omega S_\mu(\omega)$，对指数函数作级数展开，逐项交换 $H_S$ 与 $S_\mu(\omega)$，得到
\begin{equation}
\boxed{
S_\mu(\omega)\rho_{S,\beta}
=\ee^{-\beta\hbar\omega}
\rho_{S,\beta}S_\mu(\omega).
}
\label{eq:phys-oqs-gibbs-intertwining}
\end{equation}
在 Hermitian 耦合基底中 $S_\mu(-\omega)=S_\mu^\dagger(\omega)$。于是 complete-secular GKSL 耗散子中 $+\omega$ 与 $-\omega$ 的贡献可成对相消；零频块则由\eqnrefs{eq:phys-oqs-kms-detailed-balance}在 $\omega=0$ 处给出的对称性相消。再结合 $[H_S,H_{\rm LS}]=0$，得到
\begin{equation}
\boxed{
\mathcal L_{\rm sec}[\rho_{S,\beta}]=0.
}
\label{eq:phys-oqs-gibbs-stationary}
\end{equation}
因此在当前弱耦合、Markov、complete-secular 的微观极限中，裸系统 Hamiltonian $H_S$ 的 Gibbs 态由 KMS 关系与 Bohr 频率分解严格推出为一个平衡稳态，而不是额外指定的 phenomenological ansatz。这里得到的是驻定性而非自动得到全局热化：只有当 Liouvillian 在相关态空间上具有唯一零模并且其余谱的实部严格为负（例如满足适当的 irreducibility/primitivity 条件）时，任意初态才收敛到该 Gibbs 态；若存在暗子空间、额外守恒量或多个不变扇区，热 detailed balance 仍可成立而稳态并不唯一。

\begin{derivation}[从 KMS 谱关系到 Gibbs 稳态]
由\eqnrefs{eq:phys-oqs-gamma-spectral-representation}，
\[
\Gamma_{\nu\mu}(-\omega)
=\frac{2\pi\lambda^2}{\hbar^2}
\sum_{nm}p_n
(B_\nu)_{nm}(B_\mu)_{nm}^*
\delta(-\omega-\omega_{mn}).
\]
交换 $n\leftrightarrow m$ 后，Hermitian 性给出
\[
(B_\nu)_{mn}(B_\mu)_{mn}^*
=(B_\mu)_{nm}(B_\nu)_{nm}^*,
\]
而在 delta 函数的支撑上
\[
\frac{p_m}{p_n}
=\ee^{-\beta(E_m-E_n)}
=\ee^{-\beta\hbar\omega}.
\]
因此
\[
\Gamma_{\nu\mu}(-\omega)
=\ee^{-\beta\hbar\omega}\Gamma_{\mu\nu}(\omega),
\]
交换指标即得\eqnrefs{eq:phys-oqs-kms-detailed-balance}。

接着从
\[
[H_S,S_\mu(\omega)]=-\hbar\omega S_\mu(\omega)
\]
出发。对伴随作用的指数展开有
\[
\ee^{-\beta H_S}S_\mu(\omega)\ee^{+\beta H_S}
=\sum_{n=0}^{\infty}
\frac{(-\beta)^n}{n!}\operatorname{ad}_{H_S}^nS_\mu(\omega)
=\ee^{+\beta\hbar\omega}S_\mu(\omega).
\]
右乘 $\ee^{-\beta H_S}$ 并重新排列，得到
\[
S_\mu(\omega)\ee^{-\beta H_S}
=\ee^{-\beta\hbar\omega}
\ee^{-\beta H_S}S_\mu(\omega),
\]
除以 $Z_S$ 即为\eqnrefs{eq:phys-oqs-gibbs-intertwining}。

再取任意 $\omega>0$。因为 $S_\mu^\dagger(\omega)S_\nu(\omega)$ 与 $H_S$ 对易，$+\omega$ 块作用于 $\rho_{S,\beta}$ 时可写成
\[
\mathcal D_{+\omega}[\rho_{S,\beta}]
=\rho_{S,\beta}
\sum_{\mu\nu}\Gamma_{\mu\nu}(\omega)
\left[
\ee^{-\beta\hbar\omega}S_\nu S_\mu^\dagger
-S_\mu^\dagger S_\nu
\right].
\]
另一方面，用 $S_\mu(-\omega)=S_\mu^\dagger(\omega)$ 与\eqnrefs{eq:phys-oqs-kms-detailed-balance}，负频率块为
\[
\begin{aligned}
\mathcal D_{-\omega}[\rho_{S,\beta}]
={}&\rho_{S,\beta}
\sum_{\mu\nu}\Gamma_{\mu\nu}(-\omega)
\left[
\ee^{+\beta\hbar\omega}S_\nu^\dagger S_\mu
-S_\mu S_\nu^\dagger
\right]\\
={}&-\mathcal D_{+\omega}[\rho_{S,\beta}],
\end{aligned}
\]
其中第二步只需交换哑指标 $\mu\leftrightarrow\nu$。$\omega=0$ 时，$S_\mu(0)$ 与 $\rho_{S,\beta}$ 对易，且 KMS 与 Hermitian 性共同使 $\Gamma_{\mu\nu}(0)$ 为实对称矩阵，所以零频块也湮灭 $\rho_{S,\beta}$。Hamiltonian 部分则因
\[
[H_S,\rho_{S,\beta}]=0,
\qquad
[H_{\rm LS},\rho_{S,\beta}]=0
\]
而为零，故得到\eqnrefs{eq:phys-oqs-gibbs-stationary}。

\textbf{物理意义。}Gibbs 态 $\rho_{S,\beta}=e^{-\beta H_S}/Z$ 是开放系统在 KMS 平衡条件下的唯一稳态：Lindblad 跳算符由 KMS 关系 $A(\omega)=e^{-\beta\hbar\omega}A(-\omega)$ 约束，确保上跃迁和下跃迁速率满足细致平衡 $\Gamma_+/\Gamma_-=e^{-\beta\hbar\omega}$。这是量子 H 定理的固定点定理：任何初态在 Lindblad 演化下渐近趋于 Gibbs 态。实例：原子与黑体辐射场耦合，稳态布居 $\rho_{ee}/\rho_{gg}=e^{-\beta\hbar\omega_0}$ 恰为 Boltzmann 分布，与 Einstein $A,B$ 系数的细致平衡一致。
\end{derivation}

\begin{note}[完全正性不等于热 detailed balance]
有限粗粒化构造\eqnrefs{eq:phys-oqs-coarse-grained-kossakowski}保证的是 Kossakowski 矩阵半正定，因此 partial-secular 动力学可以保持完全正；但若一个频率簇同时保留 $\omega\neq\omega'$ 的交叉项，\eqnrefs{eq:phys-oqs-kms-detailed-balance}一般不足以让裸 Gibbs 态成为\emph{精确}稳态。只有在完整 Davies/complete-secular 极限，或所保留簇结构额外满足相应量子 detailed-balance 对称性时，\eqnrefs{eq:phys-oqs-gibbs-stationary}才严格成立。另一方面，有限耦合下总体系的精确平衡约化态是
\[
\rho_{S,\beta}^{\rm exact}
=\frac{\Tr_B\ee^{-\beta H_{\rm tot}}}
{\Tr_{SB}\ee^{-\beta H_{\rm tot}}},
\]
一般可由 Hamiltonian of mean force 描述，并不等于裸 $\ee^{-\beta H_S}/Z_S$。因此“GKSL”“CP”“Markov”“热详细平衡”是四个不同层次的性质，不能相互替代。
\end{note}

\section{一般生成元的 Liouvillian、jump 与对称性}
\label{sec:phys-oqs-unified}

标准 GKSL 方程~\eqnrefs{eq:phys-oqs-gksl-standard} 把微观环境信息压缩为 Hamiltonian $H=H_S+H_{\rm LS}$ 与一组 collapse operators $C_q$。它的 jump、Liouvillian 与对称分级结构与具体能级数无关，可以直接在一般算符层面建立。

这一生成元还可以按 jump、recycle 与无跳跃传播三部分重排。


令
\begin{equation}
K=\sum_q C_q^\dagger C_q,
\qquad
\boxed{H_{\rm eff}=H-\frac{\ii\hbar}{2}K.}
\label{eq:phys-oqs-heff-general}
\end{equation}
对一个无穷小时间步 $\dd t$，定义
\begin{equation}
M_0=\id()-\frac{\ii}{\hbar}H_{\rm eff}\dd t,
\qquad
M_q=\sqrt{\dd t}\,C_q\quad(q\ge1).
\label{eq:phys-oqs-infinitesimal-kraus}
\end{equation}
先检查概率守恒。由
\[
H_{\rm eff}^\dagger-H_{\rm eff}=\ii\hbar K
\]
可得
\begin{align}
M_0^\dagger M_0
&=\id()+\frac{\ii\dd t}{\hbar}
(H_{\rm eff}^\dagger-H_{\rm eff})+O(\dd t^2)\nonumber\\
&=\id()-K\dd t+O(\dd t^2),
\end{align}
而
\[
\sum_{q\ge1}M_q^\dagger M_q
=K\dd t.
\]
因此
\begin{equation}
\boxed{
\sum_{\alpha\ge0}M_\alpha^\dagger M_\alpha
=\id()+O(\dd t^2).}
\label{eq:phys-oqs-infinitesimal-completeness}
\end{equation}
也就是说，GKSL 动力学在每个无穷小时间步都具有精确到 $O(\dd t^2)$ 的 Kraus 表示；精确有限时间映射则由 CPTP 半群 $\ee^{\Liouv t}$ 给出。

把
\[
\rho(t+\dd t)=M_0\rho(t)M_0^\dagger
+\sum_{q\ge1}M_q\rho(t)M_q^\dagger+O(\dd t^2)
\]
展开到一阶：
\begin{align}
M_0\rho M_0^\dagger
&=\rho-\frac{\ii\dd t}{\hbar}
(H_{\rm eff}\rho-\rho H_{\rm eff}^\dagger)+O(\dd t^2),\\
\sum_{q\ge1}M_q\rho M_q^\dagger
&=\dd t\sum_q C_q\rho C_q^\dagger.
\end{align}
除以 $\dd t$ 并令 $\dd t\to0$，就得到与 GKSL 完全等价的分解
\begin{equation}
\boxed{
\dot\rho
=-\frac{\ii}{\hbar}
(H_{\rm eff}\rho-\rho H_{\rm eff}^\dagger)
+\mathcal R[\rho],
\qquad
\mathcal R[\rho]=\sum_q\mathcal J_q[\rho],}
\label{eq:phys-oqs-heff-recycle}
\end{equation}
其中
\begin{equation}
\boxed{\mathcal J_q[\rho]=C_q\rho C_q^\dagger}
\label{eq:phys-oqs-jump-superoperator}
\end{equation}
称为第 $q$ 个 jump 超算符，而 $\mathcal R$ 称为 recycle 超算符。在给定 unraveling 后，前者给“发生第 $q$ 个 jump 后”的未归一化状态，后者把所有 jump 分支的概率全部回注到无条件密度矩阵。

若当前归一化条件态为 $\rho_c$，则\eqnrefs{eq:phys-oqs-infinitesimal-kraus} 立即给出第 $q$ 个 jump 的概率
\begin{equation}
\boxed{
\dd p_q
=\Tr(M_q\rho_cM_q^\dagger)
=\dd t\,\Tr(C_q^\dagger C_q\rho_c),}
\label{eq:phys-oqs-jump-probability}
\end{equation}
以及 jump 后的归一化条件态
\begin{equation}
\boxed{
\rho_c\longrightarrow
\frac{C_q\rho_cC_q^\dagger}
{\Tr(C_q^\dagger C_q\rho_c)}.}
\label{eq:phys-oqs-jump-update}
\end{equation}
这里的 ``jump channel'' 依赖所选的 unraveling，而不是无条件 GKSL 生成元唯一指定的经典事件。若 $U$ 是作用在 jump 指标上的酉矩阵，并定义
\[
C'_a=\sum_qU_{aq}C_q,
\]
则
\[
\begin{aligned}
\sum_aC'_a\rho C_a'^\dagger
&=\sum_{q,q'}(U^\dagger U)_{q'q}C_q\rho C_{q'}^\dagger
=\sum_qC_q\rho C_q^\dagger,\\
\sum_aC_a'^\dagger C'_a
&=\sum_qC_q^\dagger C_q.
\end{aligned}
\]
因此这种酉混合不改变 $\Liouv$、$H_{\rm eff}$ 或无条件态 $\rho(t)$，却会改变单条随机轨迹中“哪一种 jump 发生”的标签。只有在环境测量方案、输出模基底或其他监测协议同时给定以后，\eqnrefs{eq:phys-oqs-jump-probability}--\eqnrefs{eq:phys-oqs-jump-update} 才对应一个具体的实验 unraveling。

无 jump 时，未归一化条件态为 $M_0\rho_cM_0^\dagger$，其迹
\[
1-\dd t\sum_q\Tr(C_q^\dagger C_q\rho_c)+O(\dd t^2)
\]
正好是“没有任何环境事件”的概率。因此 $H_{\rm eff}$ 的非厄密部分不是额外引入的耗散，而是条件在“尚未探测到 jump”这一信息下的概率权重。

对纯态 no-jump 分支，
\begin{equation}
\ii\hbar\frac{\dd}{\dd t}|\psi_{\rm nj}(t)\rangle
=H_{\rm eff}|\psi_{\rm nj}(t)\rangle.
\label{eq:phys-oqs-nojump-schrodinger}
\end{equation}
若 $H_{\rm eff}$ 与时间无关，
\[
|\psi_{\rm nj}(t)\rangle
=\ee^{-\ii H_{\rm eff}t/\hbar}|\psi_0\rangle.
\]
若其某个右本征值写成 $E_r-\ii\hbar\kappa_r/2$，则相应条件振幅包含 $\ee^{-\ii E_rt/\hbar}\ee^{-\kappa_rt/2}$。

接下来考察Liouvillian。

对有限维系统，主方程可以转化成普通线性常微分方程。采用按列堆叠约定
\[
|X\!\rangle\!\rangle\equiv\operatorname{vec}(X),
\qquad
\operatorname{vec}(X)_{i+(j-1)d}=X_{ij}.
\]
需要的核心恒等式是
\begin{equation}
\boxed{
\operatorname{vec}(AXB)
=(B^{\mathsf T}\Kron A)\operatorname{vec}(X).}
\label{eq:phys-oqs-vec-axb}
\end{equation}
这个式子可以直接从指标证明。左边第 $(i,j)$ 个矩阵元为
\[
(AXB)_{ij}=\sum_{k,l}A_{ik}X_{kl}B_{lj}.
\]
右边对应向量分量则是
\begin{align}
[(B^{\mathsf T}\Kron A)\operatorname{vec}(X)]_{i+(j-1)d}
&=\sum_{k,l}(B^{\mathsf T})_{jl}A_{ik}X_{kl}\nonumber\\
&=\sum_{k,l}A_{ik}X_{kl}B_{lj},
\end{align}
与左边逐项相同。

将\eqnrefs{eq:phys-oqs-vec-axb} 分别用于 Hamiltonian、jump 与 loss 项：
\[
\operatorname{vec}(H\rho)=(\id()\Kron H)|\rho\!\rangle\!\rangle,
\qquad
\operatorname{vec}(\rho H)=(H^{\mathsf T}\Kron\id())|\rho\!\rangle\!\rangle,
\]
\[
\operatorname{vec}(C_q\rho C_q^\dagger)
=(C_q^*\Kron C_q)|\rho\!\rangle\!\rangle.
\]
于是任意 $d$ 维 GKSL 方程统一成为
\begin{equation}
\boxed{
\frac{\dd}{\dd t}|\rho\!\rangle\!\rangle
=\Liouv|\rho\!\rangle\!\rangle.}
\label{eq:phys-oqs-liouville-evolution}
\end{equation}
其中
\begin{equation}
\boxed{
\begin{aligned}
\Liouv={}&-\frac{\ii}{\hbar}
(\id()\Kron H-H^{\mathsf T}\Kron\id())\\
&+\sum_q\left[
C_q^*\Kron C_q
-\frac12\id()\Kron C_q^\dagger C_q
-\frac12(C_q^\dagger C_q)^{\mathsf T}\Kron\id()
\right].
\end{aligned}}
\label{eq:phys-oqs-liouvillian-general}
\end{equation}
利用\eqnrefs{eq:phys-oqs-heff-general} 也可写成
\begin{equation}
\boxed{
\Liouv
=-\frac{\ii}{\hbar}
(\id()\Kron H_{\rm eff}-H_{\rm eff}^*\Kron\id())
+\sum_qC_q^*\Kron C_q.}
\label{eq:phys-oqs-liouvillian-heff}
\end{equation}

时间无关时，严格线性解为
\begin{equation}
\boxed{
|\rho(t)\!\rangle\!\rangle
=\ee^{\Liouv t}|\rho(0)\!\rangle\!\rangle.}
\label{eq:phys-oqs-liouvillian-exponential}
\end{equation}
若 $\Liouv(t)$ 随时间变化，则应使用
\begin{equation}
|\rho(t)\!\rangle\!\rangle
=\mathcal T\exp\!\left[\int_0^t\Liouv(s)\,\dd s\right]
|\rho(0)\!\rangle\!\rangle.
\label{eq:phys-oqs-time-dependent-liouvillian}
\end{equation}

迹守恒在 Liouville 空间中表现为一个左零模。定义
\[
\langle\!\langle\id()|X\!\rangle\!\rangle=\Tr X.
\]
因为 $\dd\Tr\rho/\dd t=0$ 对任意 $\rho$ 都成立，所以
\begin{equation}
\boxed{\langle\!\langle\id()|\Liouv=0.}
\label{eq:phys-oqs-left-zero-mode}
\end{equation}
稳态则是右零模
\begin{equation}
\boxed{\Liouv|\rho_{\rm ss}\!\rangle\!\rangle=0,
\qquad \Tr\rho_{\rm ss}=1.}
\label{eq:phys-oqs-steady-state-liouvillian}
\end{equation}
若 $\Liouv$ 可对角化，其他右本征模满足
\[
\Liouv|R_\alpha\!\rangle\!\rangle
=\lambda_\alpha|R_\alpha\!\rangle\!\rangle.
\]
有限维 CPTP 半群的生成元满足 $\operatorname{Re}\lambda_\alpha\le0$；$-\operatorname{Re}\lambda_\alpha$ 给相应无条件模的衰减率，$\operatorname{Im}\lambda_\alpha$ 给其振荡频率。最靠近零的非零实部决定长时间弛豫尺度，这与后面 $H_{\rm eff}$ 的条件极点是不同的谱对象。

接下来考察对称分级与荷扇区。

流形方法的本质不是“二能级技巧”，而是 Liouvillian 对某个荷算符诱导的伴随作用具有协变性。这里必须区分“对称分级”与“荷本身守恒”：jump 可以按确定步长改变荷，却仍保持 operator space 的荷差分块。

取任意 Hermitian 算符 $Q$，其离散谱投影为
\begin{equation}
Q=\sum_n q_nP_n,
\qquad
P_nP_m=\delta_{nm}P_n,
\qquad
\sum_nP_n=\id().
\label{eq:phys-oqs-general-charge-projectors}
\end{equation}
若相干部分满足
\begin{equation}
\boxed{[Q,H]=0,}
\label{eq:phys-oqs-charge-conserving-h}
\end{equation}
并且每个 collapse operator 都是 $Q$ 的伴随本征算符，即存在实数 $s_q$ 使
\begin{equation}
\boxed{[Q,C_q]=-s_qC_q,}
\label{eq:phys-oqs-charge-jump}
\end{equation}
则 $C_q$ 把荷为 $q_n$ 的子空间送到荷为 $q_n-s_q$ 的子空间。由
\[
[Q,C_q^\dagger]=+s_qC_q^\dagger
\]
以及 Leibniz 规则，
\[
[Q,C_q^\dagger C_q]
=[Q,C_q^\dagger]C_q+C_q^\dagger[Q,C_q]=0,
\]
所以 $H_{\rm eff}$ 与 $Q$ 对易。

定义作用在 operator space 上的荷超算符
\begin{equation}
\boxed{
\mathcal C_Q[X]\equiv[Q,X].
}
\label{eq:phys-oqs-charge-superoperator}
\end{equation}
Hamiltonian 部分满足
\[
[Q,[H,X]]=[H,[Q,X]],
\]
而对任一 jump，先计算 recycle 项
\[
\begin{aligned}
[Q,C_qXC_q^\dagger]
={}&[Q,C_q]XC_q^\dagger
+C_q[Q,X]C_q^\dagger
+C_qX[Q,C_q^\dagger]\\
={}&C_q[Q,X]C_q^\dagger,
\end{aligned}
\]
其中首尾两个 $s_q$ 项精确抵消。又因 $[Q,C_q^\dagger C_q]=0$，
\[
\left[Q,\{C_q^\dagger C_q,X\}\right]
=\{C_q^\dagger C_q,[Q,X]\}.
\]
故整个 GKSL 生成元满足
\begin{equation}
\boxed{
[\mathcal C_Q,\Liouv]=0.
}
\label{eq:phys-oqs-liouvillian-charge-covariance}
\end{equation}
这才是后续分块传播真正需要的对称性：若 $X$ 满足 $[Q,X]=\Delta q\,X$，则 $\Liouv X$ 仍处在同一个荷差 $\Delta q$ 扇区。

与此同时，$Q$ 本身通常\emph{不守恒}。伴随生成元作用于 $Q$ 时，Hamiltonian 部分为零，而
\[
\begin{aligned}
\mathcal D^\dagger[C_q]Q
&=C_q^\dagger Q C_q
-\frac12\{C_q^\dagger C_q,Q\}\\
&=C_q^\dagger C_q(Q-s_q)-C_q^\dagger C_qQ\\
&=-s_qC_q^\dagger C_q.
\end{aligned}
\]
因此
\begin{equation}
\boxed{
\frac{\dd}{\dd t}\langle Q\rangle
=-\sum_qs_q\,\langle C_q^\dagger C_q\rangle.
}
\label{eq:phys-oqs-charge-expectation-flow}
\end{equation}
只有当所有有效 jump 都满足 $s_q=0$，或更一般地右边因额外结构恒等为零时，$Q$ 才是完整开放动力学的守恒量。特别地，$s_q=0$ 对所有 $q$ 时，$[Q,H]=[Q,C_q]=0$ 给出通常所称的 strong symmetry；此时每个荷本征子空间都由完整动力学保持。

为了写出具有固定步长的分级结构，以下假设相关荷本征值可以等距编号为整数 $N$，并把每个 jump 的荷变化 $s_q$ 用同一基本间隔计为整数，相应投影记为 $P_N$；若谱并非等距，只需保留实际的 $q_n$ 标签与真实荷差，下面的投影论证不变。定义
\begin{equation}
\boxed{R_{NM}=P_N\rho P_M,}
\label{eq:phys-oqs-density-block-general}
\end{equation}
以及
\begin{equation}
\boxed{
M_N=P_NH_{\rm eff}P_N,
\qquad
J_q^{(N)}=P_{N-s_q}C_qP_N.}
\label{eq:phys-oqs-general-block-operators}
\end{equation}
相干部分逐项投影为
\[
P_NH_{\rm eff}\rho P_M=M_NR_{NM},
\qquad
P_N\rho H_{\rm eff}^\dagger P_M=R_{NM}M_M^\dagger.
\]
对 recycle 项，在 $\rho=\sum_{KL}P_K\rho P_L$ 中插入完备分解：
\begin{align}
P_NC_q\rho C_q^\dagger P_M
&=\sum_{K,L}P_NC_qP_K\,R_{KL}\,P_LC_q^\dagger P_M.
\end{align}
\eqnrefs{eq:phys-oqs-charge-jump} 要求只有
\[
K=N+s_q,
\qquad
L=M+s_q
\]
能留下非零贡献，因此
\[
P_NC_q\rho C_q^\dagger P_M
=J_q^{(N+s_q)}R_{N+s_q,M+s_q}J_q^{(M+s_q)\dagger}.
\]
最终得到一般分级主方程
\begin{equation}
\boxed{
\begin{aligned}
\dot R_{NM}
={}&-\frac{\ii}{\hbar}
(M_NR_{NM}-R_{NM}M_M^\dagger)\\
&+\sum_q
J_q^{(N+s_q)}R_{N+s_q,M+s_q}
J_q^{(M+s_q)\dagger}.
\end{aligned}}
\label{eq:phys-oqs-general-block-master}
\end{equation}
因为 recycle 同时把左右荷标签各降低同一个 $s_q$，荷差 $N-M$ 在\eqnrefs{eq:phys-oqs-general-block-master}中保持不变，这正是\eqnrefs{eq:phys-oqs-liouvillian-charge-covariance}的投影表示。若 $s_q>0$，population 可以沿 $N$ 方向逐级流向低荷扇区，但不同荷差的 coherence sector 仍不会混合。

对一般分级只需定义
\[
\mathcal H_N=\operatorname{Ran}P_N,
\qquad
d_N=\operatorname{rank}P_N.
\]
低荷扇区的维数可以与高荷扇区不同；任何矩阵传播都必须从实际的 $P_N$ 与 $d_N$ 构造，而不能从某个具体腔模型的高激发矩阵反推边界扇区。单模 RWA 中如何由原子跃迁图求出整数激发等级 $\nu_i$，将在最后的腔 QED 应用中再具体化。

接下来考察矩阵传播。

选定某个 $N$ 流形的有序基
\[
\mathcal B_N=
\{|\chi_1^{(N)}\rangle,\ldots,|\chi_{d_N}^{(N)}\rangle\}.
\]
对 no-jump 纯态写
\[
|\psi_N(t)\rangle
=\sum_{r=1}^{d_N}c_r(t)|\chi_r^{(N)}\rangle,
\qquad
\mathbf c_N=(c_1,\ldots,c_{d_N})^{\mathsf T}.
\]
将\eqnrefs{eq:phys-oqs-nojump-schrodinger} 左乘 $P_N$，并用 $P_NH_{\rm eff}=P_NH_{\rm eff}P_N$，得到
\[
\ii\hbar P_N\frac{\dd}{\dd t}|\psi_N\rangle
=M_N|\psi_N\rangle.
\]
再左乘每个 $\langle\chi_r^{(N)}|$，就得到有限维矩阵方程
\begin{equation}
\boxed{
\ii\hbar\dot{\mathbf c}_N=M_N\mathbf c_N,
\qquad
\mathbf c_N(t)=\ee^{-\ii M_Nt/\hbar}\mathbf c_N(0).}
\label{eq:phys-oqs-manifold-propagator}
\end{equation}

下面只证明一次任意维矩阵指数的谱形式。定义
\[
p_N(z)=\det(z\id(d_N)-M_N).
\]
若 $M_N$ 有 $d_N$ 个互异本征值 $z_{r,N}$，定义 Lagrange 谱投影多项式
\begin{equation}
\Pi_{r,N}
=\prod_{s\ne r}
\frac{M_N-z_{s,N}\id()}{z_{r,N}-z_{s,N}}.
\label{eq:phys-oqs-lagrange-projector}
\end{equation}
对任一本征向量 $|v_{k,N}\rangle$，有
\[
\Pi_{r,N}|v_{k,N}\rangle
=\prod_{s\ne r}
\frac{z_{k,N}-z_{s,N}}{z_{r,N}-z_{s,N}}
|v_{k,N}\rangle
=\delta_{rk}|v_{k,N}\rangle.
\]
因此
\[
\Pi_{r,N}\Pi_{s,N}=\delta_{rs}\Pi_{r,N},
\qquad
\sum_r\Pi_{r,N}=\id(),
\qquad
M_N=\sum_r z_{r,N}\Pi_{r,N}.
\]
对任意解析函数 $f$，幂级数逐项作用后便有
\[
f(M_N)=\sum_r f(z_{r,N})\Pi_{r,N}.
\]
取 $f(z)=\ee^{-\ii zt/\hbar}$，得到
\begin{equation}
\boxed{
\ee^{-\ii M_Nt/\hbar}
=\sum_{r=1}^{d_N}
\ee^{-\ii z_{r,N}t/\hbar}
\prod_{s\ne r}
\frac{M_N-z_{s,N}\id(d_N)}{z_{r,N}-z_{s,N}}.}
\label{eq:phys-oqs-spectral-propagator}
\end{equation}
若本征值简并但矩阵仍可对角化，只需把同一简并本征值对应的投影合并。若矩阵在 exceptional point 处不可对角化，设某个 Jordan 块为
\[
J_r=z_r\id()+N_r,
\qquad N_r^{m_r}=0,
\]
则
\begin{equation}
\ee^{-\ii J_rt/\hbar}
=\ee^{-\ii z_rt/\hbar}
\sum_{k=0}^{m_r-1}\frac1{k!}
\left(-\frac{\ii t}{\hbar}N_r\right)^k,
\label{eq:phys-oqs-jordan-propagator}
\end{equation}
所以矩阵指数在简并点仍然完全有定义，只是会多出多项式乘指数的时间因子。

无条件块方程~\eqnrefs{eq:phys-oqs-general-block-master} 具有相同的齐次传播子。定义
\[
U_N(t)=\ee^{-\ii M_Nt/\hbar},
\]
以及 recycle 源
\begin{equation}
S_{NM}(t)=\sum_q
J_q^{(N+s_q)}R_{N+s_q,M+s_q}(t)
J_q^{(M+s_q)\dagger}.
\label{eq:phys-oqs-block-source}
\end{equation}
由 $\dot U_N=-(\ii/\hbar)M_NU_N$ 可得
\[
\frac{\dd}{\dd t}U_N^{-1}(t)
=+\frac{\ii}{\hbar}U_N^{-1}(t)M_N.
\]
因此对乘积逐项求导：
\begin{align}
\frac{\dd}{\dd t}
[U_N^{-1}R_{NM}U_M^{-\dagger}]
={}&\frac{\ii}{\hbar}U_N^{-1}M_NR_{NM}U_M^{-\dagger}
+U_N^{-1}\dot R_{NM}U_M^{-\dagger}\nonumber\\
&-\frac{\ii}{\hbar}U_N^{-1}R_{NM}M_M^\dagger U_M^{-\dagger}.
\end{align}
代入\eqnrefs{eq:phys-oqs-general-block-master} 后，两个齐次项精确抵消，只剩
\[
\frac{\dd}{\dd t}
[U_N^{-1}R_{NM}U_M^{-\dagger}]
=U_N^{-1}S_{NM}U_M^{-\dagger}.
\]
从 $0$ 到 $t$ 积分并恢复左右传播子，得到
\begin{equation}
\boxed{
\begin{aligned}
R_{NM}(t)
={}&U_N(t)R_{NM}(0)U_M^\dagger(t)\\
&+\int_0^t\dd\tau\,
U_N(t-\tau)S_{NM}(\tau)U_M^\dagger(t-\tau).
\end{aligned}}
\label{eq:phys-oqs-duhamel-block}
\end{equation}
第一项是当前块的 no-jump 齐次传播，第二项是所有其他流形经 jump 后在过去各时刻注入该块的历史总和。

若 $M_N(t)$ 随时间变化，结论仍成立，只需把 $U_N(t-\tau)$ 换成满足
\[
\ii\hbar\partial_tU_N(t,\tau)=M_N(t)U_N(t,\tau),
\qquad U_N(\tau,\tau)=\id()
\]
的两时传播子。也就是说，谱分解是求时间无关矩阵指数的一种计算工具，而“齐次传播 + Duhamel 源项”这一结构本身并不依赖时间无关假设。

由精确约化映射到 Liouvillian 的一般链条至此闭合。Born 截断控制耦合阶数，TCL 重组记忆为时间局域生成元，Markov 控制环境记忆时间，有限粗粒化与 secular 决定频率分辨率；这些条件彼此独立。具体环境谱、腔模与有限能级结构只改变一般公式中的谱矩阵、Hamiltonian 和 collapse operators。

\section{应用：环境谱、微观跃迁率与腔 QED}
\label{sec:phys-oqs-bath-examples}

一般生成元、近似层级和求解表示已经建立以后，才需要把具体环境相关函数代入。下面三个例子分别展示有限带宽相关、热玻色环境与自由空间电磁真空如何生成速率矩阵和 Lamb shift；它们不改变前面的动力学结构。本章一直保留的 $\lambda$ 只是系统--环境耦合阶数的 bookkeeping 参数；以下物理耦合强度已经写入 $g_{\alpha q}$、偶极矩阵元等实际参数中，因此完成微扰阶数审计后可令 $\lambda=1$。

\begin{exbox}[Lorentzian 环境]
取与前文 Markov 分析一致的相关时间 $\tau_B>0$，并令
\[
C(\tau)
=A\,\ee^{-|\tau|/\tau_B}
\ee^{-\ii\omega_B\tau},
\qquad A>0.
\]
其中 $\omega_B$ 是谱峰中心，$\tau_B$ 正是环境记忆衰减的时间尺度。双边谱为
\[
\begin{aligned}
\widetilde C(\omega)
&=A\int_{-\infty}^{+\infty}\dd\tau\,
\ee^{-|\tau|/\tau_B}
\ee^{\ii(\omega-\omega_B)\tau}\\
&=\frac{2A\tau_B}
{1+(\omega-\omega_B)^2\tau_B^2},
\end{aligned}
\]
而单边响应为
\[
G(\omega)
=\frac{A}{\tau_B^{-1}-\ii(\omega-\omega_B)}
=A\tau_B\frac{1+\ii(\omega-\omega_B)\tau_B}
{1+(\omega-\omega_B)^2\tau_B^2}.
\]
代入\eqnrefs{eq:phys-oqs-gamma-from-g} 与\eqnrefs{eq:phys-oqs-delta-ls}，
\[
\Gamma(\omega)
=\frac{2\lambda^2A\tau_B/\hbar^2}
{1+(\omega-\omega_B)^2\tau_B^2},
\qquad
\Delta^{\rm LS}(\omega)
=\frac{\lambda^2A\tau_B^2}{\hbar}
\frac{\omega-\omega_B}
{1+(\omega-\omega_B)^2\tau_B^2}.
\]
若某个频率块只有一条下降跃迁 $S(\omega_0)=|j\rangle\langle i|$，则
\[
S^\dagger(\omega_0)S(\omega_0)=|i\rangle\langle i|,
\]
从而
\begin{equation}
\boxed{
H_{\rm LS}^{(\omega_0)}
=\Delta^{\rm LS}(\omega_0)|i\rangle\langle i|.
}
\label{eq:phys-oqs-lorentzian-lamb-shift}
\end{equation}
若同一频率块包含多条不可分辨跃迁，则应保留完整的 $\Delta_{\mu\nu}^{\rm LS}(\omega_0)$ 矩阵，并按\eqnrefs{eq:phys-oqs-hls-general} 形成跃迁之间的环境诱导相干耦合。
\end{exbox}

\begin{exbox}[热玻色环境]
考虑一组系统下降算符 $A_\alpha$ 与玻色环境交换能量，
\begin{equation}
\boxed{
H_B=\sum_q\hbar\nu_qb_q^\dagger b_q,
\qquad
H_I^{\rm ex}
=\hbar\sum_{\alpha q}
\left(
 g_{\alpha q}A_\alpha b_q^\dagger
+g_{\alpha q}^*A_\alpha^\dagger b_q
\right).
}
\label{eq:phys-oqs-boson-exchange}
\end{equation}
沿用\eqnrefs{eq:phys-oqs-bath-gibbs-state} 的热平衡记号，环境状态为
\[
\rho_B=\frac{\ee^{-\beta H_B}}{Z_B},
\qquad
\beta=\frac1{k_BT}.
\]
单模平均占据数为
\begin{equation}
\boxed{
\bar n_T(\nu)=\frac1{\exp[\hbar\nu/(k_BT)]-1}.
}
\label{eq:phys-oqs-bose-occupation}
\end{equation}
从热态的 Fock 权重 $p_{n_q}$ 直接有
\[
\sum_{n_q}p_{n_q}n_q=\bar n_T(\nu_q),
\qquad
\sum_{n_q}p_{n_q}(n_q+1)=\bar n_T(\nu_q)+1.
\]
因此环境相关函数满足
\[
\langle b_q(\tau)b_{q'}^\dagger(0)\rangle
=\delta_{qq'}[\bar n_T(\nu_q)+1]\ee^{-\ii\nu_q\tau},
\]
\[
\langle b_q^\dagger(\tau)b_{q'}(0)\rangle
=\delta_{qq'}\bar n_T(\nu_q)\ee^{+\ii\nu_q\tau}.
\]
定义正频谱密度矩阵
\begin{equation}
\boxed{
J_{\alpha\beta}(\nu)
=\sum_qg_{\alpha q}g_{\beta q}^*\delta(\nu-\nu_q),
\qquad \nu>0.
}
\label{eq:phys-oqs-boson-j}
\end{equation}
对一个正 Bohr 频率 $\omega>0$，若 $\alpha,\beta$ 标记该频率块内的所有下降跃迁，则在壳的 $\delta$ 项给出
\begin{equation}
\boxed{
\Gamma_{\alpha\beta}^{\downarrow}(\omega,T)
=2\pi\lambda^2J_{\alpha\beta}(\omega)
[\bar n_T(\omega)+1],
}
\label{eq:phys-oqs-thermal-down-matrix}
\end{equation}
而反向吸收块为
\begin{equation}
\boxed{
\Gamma_{\alpha\beta}^{\uparrow}(\omega,T)
=2\pi\lambda^2J_{\beta\alpha}(\omega)
\bar n_T(\omega).
}
\label{eq:phys-oqs-thermal-up-matrix}
\end{equation}
这两式是任意多能级系统的矩阵形式。若该频率块只有一条独立跃迁 $i\to j$，$E_i>E_j$，定义
\[
\Gamma_{i\to j}^{(0)}
=2\pi\lambda^2J_{i\to j}(\omega_{ij}),
\]
则矩阵公式退化为
\begin{equation}
\boxed{
\Gamma_{i\to j}(T)
=\Gamma_{i\to j}^{(0)}[\bar n_T(\omega_{ij})+1],
\qquad
\Gamma_{j\to i}(T)
=\Gamma_{i\to j}^{(0)}\bar n_T(\omega_{ij}).
}
\label{eq:phys-oqs-thermal-transition-rates}
\end{equation}
并满足
\[
\frac{\Gamma_{j\to i}(T)}{\Gamma_{i\to j}(T)}
=\ee^{-\hbar\omega_{ij}/(k_BT)}.
\]
所以二能级热跃迁率只是\eqnrefs{eq:phys-oqs-thermal-down-matrix}--\eqnrefs{eq:phys-oqs-thermal-up-matrix} 在只有一条跃迁通道时的标量特例。
\end{exbox}

\begin{exbox}[自由空间多能级电磁场]
设物质 Hamiltonian 为
\[
H_A=\sum_iE_i|i\rangle\langle i|,
\]
并通过 $H_I=-\mathbf d\cdot\mathbf E(\mathbf r_0)$ 与三维自由空间电磁场耦合。对任意允许下降跃迁
\[
\alpha=(i\to j),
\qquad
E_i>E_j,
\qquad
\omega_\alpha=\omega_{ij},
\]
定义
\begin{equation}
\boxed{
A_\alpha=|j\rangle\langle i|,
\qquad
\mathbf d_\alpha
=\langle j|\mathbf d|i\rangle.
}
\label{eq:phys-oqs-transition-alpha}
\end{equation}
在交换型 RWA 中，一个模式 $(\mathbf k,\lambda)$ 的耦合常数可以写成
\[
g_{\alpha,\mathbf k\lambda}
=\ii
\sqrt{\frac{\omega_k}{2\hbar\epsilon_0V_{\rm q}}}
\,\mathbf d_\alpha\cdot\mathbf e_{\mathbf k\lambda}^*
\ee^{-\ii\mathbf k\cdot\mathbf r_0}.
\]
固定一个正 Bohr 频率 $\omega$，把所有满足 $\omega_\alpha=\omega$ 的跃迁组成集合 $\mathcal T_\omega$。由\eqnrefs{eq:phys-oqs-thermal-down-matrix} 的零温极限，
\[
\Gamma_{\alpha\beta}^{(0)}(\omega)
=2\pi\lambda^2
\sum_{\mathbf k,\lambda}
 g_{\alpha,\mathbf k\lambda}
 g_{\beta,\mathbf k\lambda}^*
\delta(\omega-\omega_k),
\qquad
\alpha,\beta\in\mathcal T_\omega.
\]
采用连续极限
\[
\sum_{\mathbf k}
\longrightarrow
\frac{V_{\rm q}}{(2\pi)^3}
\int_0^\infty k^2\dd k\int\dd\Omega,
\]
横向偏振完备关系
\[
\sum_\lambda e_{\mathbf k\lambda,a}e_{\mathbf k\lambda,b}^*
=\delta_{ab}-\hat k_a\hat k_b,
\]
以及角积分
\[
\int\dd\Omega\,
(\delta_{ab}-\hat k_a\hat k_b)
=\frac{8\pi}{3}\delta_{ab},
\]
可得
\[
\int\dd\Omega\sum_\lambda
(\mathbf d_\alpha\cdot\mathbf e_{\mathbf k\lambda}^*)
(\mathbf d_\beta^*\cdot\mathbf e_{\mathbf k\lambda})
=\frac{8\pi}{3}\mathbf d_\alpha\cdot\mathbf d_\beta^*.
\]
径向积分为
\[
\int_0^\infty k^2\dd k\,(ck)\delta(\omega-ck)
=\frac{\omega^3}{c^3}.
\]
代回模式和，得到任意同频跃迁的自由空间零温速率矩阵
\begin{equation}
\boxed{
\Gamma_{\alpha\beta}^{(0)}(\omega)
=\lambda^2
\frac{\omega^3}{3\pi\epsilon_0\hbar c^3}
\,\mathbf d_\alpha\cdot\mathbf d_\beta^*,
\qquad
\alpha,\beta\in\mathcal T_\omega.
}
\label{eq:phys-oqs-free-space-gamma-matrix-zerot}
\end{equation}
它是跃迁偶极向量的 Gram 矩阵乘以非负因子。对任意系数 $c_\alpha$，
\[
\sum_{\alpha\beta}c_\alpha^*
\Gamma_{\alpha\beta}^{(0)}c_\beta
=\lambda^2\frac{\omega^3}{3\pi\epsilon_0\hbar c^3}
\left|\sum_\alpha c_\alpha^*\mathbf d_\alpha\right|^2\ge0,
\]
这与一般的半正定性\eqnrefs{eq:phys-oqs-gamma-psd} 一致。若同一频率块中两条跃迁的偶极矩阵元不正交，\eqnrefs{eq:phys-oqs-free-space-gamma-matrix-zerot} 的非对角元就是共同真空产生的交叉阻尼；若偶极正交，该交叉项消失。

在温度 $T$ 的自由空间热光子场中，直接乘上\eqnrefs{eq:phys-oqs-thermal-down-matrix}--\eqnrefs{eq:phys-oqs-thermal-up-matrix} 的 Bose 因子：
\begin{equation}
\boxed{
\Gamma_{\alpha\beta}^{\downarrow}(\omega,T)
=\Gamma_{\alpha\beta}^{(0)}(\omega)
[\bar n_T(\omega)+1],
\qquad
\Gamma_{\alpha\beta}^{\uparrow}(\omega,T)
=\Gamma_{\beta\alpha}^{(0)}(\omega)
\bar n_T(\omega).
}
\label{eq:phys-oqs-free-space-thermal-matrix}
\end{equation}
若某个 Bohr 频率只有一条跃迁 $i\to j$，则
\begin{equation}
\boxed{
\Gamma_{i\to j}^{(0)}
=\lambda^2
\frac{\omega_{ij}^3|\mathbf d_{ji}|^2}
{3\pi\epsilon_0\hbar c^3},
}
\label{eq:phys-oqs-free-space-single-rate}
\end{equation}
并与\eqnrefs{eq:phys-oqs-thermal-transition-rates} 直接组合得到有限温度的上下跃迁率。

同一自由空间谱的主值部分产生 Lamb shift。对交换型耦合的正频块，定义
\[
J_{\alpha\beta}^{\rm fs}(\nu)
=\frac{\nu^3}{6\pi^2\epsilon_0\hbar c^3}
\mathbf d_\alpha\cdot\mathbf d_\beta^*,
\]
使 $\Gamma_{\alpha\beta}^{(0)}(\omega)=2\pi\lambda^2J_{\alpha\beta}^{\rm fs}(\omega)$。则带高频截断 $\varpi_{\rm UV}$ 的正频主值贡献为
\begin{equation}
\boxed{
\Delta_{\alpha\beta}^{\rm LS,+}(\omega;\varpi_{\rm UV})
=\hbar\lambda^2\,\mathcal P
\int_0^{\varpi_{\rm UV}}\dd\varpi\,
\frac{J_{\alpha\beta}^{\rm fs}(\varpi)}{\omega-\varpi}.
}
\label{eq:phys-oqs-free-space-ls-general}
\end{equation}
由上式 $J_{\alpha\beta}^{\rm fs}(\varpi)=A_{\alpha\beta}\varpi^3$，其中 $A_{\alpha\beta}=\mathbf d_\alpha\cdot\mathbf d_\beta^*/(6\pi^2\epsilon_0\hbar c^3)$，故主值积分的频率部分为
\[
\frac{\varpi^3}{\omega-\varpi}
=-(\varpi^2+\omega\varpi+\omega^2)
+\frac{\omega^3}{\omega-\varpi},
\]
故
\[
\begin{aligned}
\mathcal P\int_0^{\varpi_{\rm UV}}\dd\varpi\,
\frac{\varpi^3}{\omega-\varpi}
={}&-\frac{\varpi_{\rm UV}^3}{3}
-\frac{\omega\varpi_{\rm UV}^2}{2}
-\omega^2\varpi_{\rm UV}\\
&-\omega^3
\ln\!\left|\frac{\varpi_{\rm UV}-\omega}{\omega}\right|.
\end{aligned}
\]
因此点偶极连续谱的绝对能级位移具有紫外敏感性。若从完整偶极 Hamiltonian 而不是交换型 RWA 出发，负频/反旋虚过程还会给出额外主值项，但高频重整化问题仍然存在。可观测量应取完整原子结构与重整化后得到的能级差或跃迁频移，而\eqnrefs{eq:phys-oqs-hls-general} 仍是把这些主值系数组装成多能级 $H_{\rm LS}$ 的统一公式。
\end{exbox}


\label{sec:phys-oqs-manifoldmodels}
接下来考察多能级物质与单模腔的 Hamiltonian。

这一节开始才把一般生成元代入腔 QED。物质 Hamiltonian 与跃迁算符定义为
\begin{equation}
\boxed{
H_A=\sum_iE_i\sigma_{ii},
\qquad
\sigma_{ij}=|i\rangle\langle j|.
}
\label{eq:phys-oqs-atomic-h-general}
\end{equation}
它们满足
\[
\sigma_{ij}\sigma_{kl}=\delta_{jk}\sigma_{il},
\qquad
[\sigma_{ij},\sigma_{kl}]
=\delta_{jk}\sigma_{il}-\delta_{li}\sigma_{kj}.
\]

单一量子模式的自由 Hamiltonian 为
\begin{equation}
\boxed{
H_c=\hbar\omega_c\left(a^\dagger a+\frac12\right),
\qquad [a,a^\dagger]=1.}
\label{eq:phys-oqs-cavity-h}
\end{equation}
物质偶极算符写成
\[
\mathbf d=\sum_{ij}\mathbf d_{ij}\sigma_{ij},
\qquad
\mathbf d_{ij}=\langle i|\mathbf d|j\rangle,
\qquad
\mathbf d_{ji}=\mathbf d_{ij}^*.
\]
为了使物质与场使用同一套频率符号，先按自由物质 Hamiltonian $H_A$ 的 Heisenberg 时间依赖定义偶极的正、负频率部分。对 $E_i>E_j$，算符 $\sigma_{ji}$ 随 $\ee^{-\ii\omega_{ij}t}$ 演化，因此正频率部分必须是降能算符，而负频率部分是它的厄密共轭：
\begin{equation}
\boxed{
\mathbf d^{(+)}=\sum_{E_i>E_j}\mathbf d_{ji}\sigma_{ji},
\qquad
\mathbf d^{(-)}=(\mathbf d^{(+)})^\dagger
=\sum_{E_i>E_j}\mathbf d_{ij}\sigma_{ij}.}
\label{eq:phys-oqs-dipole-positive-negative}
\end{equation}
这里上标 $(\pm)$ 表示 Fourier 正、负频率，而不是“升能/降能”的符号；按上述约定，$\mathbf d^{(+)}$ 恰好降低物质能量。若体系具有永久偶极矩，
\[
\mathbf d^{(0)}=\sum_i\mathbf d_{ii}\sigma_{ii}
\]
是额外的零频部分。它与腔场耦合时产生 $-\mathbf d^{(0)}\cdot(\mathbf E^{(+)}+\mathbf E^{(-)})$ 的纵向/位移型耦合，并不会自动化成下面的激发交换项。对具有反演或适当宇称对称性的原子常有 $\mathbf d_{ii}=0$；若永久偶极不为零，则必须显式保留该纵向项，或在给出独立小参数后再消去，不能把它无条件并入 exchange RWA。若 $\mathbf d^{(0)}$ 耦合的是外部零频噪声，它还会产生纯退相干或静态位移通道。

在物质位置 $\mathbf r_0$，令 $\mathbf e_c$ 为腔模的单位偏振向量，$f_c(\mathbf r)$ 为无量纲复模函数，$V_{\rm eff}$ 为与该模函数归一化配套的有效模体积；模函数的整体归一化约定只允许在 $f_c$ 与 $V_{\rm eff}$ 之间重新分配，不能重复计入。单模电场正频率部分写为
\begin{equation}
\mathbf E^{(+)}(\mathbf r_0)
=\ii\sqrt{\frac{\hbar\omega_c}{2\epsilon_0V_{\rm eff}}}
\,\mathbf e_c f_c(\mathbf r_0)a,
\qquad
\mathbf E^{(-)}=(\mathbf E^{(+)})^\dagger.
\label{eq:phys-oqs-cavity-field}
\end{equation}
其中 $\mathbf E^{(+)}$ 随自由演化携带 $\ee^{-\ii\omega_ct}$，所以与\eqnrefs{eq:phys-oqs-dipole-positive-negative} 的正频率约定一致。相应地，\eqnrefs{eq:phys-oqs-gij-def} 中的 $g_{ij}$ 具有角频率量纲。
电偶极相互作用为
\begin{equation}
\boxed{H_{Ac}=-\mathbf d\cdot(\mathbf E^{(+)}+\mathbf E^{(-)}).}
\label{eq:phys-oqs-dipole-interaction}
\end{equation}
尚未施加系统内部 RWA 时，
\begin{equation}
\boxed{H_S^{\rm full}=H_A+H_c+H_{Ac}.}
\label{eq:phys-oqs-hs-full}
\end{equation}

以
\[
H_{0,S}=H_A+H_c
\]
定义系统内部相互作用绘景。对每条 $E_i>E_j$ 的跃迁，$\mathbf d^{(-)}\mathbf E^{(+)}$ 与 $\mathbf d^{(+)}\mathbf E^{(-)}$ 携带慢相位 $\ee^{\pm\ii(\omega_{ij}-\omega_c)t}$，描述吸收一个腔光子与发射一个腔光子；$\mathbf d^{(-)}\mathbf E^{(-)}$ 与 $\mathbf d^{(+)}\mathbf E^{(+)}$ 则携带 $\ee^{\pm\ii(\omega_{ij}+\omega_c)t}$ 的反旋相位。对一般复偏振或行波模，共旋与反旋项读取的是不同的复偶极投影，因此分别定义
\begin{equation}
\boxed{
\begin{aligned}
 g_{ij}
 &=-\ii\sqrt{\frac{\omega_c}{2\hbar\epsilon_0V_{\rm eff}}}
 \,\mathbf d_{ij}\!\cdot\!\mathbf e_c f_c(\mathbf r_0),\\
 g_{ij}^{\rm cr}
 &=+\ii\sqrt{\frac{\omega_c}{2\hbar\epsilon_0V_{\rm eff}}}
 \,\mathbf d_{ij}\!\cdot\!\mathbf e_c^* f_c^*(\mathbf r_0),
 \qquad E_i>E_j.
\end{aligned}}
\label{eq:phys-oqs-gij-def}
\end{equation}
忽略永久偶极项后，跃迁部分的完整偶极耦合可写成
\begin{equation}
\boxed{
H_{Ac}^{\rm tr}
=\hbar\sum_{E_i>E_j}
\left(
 g_{ij}a\sigma_{ij}+g_{ij}^*a^\dagger\sigma_{ji}
 +g_{ij}^{\rm cr}a^\dagger\sigma_{ij}
 +g_{ij}^{{\rm cr}*}a\sigma_{ji}
\right).}
\label{eq:phys-oqs-cavity-rot-counter-coupling}
\end{equation}
若偏振与模函数可同时取实（差一个整体相位也可吸收入 $a$），则 $|g_{ij}^{\rm cr}|=|g_{ij}|$；一般复模中不能预先假定这一等式。在系统内部旋波近似可控时，保留交换激发的慢相位对，得到
\begin{equation}
\boxed{
H_S^{\rm RWA}
=\sum_iE_i\sigma_{ii}
+\hbar\omega_c\left(a^\dagger a+\frac12\right)
+\hbar\sum_{E_i>E_j}
\left(g_{ij}a\sigma_{ij}+g_{ij}^*a^\dagger\sigma_{ji}\right).}
\label{eq:phys-oqs-hs-rwa}
\end{equation}
这里仍是一般多能级单模 Hamiltonian；具体构形只通过哪些 $g_{ij}$ 非零来指定。RWA 的控制量不是“小失谐”本身。若所研究态的相关光子数不超过 $n$，反旋项的最大一阶混合振幅由
\begin{equation}
\boxed{
\epsilon_{\rm RWA}^{(ij,n)}
=\frac{|g_{ij}^{\rm cr}|\sqrt{n+1}}{\omega_{ij}+\omega_c}
\ll1}
\label{eq:phys-oqs-rwa-small-parameter}
\end{equation}
控制。此时反旋项产生的态矢修正为 $O(\epsilon_{\rm RWA})$，而 Bloch--Siegert 型频率修正的量级为
\[
\delta\omega_{\rm BS}
=O\!\left(\frac{|g_{ij}^{\rm cr}|^2(n+1)}{\omega_{ij}+\omega_c}\right).
\]
因此 RWA 与近共振是两个不同条件：即使 $|\omega_{ij}-\omega_c|$ 很大，只要\eqnrefs{eq:phys-oqs-rwa-small-parameter} 足够小，反旋项仍可受控消去，而交换项进入色散区；反之在 ultrastrong-coupling 区，即使共振也不能直接删除反旋项。

\begin{derivation}[从偶极频率分解到交换型 RWA]
由 $H_A\sigma_{ij}-\sigma_{ij}H_A=(E_i-E_j)\sigma_{ij}$，这是算符代数的基本对易关系，\(\sigma_{ij}=|i\rangle\langle j|\) 为原子跃迁算符。对 $E_i>E_j$ 有
\[
\ee^{+\ii H_At/\hbar}\sigma_{ij}\ee^{-\ii H_At/\hbar}
=\ee^{+\ii\omega_{ij}t}\sigma_{ij},
\qquad
\ee^{+\ii H_At/\hbar}\sigma_{ji}\ee^{-\ii H_At/\hbar}
=\ee^{-\ii\omega_{ij}t}\sigma_{ji}.
\]
该关系由 Baker--Campbell--Hausdorff 公式给出：\(e^A B e^{-A}=B+[A,B]+\frac12[A,[A,B]]+\cdots\)，因 \([H_A,\sigma_{ij}]=(E_i-E_j)\sigma_{ij}\) 使高阶换位子闭合，级数求和即得 \(e^{\ii\omega_{ij}t}\sigma_{ij}\)。另一方面，
\[
a(t)=\ee^{-\ii\omega_ct}a,
\qquad
a^\dagger(t)=\ee^{+\ii\omega_ct}a^\dagger.
\]
这些是腔场算符在自由 Hamiltonian \(H_c=\hbar\omega_c a^\dagger a\) 下的 Heisenberg 演化，由 \([H_c,a]=-\hbar\omega_c a\) 与 \([H_c,a^\dagger]=+\hbar\omega_c a^\dagger\) 给出。于是四类非零频乘积分别为
\begin{align*}
\sigma_{ij}(t)a(t)
&=\ee^{+\ii(\omega_{ij}-\omega_c)t}\sigma_{ij}a,
&
\sigma_{ji}(t)a^\dagger(t)
&=\ee^{-\ii(\omega_{ij}-\omega_c)t}\sigma_{ji}a^\dagger,\\
\sigma_{ij}(t)a^\dagger(t)
&=\ee^{+\ii(\omega_{ij}+\omega_c)t}\sigma_{ij}a^\dagger,
&
\sigma_{ji}(t)a(t)
&=\ee^{-\ii(\omega_{ij}+\omega_c)t}\sigma_{ji}a.
\end{align*}
前两式的相位 \(\omega_{ij}-\omega_c\) 在共振时为零，对应能量守恒跃迁；后两式的相位 \(\omega_{ij}+\omega_c\) 始终为正，对应反旋跃迁。前一行来自 $\mathbf d^{(-)}\mathbf E^{(+)}+\mathbf d^{(+)}\mathbf E^{(-)}$，保持裸激发数；后一行来自 $\mathbf d^{(-)}\mathbf E^{(-)}+\mathbf d^{(+)}\mathbf E^{(+)}$，一次改变两个裸激发。这一区分是 RWA 的物理基础：在近共振 \(\omega_{ij}\approx\omega_c\) 下，前一行相位因子 \(\ee^{\pm\ii(\omega_{ij}-\omega_c)t}\) 缓慢变化，后一行 \(\ee^{\pm\ii(\omega_{ij}+\omega_c)t}\) 快速振荡，在旋转框架下平均消去。把\eqnrefs{eq:phys-oqs-cavity-field} 代入 $-\mathbf d\cdot\mathbf E$，四项的系数分别组织成\eqnrefs{eq:phys-oqs-cavity-rot-counter-coupling}。其中交换项的系数为
\[
-\ii\sqrt{\frac{\hbar\omega_c}{2\epsilon_0V_{\rm eff}}}
\,\mathbf d_{ij}\!\cdot\!\mathbf e_c f_c(\mathbf r_0)
=\hbar g_{ij},
\]
其厄密共轭系数为 $\hbar g_{ij}^*$，从而在删除 $g_{ij}^{\rm cr}$ 两个反旋项后得到\eqnrefs{eq:phys-oqs-hs-rwa}。RWA 的物理依据是：交换项相位 \(\ee^{\pm\ii(\omega_{ij}-\omega_c)t}\) 在共振附近缓慢变化，反旋项相位 \(\ee^{\pm\ii(\omega_{ij}+\omega_c)t}\) 快速振荡，在旋转框架下平均消去。这一平均对应于时间积分
\[
 \int_0^T\ee^{\pm\ii(\omega_{ij}+\omega_c)t}\dd t
 =\frac{\ee^{\pm\ii(\omega_{ij}+\omega_c)T}-1}{\pm\ii(\omega_{ij}+\omega_c)}
 \sim O\!\left(\frac{1}{\omega_{ij}+\omega_c}\right),
\]
在 \(T\gg1/(\omega_{ij}+\omega_c)\) 时趋于零。若用平均 Hamiltonian 或 Schrieffer--Wolff 消元处理反旋项，虚跃迁的能量分母为 $\hbar(\omega_{ij}+\omega_c)$，而矩阵元平方至多为 $\hbar^2|g_{ij}^{\rm cr}|^2(n+1)$，故频率修正量级正是 $|g_{ij}^{\rm cr}|^2(n+1)/(\omega_{ij}+\omega_c)$。这同时说明\eqnrefs{eq:phys-oqs-rwa-small-parameter} 控制一阶态混合，而频率修正从二阶开始；只有在 $|g_{ij}^{\rm cr}|=|g_{ij}|$ 的特殊模约定下，才可用同一个耦合模长估计两者。

作为应用举例，考虑二能级原子与单模腔场的共振耦合 \(\omega_{ij}=\omega_c\)：在 RWA 下得到 Jaynes--Cummings 模型 \(H_{\rm JC}=\hbar\omega_c a^\dagger a+\frac12\hbar\omega_{ij}\sigma_z+\hbar g(\sigma_{+}a+\sigma_{-}a^\dagger)\)，其激发数 \(\mathcal N=a^\dagger a+\sigma_{+}\sigma_{-}\) 严格守恒。这正是从经典偶极辐射到量子电动力学自发辐射的桥梁：在 Markov 极限下对腔场求迹，即得到原子自发辐射的 Lindblad 方程 \(\dot\rho=-\ii[H_A,\rho]/\hbar+\Gamma(\sigma_{-}\rho\sigma_{+}-\frac12\{\sigma_{+}\sigma_{-},\rho\})\)，其中 \(\Gamma=2\pi g^2D(\omega_{ij})\) 由 Fermi 黄金法则给出。该 Lindblad 方程的稳态为基态 \(\rho_{\rm ss}=|g\rangle\langle g|\)，激发态布居按 \(P_e(t)=P_e(0)e^{-\Gamma t}\) 指数衰减，对应自发辐射寿命 \(\tau=1/\Gamma\)。更一般地，若原子置于连续模式电磁场中而非单模腔，则耦合强度需替换为 \(g(\omega)\to\mathbf d_{ij}\cdot\mathbf e_\lambda\sqrt{\omega/(2\hbar\epsilon_0\pi c^3)}\)，对模式求和后给出真空自发辐射率 \(\Gamma=\omega_{ij}^3|\mathbf d_{ij}|^2/(3\pi\hbar\epsilon_0c^3)\)，这正是经典 Larmor 公式经量子化后的对应结果，体现了从经典辐射到量子电动力学的桥梁。反旋项在 ultrastrong-coupling 区（\(g/\omega_c\sim1\)）不可忽略：此时 RWA 失效，完整 Rabi 模型 \(H_{\rm Rabi}=\hbar\omega_c a^\dagger a+\frac12\hbar\omega_{ij}\sigma_z+\hbar g(\sigma_{+}+\sigma_{-})(a+a^\dagger)\) 需直接处理，激发数不守恒，基态含虚光子成分，对应于电路量子电动力学中超导量子比特与微波腔强耦合实验所观测到的 Bloch--Siegert 频移。

\emph{物理意义。}旋转波近似（RWA）的物理本质是"近共振时快振荡项的时间平均为零"：相互作用绘景中，共振项 \(\sigma_+ a\ee^{-i(\omega_{ij}-\omega_c)t}\) 缓变（相位近似冻结），反旋项 \(\sigma_+ a^\dagger\ee^{i(\omega_{ij}+\omega_c)t}\) 以 \(\omega_{ij}+\omega_c\) 快速振荡，长时间积分后平均为零。RWA 的适用条件 \(|g|/\omega_c\ll1\) 确保反旋项的累积效应远小于共振项。物理后果是激发数守恒——Jaynes--Cummings 模型中 \(\mathcal N=a^\dagger a+\sigma_+\sigma_-\) 严格守恒，使每个激发数子空间独立演化，这是腔 QED 中 Rabi 振荡、真空 Rabi 分裂和量子态制备的理论基础。RWA 在超强耦合区（\(g/\omega_c\sim1\)）失效：反旋项不可忽略，激发数不守恒，基态含虚光子成分（" dressed vacuum"），产生 Bloch--Siegert 频移和超辐射相变——这些效应在电路 QED 的超导量子比特-微波腔系统中已实验观测。从 RWA 到完整 Rabi 模型的过渡是量子光学从弱耦合到强耦合的标志性转变。
\end{derivation}

对这里的单模 RWA Hamiltonian~\eqnrefs{eq:phys-oqs-hs-rwa}，最常用的选择是
\begin{equation}
\boxed{
\mathcal N=a^\dagger a+\sum_i\nu_i\sigma_{ii}.}
\label{eq:phys-oqs-excitation-number-general}
\end{equation}
由
\[
[a^\dagger a,a]=-a,
\qquad
\left[\sum_k\nu_k\sigma_{kk},\sigma_{ij}\right]
=(\nu_i-\nu_j)\sigma_{ij}
\]
可得
\[
[\mathcal N,a\sigma_{ij}]
=(-1+\nu_i-\nu_j)a\sigma_{ij}.
\]
因此每条被保留的单光子边都必须满足
\begin{equation}
\boxed{\nu_i-\nu_j=1\qquad(g_{ij}\ne0).}
\label{eq:phys-oqs-grading-edge-condition}
\end{equation}
若这组线性约束存在解，就有 $[\mathcal N,H_S^{\rm RWA}]=0$。若不存在解，则不能强行使用固定激发流形。例如三条单光子边 $1\leftrightarrow2$、$2\leftrightarrow3$、$1\leftrightarrow3$ 同时保留时，会要求
\[
\nu_2-\nu_1=1,
\qquad
\nu_3-\nu_2=1,
\qquad
\nu_3-\nu_1=1,
\]
前两式相加却给 $\nu_3-\nu_1=2$，与第三式矛盾；这类闭环模型应直接使用完整 Liouvillian 或引入更合适的多模守恒荷。

当\eqnrefs{eq:phys-oqs-grading-edge-condition} 成立时，$\mathcal N$ 的 $N$ 本征子空间为
\begin{equation}
\boxed{
\mathcal H_N
=\operatorname{span}
\{|i,N-\nu_i\rangle:\;N-\nu_i\in\mathbb N_0\}.}
\label{eq:phys-oqs-manifold-span}
\end{equation}
其维数为
\begin{equation}
d_N=\#\{i:N-\nu_i\ge0\}.
\label{eq:phys-oqs-manifold-dimension}
\end{equation}
低激发区不能默认具有与高激发区相同的维数。若 $\min_i\nu_i=0$，则
\[
\mathcal H_0
=\operatorname{span}\{|i,0\rangle:\nu_i=0\}.
\]
二能级模型通常只有 $|g,0\rangle$；$\Lambda$ 型可有两个零等级下能级；V 型与 $\Xi$ 型通常只有最低态。随着 $N$ 增大，只有满足 $N\ge\nu_i$ 的原子态才进入该流形。



单模 RWA Hamiltonian 与激发流形确定以后，条件传播由\eqnrefs{eq:phys-oqs-manifold-propagator} 给出，无条件传播由\eqnrefs{eq:phys-oqs-duhamel-block} 给出。若经典驱动使生成元显含时间，则传播子必须相应采用时间有序形式；静态 Liouvillian 的普通矩阵指数只适用于时间无关生成元。

接下来考察二能级。

取
\[
|g\rangle,\ |e\rangle,
\qquad
\nu_g=0,\quad \nu_e=1,
\qquad
g_{eg}=g.
\]
由\eqnrefs{eq:phys-oqs-hs-rwa} 取 $\nu_g=0,\nu_e=1$，并记 $g_{eg}=g$。零激发扇区为
\[
\mathcal H_0^{(2)}=\operatorname{span}\{|g,0\rangle\},
\qquad
P_0=|g,0\rangle\langle g,0|,
\]
它没有可由 RWA 相互作用继续降低的伙伴态。对 $N\ge1$，
\[
\mathcal H_N^{(2)}
=\operatorname{span}\{|e,N-1\rangle,|g,N\rangle\},
\]
并且
\begin{equation}
P_N=|e,N-1\rangle\langle e,N-1|
+|g,N\rangle\langle g,N|.
\label{eq:phys-oqs-pn-two}
\end{equation}
在基底 $\mathcal B_N^{(2)}=\{|e,N-1\rangle,|g,N\rangle\}$ 中，逐个矩阵元做投影。对角元为
\[
\begin{aligned}
\langle e,N-1|H_S^{\rm RWA}|e,N-1\rangle
&=E_e+\hbar\omega_c\left(N-\frac12\right),\\
\langle g,N|H_S^{\rm RWA}|g,N\rangle
&=E_g+\hbar\omega_c\left(N+\frac12\right).
\end{aligned}
\]
非对角元使用 $a|N\rangle=\sqrt N|N-1\rangle$ 与 $\sigma_+|g\rangle=|e\rangle$：
\[
\begin{aligned}
\langle e,N-1|\hbar g a\sigma_+|g,N\rangle
&=\hbar g\sqrt N,\\
\langle g,N|\hbar g^*a^\dagger\sigma_-|e,N-1\rangle
&=\hbar g^*\sqrt N.
\end{aligned}
\]
其余交叉矩阵元为零。因此
\begin{equation}
\boxed{
H_{S,N}^{(2)}
=P_NH_S^{\rm RWA}P_N
=\begin{pmatrix}
E_e+\hbar\omega_c(N-\frac12)&\hbar g\sqrt N\\
\hbar g^*\sqrt N&E_g+\hbar\omega_c(N+\frac12)
\end{pmatrix}.}
\label{eq:phys-oqs-hn2-bare}
\end{equation}
\eqnrefs{eq:phys-oqs-hn2-bare} 还可以把“无关的共同能量”和真正控制两态混合的部分分开。定义
\begin{equation}
E_{0,N}=\frac{E_e+E_g}{2}+\hbar\omega_cN,
\qquad
\Delta=\omega_{eg}-\omega_c,
\label{eq:phys-oqs-two-detuning}
\end{equation}
则
\begin{equation}
\boxed{
H_{S,N}^{(2)}
=E_{0,N}\id(2)
+\frac{\hbar}{2}
\begin{pmatrix}
\Delta&2g\sqrt N\\
2g^*\sqrt N&-\Delta
\end{pmatrix}.}
\label{eq:phys-oqs-hn2-detuning-form}
\end{equation}
第一项只给整个流形共同相位；第二项才决定 dressed-state 混合、Rabi 频率和能级劈裂。因此后面出现的 $\sqrt N$ 增强并不是额外假设，而是投影 $a|N\rangle=\sqrt N|N-1\rangle$ 的直接结果。

加入\secref{sec:phys-oqs-gksl}得到的 Lamb shift 后，相干块统一为
\begin{equation}
\boxed{
H_N^{(2)}
=P_N(H_S^{\rm RWA}+H_{\rm LS})P_N
=H_{S,N}^{(2)}+P_NH_{\rm LS}P_N.}
\label{eq:phys-oqs-hn2-full}
\end{equation}
例如若 $H_{\rm LS}=\Delta_e^{\rm LS}|e\rangle\langle e|+\Delta_g^{\rm LS}|g\rangle\langle g|$，则只需在\eqnrefs{eq:phys-oqs-hn2-bare} 的两个对角元上分别加 $\Delta_e^{\rm LS}$ 与 $\Delta_g^{\rm LS}$；可观测的跃迁频率位移为 $(\Delta_e^{\rm LS}-\Delta_g^{\rm LS})/\hbar$。

下面把一般 GKSL 生成元特殊化为常用的“裸原子 + 裸腔模”耗散表示。若这些耗散通道来自对外部环境的微观消元，它们并不由\eqnrefs{eq:phys-oqs-hs-rwa} 自动推出。严格做法是以实际采用的耦合系统 Hamiltonian 求 Bohr 算符 $S_\mu(\omega)$，其 jump 一般作用在 dressed states 之间。只有当一组 dressed transition 的最大频率分裂 $\delta\omega_{\rm dress}$ 同时满足
\begin{equation}
\boxed{
\delta\omega_{\rm dress}\,\tau_B\ll1,
\qquad
\delta\omega_{\rm dress}\,\Delta t\lesssim1,
\qquad
\tau_B\ll\Delta t\ll\tau_R,}
\label{eq:phys-oqs-local-dissipator-condition}
\end{equation}
并且环境谱与耦合矩阵在该频率簇内近似常数时，有限粗粒化才既看不见谱函数的局部变化，也不把该簇内部频率完全 secular 分开；此时保留完整交叉块后，簇内 dressed Bohr 算符之和可重新组合成 $a$、$\sigma_-$ 等裸算符。若任一条件失效，就应回到 dressed $S_\mu(\omega)$ 与相应 Kossakowski 块，而不能把“裸耗散”当作与 Hamiltonian 无关的普适形式。

在满足上述局域耗散条件，或把下列通道直接作为现象学 GKSL 模型定义时，考虑原子上下跃迁、原子纯退相干、腔的损失与热激发以及腔频率噪声：
\[
C_{e\to g}=\sqrt{\Gamma_{e\to g}}\,\sigma_-,
\qquad
C_{g\to e}=\sqrt{\Gamma_{g\to e}}\,\sigma_+,
\]
\[
C_A^\varphi=\sqrt{\frac{\gamma_\varphi}{2}}\,\sigma_z,
\qquad
C_c^-=\sqrt{\kappa_\downarrow}\,a,
\qquad
C_c^+=\sqrt{\kappa_\uparrow}\,a^\dagger,
\qquad
C_c^\varphi=\sqrt{\kappa_\varphi}\,a^\dagger a.
\]
这里的 $\Gamma_{e\to g}$ 与 $\Gamma_{g\to e}$ 是\secref{sec:phys-oqs-gksl}一般频率块在“只有一条 $e\leftrightarrow g$ 跃迁”时的标量极限；热环境中可直接使用\eqnrefs{eq:phys-oqs-thermal-transition-rates}，而不是重新定义另一套速率符号。

这些 collapse operators 对流形的作用也可以在构造 hazard 之前直接写出。对 $N\ge1$，
\begin{align}
J_{e\to g}^{(N)}
&=\sqrt{\Gamma_{e\to g}}\,
|g,N-1\rangle\langle e,N-1|,
& s_{e\to g}&=1,\\
J_{g\to e}^{(N)}
&=\sqrt{\Gamma_{g\to e}}\,
|e,N\rangle\langle g,N|,
& s_{g\to e}&=-1,
\end{align}
而腔 jump 为
\begin{align}
J_{c,-}^{(N)}
={}&\sqrt{\kappa_\downarrow}\big[
\sqrt{N-1}|e,N-2\rangle\langle e,N-1|
+\sqrt N|g,N-1\rangle\langle g,N|\big],\\
J_{c,+}^{(N)}
={}&\sqrt{\kappa_\uparrow}\big[
\sqrt N|e,N\rangle\langle e,N-1|
+\sqrt{N+1}|g,N+1\rangle\langle g,N|\big].
\end{align}
纯退相干不改变 $N$：
\begin{align}
J_{A,\varphi}^{(N)}
&=\sqrt{\frac{\gamma_\varphi}{2}}
\big(|e,N-1\rangle\langle e,N-1|-|g,N\rangle\langle g,N|\big),\\
J_{c,\varphi}^{(N)}
&=\sqrt{\kappa_\varphi}
\big[(N-1)|e,N-1\rangle\langle e,N-1|
+N|g,N\rangle\langle g,N|\big].
\end{align}
这些 ket-bra 形式在低激发边界的含义是：仅保留光子数指标非负的项；任何含负 Fock 指标的项都从算符和中省略，而不是把不存在的 ket 形式上乘以零。这样 $N=0,1$ 的实际流形维数始终由投影 $P_N$ 决定，不需要人为补成固定维方阵。

构造 hazard 矩阵时先计算
\[
\begin{aligned}
C_{e\to g}^\dagger C_{e\to g}&=\Gamma_{e\to g}|e\rangle\langle e|,\\
C_{g\to e}^\dagger C_{g\to e}&=\Gamma_{g\to e}|g\rangle\langle g|,\\
(C_A^\varphi)^\dagger C_A^\varphi&=\frac{\gamma_\varphi}{2}\id(),\\
(C_c^-)^\dagger C_c^-&=\kappa_\downarrow a^\dagger a,\\
(C_c^+)^\dagger C_c^+&=\kappa_\uparrow(a^\dagger a+1),\\
(C_c^\varphi)^\dagger C_c^\varphi&=\kappa_\varphi(a^\dagger a)^2.
\end{aligned}
\]
定义
\[
\boxed{
\ell_{e,N}=\langle e,N-1|K_N|e,N-1\rangle,
\qquad
\ell_{g,N}=\langle g,N|K_N|g,N\rangle,}
\]
其中
\[
K_N=P_N\left(\sum_qC_q^\dagger C_q\right)P_N.
\]
由于 $a^\dagger a|e,N-1\rangle=(N-1)|e,N-1\rangle$，逐项取期望值得
\[
\boxed{
\ell_{e,N}
=\Gamma_{e\to g}+\frac{\gamma_\varphi}{2}
+(N-1)\kappa_\downarrow+N\kappa_\uparrow
+(N-1)^2\kappa_\varphi.}
\]
同理，$a^\dagger a|g,N\rangle=N|g,N\rangle$，因此
\[
\boxed{
\ell_{g,N}
=\Gamma_{g\to e}+\frac{\gamma_\varphi}{2}
+N\kappa_\downarrow+(N+1)\kappa_\uparrow
+N^2\kappa_\varphi.}
\]
两者分别是基矢 $|e,N-1\rangle$ 与 $|g,N\rangle$ 在单位时间内离开 no-jump 分支的总 hazard。于是
\begin{equation}
\boxed{
K_N^{(2)}=\begin{pmatrix}\ell_{e,N}&0\\0&\ell_{g,N}\end{pmatrix},
\qquad
M_N^{(2)}=H_N^{(2)}-\frac{\ii\hbar}{2}K_N^{(2)}.}
\label{eq:phys-oqs-mn2}
\end{equation}

写
\[
M_N^{(2)}=
\begin{pmatrix}
m_{ee,N}&m_{eg,N}\\
m_{ge,N}&m_{gg,N}
\end{pmatrix}.
\]
其两个复极点由二次特征方程
\[
(z-m_{ee,N})(z-m_{gg,N})-m_{eg,N}m_{ge,N}=0
\]
得到。令
\[
\bar m_N=\frac{m_{ee,N}+m_{gg,N}}2,
\qquad
\chi_N=\sqrt{\left(\frac{m_{ee,N}-m_{gg,N}}2\right)^2+m_{eg,N}m_{ge,N}},
\]
则
\[
\boxed{z_{\pm,N}=\bar m_N\pm\chi_N.}
\]
把 $z_{\pm,N}$ 代入普适谱公\eqnrefs{eq:phys-oqs-spectral-propagator} 即得条件传播子。对非简并情形，可写成
\[
U_N(t)=\ee^{-\ii z_{+,N}t/\hbar}
\frac{M_N^{(2)}-z_{-,N}\id(2)}{z_{+,N}-z_{-,N}}
+\ee^{-\ii z_{-,N}t/\hbar}
\frac{M_N^{(2)}-z_{+,N}\id(2)}{z_{-,N}-z_{+,N}}.
\]
闭合极限 $C_q=0$、$H_{\rm LS}=0$ 下，\eqnrefs{eq:phys-oqs-hn2-detuning-form} 直接给出
\begin{equation}
\boxed{\Omega_N=\sqrt{\Delta^2+4|g|^2N}.}
\label{eq:phys-oqs-jc-generalized-rabi}
\end{equation}
把\eqnrefs{eq:phys-oqs-hn2-detuning-form} 的两条谱投影代入普适传播子~\eqnrefs{eq:phys-oqs-spectral-propagator}，分别乘上相位因子后得到
\begin{equation}
\boxed{
U_N^{\rm JC}(t)
=\ee^{-\ii E_{0,N}t/\hbar}
\left[
\cos\frac{\Omega_Nt}{2}\,\id(2)
-\frac{\ii\sin(\Omega_Nt/2)}{\Omega_N}
\begin{pmatrix}
\Delta&2g\sqrt N\\
2g^*\sqrt N&-\Delta
\end{pmatrix}
\right].}
\label{eq:phys-oqs-jc-propagator}
\end{equation}
例如初态为 $|e,N-1\rangle$ 时，
\begin{align}
c_e(t)
&=\ee^{-\ii E_{0,N}t/\hbar}
\left[\cos\frac{\Omega_Nt}{2}
-\ii\frac{\Delta}{\Omega_N}\sin\frac{\Omega_Nt}{2}\right],\\
c_g(t)
&=-\ii\,\ee^{-\ii E_{0,N}t/\hbar}
\frac{2g^*\sqrt N}{\Omega_N}
\sin\frac{\Omega_Nt}{2}.
\end{align}
因此共振 $\Delta=0$ 时的 population 交换频率就是 $2|g|\sqrt N$；失谐只改变同一普适二维传播子的极点和混合角。

密度矩阵块的物理内容可直接写成 Fock 矩阵元：
\[
\boxed{
R_{NM}^{(2)}=
\begin{pmatrix}
\rho_{N-1,M-1}^{ee}&\rho_{N-1,M}^{eg}\\
\rho_{N,M-1}^{ge}&\rho_{N,M}^{gg}
\end{pmatrix},
\qquad
\rho_{nm}^{\alpha\beta}=\langle\alpha,n|\rho|\beta,m\rangle.}
\]
记 recycle 源矩阵为
\[
\mathcal S_{NM}^{(2)}=
\sum_qJ_q^{(N+s_q)}R_{N+s_q,M+s_q}^{(2)}J_q^{(M+s_q)\dagger}
\equiv
\begin{pmatrix}s_{ee}&s_{eg}\\s_{ge}&s_{gg}\end{pmatrix}_{NM}.
\]
把 $M_N^{(2)}R_{NM}^{(2)}-R_{NM}^{(2)}M_M^{(2)\dagger}$ 逐项相乘，得到四个直接可计算的方程：
\[
\boxed{
\begin{aligned}
\dot\rho_{N-1,M-1}^{ee}
={}&-\frac{\ii}{\hbar}
\left[(m_{ee,N}-m_{ee,M}^*)\rho_{N-1,M-1}^{ee}
+m_{eg,N}\rho_{N,M-1}^{ge}
-m_{eg,M}^*\rho_{N-1,M}^{eg}\right]+s_{ee}^{NM},\\
\dot\rho_{N-1,M}^{eg}
={}&-\frac{\ii}{\hbar}
\left[(m_{ee,N}-m_{gg,M}^*)\rho_{N-1,M}^{eg}
+m_{eg,N}\rho_{N,M}^{gg}
-m_{ge,M}^*\rho_{N-1,M-1}^{ee}\right]+s_{eg}^{NM},\\
\dot\rho_{N,M-1}^{ge}
={}&-\frac{\ii}{\hbar}
\left[(m_{gg,N}-m_{ee,M}^*)\rho_{N,M-1}^{ge}
+m_{ge,N}\rho_{N-1,M-1}^{ee}
-m_{eg,M}^*\rho_{N,M}^{gg}\right]+s_{ge}^{NM},\\
\dot\rho_{N,M}^{gg}
={}&-\frac{\ii}{\hbar}
\left[(m_{gg,N}-m_{gg,M}^*)\rho_{N,M}^{gg}
+m_{ge,N}\rho_{N-1,M}^{eg}
-m_{ge,M}^*\rho_{N,M-1}^{ge}\right]+s_{gg}^{NM}.
\end{aligned}}
\]

各 jump 对 recycle 源的作用可以直接从投影矩阵读出。对原子跃迁，
\[
\mathcal S_{NM}^{e\to g}
=\Gamma_{e\to g}
\begin{pmatrix}0&0\\0&\rho_{N,M}^{ee}\end{pmatrix},
\qquad
\mathcal S_{NM}^{g\to e}
=\Gamma_{g\to e}
\begin{pmatrix}\rho_{N-1,M-1}^{gg}&0\\0&0\end{pmatrix}.
\]
对腔损失，
\[
\mathcal S_{NM}^{c-}=\kappa_\downarrow
\begin{pmatrix}
\sqrt{NM}\,\rho_{N,M}^{ee}&\sqrt{N(M+1)}\,\rho_{N,M+1}^{eg}\\
\sqrt{(N+1)M}\,\rho_{N+1,M}^{ge}&\sqrt{(N+1)(M+1)}\,\rho_{N+1,M+1}^{gg}
\end{pmatrix},
\]
而腔热激发给
\[
\mathcal S_{NM}^{c+}=\kappa_\uparrow
\begin{pmatrix}
\sqrt{(N-1)(M-1)}\,\rho_{N-2,M-2}^{ee}&\sqrt{(N-1)M}\,\rho_{N-2,M-1}^{eg}\\
\sqrt{N(M-1)}\,\rho_{N-1,M-2}^{ge}&\sqrt{NM}\,\rho_{N-1,M-1}^{gg}
\end{pmatrix}.
\]
纯退相干是 $s_q=0$ 的同流形 recycle：
\[
\mathcal S_{NM}^{A\varphi}=\frac{\gamma_\varphi}{2}
\begin{pmatrix}
\rho_{N-1,M-1}^{ee}&-\rho_{N-1,M}^{eg}\\
-\rho_{N,M-1}^{ge}&\rho_{N,M}^{gg}
\end{pmatrix},
\]
\[
\mathcal S_{NM}^{c\varphi}=\kappa_\varphi
\begin{pmatrix}
(N-1)(M-1)\rho_{N-1,M-1}^{ee}&(N-1)M\rho_{N-1,M}^{eg}\\
N(M-1)\rho_{N,M-1}^{ge}&NM\rho_{N,M}^{gg}
\end{pmatrix}.
\]
这些矩阵公式只对相应 Fock 指标非负的投影子空间成立；越过低激发边界的项由流形投影 $P_N$ 或 $P_M$ 直接缺席，而不是额外规定不存在的矩阵元等于零。这组方程把 no-jump 漂移与 recycle 注入分别保留，因此可以逐项追踪 population 与 coherence 的来源。

接下来考察三能级。

取三个物质能级 $|1\rangle,|2\rangle,|3\rangle$。三能级动力学先由一般能级指标、一般耦合矩阵元和一般耗散通道建立；$\Lambda$、V、$\Xi$ 三种构形随后由允许跃迁的参数约束直接生成。

把一般结构特殊化到量子单模耦合的三能级系统，可以得到以下块矩阵。

给每个能级分配整数激发等级 $\nu_i$，并沿用\eqnrefs{eq:phys-oqs-manifold-span} 的总激发算符。固定流形中的基矢写为
\[
|i;N\rangle=|i,N-\nu_i\rangle,
\qquad N-\nu_i\ge0,
\]
并定义 $I_N=\{i:N-\nu_i\ge0\}$。投影算符与固定流形为
\begin{equation}
\boxed{
P_N=\sum_{i\in I_N}|i;N\rangle\langle i;N|,
\qquad
\mathcal H_N^{(3)}=\operatorname{span}\{|i;N\rangle:i\in I_N\}.}
\label{eq:phys-oqs-pn-three}
\end{equation}
因此低激发区的维数是 $d_N=|I_N|$；只有当三个基矢都存在时才有 $d_N=3$。

把\eqnrefs{eq:phys-oqs-hs-rwa} 投影到 $\mathcal H_N^{(3)}$。先算对角元。由于 $H_A$ 与 $H_c$ 都不改变原子态和光子数，
\begin{align}
\langle i;N|P_N(H_A+H_c)P_N|i;N\rangle
&=\langle i,N-\nu_i|H_A+H_c|i,N-\nu_i\rangle\\
&=E_i+\hbar\omega_c\left(N-\nu_i+\frac12\right).
\end{align}
再算非对角元。若 $\nu_i=\nu_j+1$，\eqnrefs{eq:phys-oqs-hs-rwa} 中连接 $|j;N\rangle$ 与 $|i;N\rangle$ 的项是 $\hbar g_{ij}a\sigma_{ij}$，于是
\begin{align}
\langle i;N|\hbar g_{ij}a\sigma_{ij}|j;N\rangle
&=\hbar g_{ij}
\langle i|\sigma_{ij}|j\rangle
\langle N-\nu_i|a|N-\nu_j\rangle\\
&=\hbar g_{ij}
\sqrt{N-\nu_j}\,
\delta_{N-\nu_i,N-\nu_j-1}\\
&=\hbar g_{ij}\sqrt{N-\nu_j}.
\end{align}
若 $\nu_j=\nu_i+1$，则由 Hermitian 共轭项得到反向矩阵元。故
\begin{equation}
\boxed{
(H_{S,N}^{(3)})_{ij}=
\begin{cases}
E_i+\hbar\omega_c(N-\nu_i+\frac12),&i=j,\\
\hbar g_{ij}\sqrt{N-\nu_j},&\nu_i=\nu_j+1,\\
\hbar g_{ji}^*\sqrt{N-\nu_i},&\nu_j=\nu_i+1,\\
0,&\text{其余情形}.
\end{cases}}
\label{eq:phys-oqs-hn3-elements}
\end{equation}
这一步给出的不是某一种三能级构形，而是任意激发等级 $\nu_i$ 和任意允许单光子边 $g_{ij}$ 的统一投影结果。

开放系统还需要把 collapse operators 投影成 jump blocks。对任意原子跃迁 $i\to j$，取
\[
C_{i\to j}=\sqrt{\Gamma_{i\to j}}\,\sigma_{ji},
\qquad
s_{i\to j}=\nu_i-\nu_j.
\]
由\eqnrefs{eq:phys-oqs-pn-three} 逐项作用，
\begin{align}
P_{N-s_{i\to j}}C_{i\to j}P_N|k;N\rangle
&=\sqrt{\Gamma_{i\to j}}\,\delta_{ik}
P_{N-s_{i\to j}}|j,N-\nu_i\rangle\\
&=\sqrt{\Gamma_{i\to j}}\,\delta_{ik}|j;N-s_{i\to j}\rangle.
\end{align}
因此
\begin{equation}
\boxed{
J_{i\to j}^{(N)}
=\sqrt{\Gamma_{i\to j}}\,
|j;N-s_{i\to j}\rangle\langle i;N|.}
\label{eq:phys-oqs-three-atomic-jump}
\end{equation}
这说明原子 jump 在固定流形语言中只是把原子指标 $i$ 改成 $j$，同时把总激发标签从 $N$ 改成 $N-s_{i\to j}$。

对腔损失 $C_{c,-}=\sqrt{\kappa_\downarrow}\,a$，有 $s_{c,-}=1$。因为
\[
a|i;N\rangle
=\sqrt{N-\nu_i}\,|i;N-1\rangle,
\]
故
\begin{equation}
\boxed{
J_{c,-}^{(N)}
=\sqrt{\kappa_\downarrow}
\sum_{i\in I_N\cap I_{N-1}}
\sqrt{N-\nu_i}\,
|i;N-1\rangle\langle i;N|.}
\label{eq:phys-oqs-three-cavity-jump-down}
\end{equation}
同理，对热激发 $C_{c,+}=\sqrt{\kappa_\uparrow}\,a^\dagger$，$s_{c,+}=-1$，
\begin{equation}
\boxed{
J_{c,+}^{(N)}
=\sqrt{\kappa_\uparrow}
\sum_{i\in I_N\cap I_{N+1}}
\sqrt{N-\nu_i+1}\,
|i;N+1\rangle\langle i;N|.}
\label{eq:phys-oqs-three-cavity-jump-up}
\end{equation}
对零频纯退相干，取 Hermitian 对角 collapse operator
\begin{equation}
C_{\varphi,\alpha}
=\sqrt{\kappa_\alpha}\sum_iq_{\alpha i}\sigma_{ii},
\qquad q_{\alpha i}\in\mathbb R.
\label{eq:phys-oqs-three-dephasing-collapse}
\end{equation}
它满足 $s_{\varphi,\alpha}=0$，所以不改变流形。直接取任意原子 coherence 的矩阵元：
\begin{align}
\langle i|\mathcal D[C_{\varphi,\alpha}]\rho|j\rangle
={}&\kappa_\alpha q_{\alpha i}q_{\alpha j}\rho_{ij}
-\frac{\kappa_\alpha}{2}
(q_{\alpha i}^2+q_{\alpha j}^2)\rho_{ij}\nonumber\\
={}&-\frac{\kappa_\alpha}{2}
(q_{\alpha i}-q_{\alpha j})^2\rho_{ij}.
\end{align}
因此一般多通道纯退相干率为
\begin{equation}
\boxed{
\gamma_{ij}^{\varphi}
=\frac12\sum_\alpha\kappa_\alpha
(q_{\alpha i}-q_{\alpha j})^2.}
\label{eq:phys-oqs-general-pure-dephasing-rate}
\end{equation}
这就是\secref{sec:phys-oqs-gksl}零频谱结论在三能级矩阵元语言中的直接实现：只有不同能级感受到不同的零频噪声位移时，相对相位才会衰减。

由这些 jump 可直接计算 $K_N=P_N\sum_qC_q^\dagger C_qP_N$。例如相互独立的原子跃迁给
\[
P_N\sum_{i\ne j}C_{i\to j}^\dagger C_{i\to j}P_N
=\sum_{i\in I_N}\Gamma_i^{\rm out}|i;N\rangle\langle i;N|,
\qquad
\Gamma_i^{\rm out}=\sum_{j\ne i}\Gamma_{i\to j},
\]
腔损失与热激发分别给
\[
\kappa_\downarrow\sum_{i\in I_N}(N-\nu_i)|i;N\rangle\langle i;N|,
\qquad
\kappa_\uparrow\sum_{i\in I_N}(N-\nu_i+1)|i;N\rangle\langle i;N|.
\]
上式只对应“各原子跃迁已经彼此可分辨”的对角耗散情形。若同一 Bohr 频率块内有多条不可分辨跃迁，应保留完整 Kossakowski 矩阵。用跃迁标签
\[
\alpha=(i_\alpha\to j_\alpha),
\qquad
F_\alpha=\sigma_{j_\alpha i_\alpha}
\]
表示该频率块中的系统算符，则耗散部分为
\begin{equation}
\mathcal D_{\Gamma^{(\omega)}}[\rho]
=\sum_{\alpha,\beta\in\omega}
\Gamma_{\alpha\beta}^{(\omega)}
\left(F_\beta\rho F_\alpha^\dagger
-\frac12\{F_\alpha^\dagger F_\beta,\rho\}\right).
\label{eq:phys-oqs-general-transition-kossakowski}
\end{equation}
相应 hazard 算符不是简单的标量速率之和，而是
\begin{equation}
K^{(\omega)}
=\sum_{\alpha,\beta\in\omega}
\Gamma_{\alpha\beta}^{(\omega)}F_\alpha^\dagger F_\beta.
\label{eq:phys-oqs-kossakowski-hazard}
\end{equation}
由于
\[
F_\alpha^\dagger F_\beta
=\delta_{j_\alpha j_\beta}
\sigma_{i_\alpha i_\beta},
\]
如果两条近简并跃迁具有同一个末态，$K^{(\omega)}$ 会在它们的初态之间产生非对角 hazard；如果它们具有同一个初态，则 recycle 项
$F_\beta\rho F_\alpha^\dagger$ 可以在不同末态之间直接生成 coherence。二者都来自同一个 $\Gamma_{\alpha\beta}^{(\omega)}$，并不是 V 型或 $\Lambda$ 型才额外定义的特殊机制。

把交叉耗散继续投影到固定激发流形之前，还必须检查它是否保留前面使用的荷分级。若
\[
[\mathcal N,F_\alpha]=-s_\alpha F_\alpha,
\]
则交叉 recycle 项 $F_\beta\rho F_\alpha^\dagger$ 会把左右荷标签分别改变 $s_\beta$ 与 $s_\alpha$。因此只有在每个非零 Kossakowski 元都满足
\begin{equation}
\boxed{
\Gamma_{\alpha\beta}^{(\omega)}\ne0
\quad\Longrightarrow\quad
s_\alpha=s_\beta
}
\label{eq:phys-oqs-kossakowski-grading-compatibility}
\end{equation}
时，$N-M$ 才保持不变，并且该频率块可在固定的 $s$ 子块内酉对角化成具有确定荷步长的 collapse operators。此时只需计算
\[
K_N^{(\omega)}=P_NK^{(\omega)}P_N,
\qquad
P_{N-s}F_\alpha P_N,
\]
再代入一般块方程~\eqnrefs{eq:phys-oqs-general-block-master}。若\eqnrefs{eq:phys-oqs-kossakowski-grading-compatibility} 不成立，物理 GKSL 方程本身没有问题，但当前 $\mathcal N$ 不再给出 Liouvillian 对称分级，必须在完整 Liouville 空间传播或寻找与耗散共同兼容的另一守恒荷。只有当 $\Gamma_{\alpha\beta}^{(\omega)}$ 进一步在所选跃迁基底中对角时，才可把它简化成彼此独立的标量 $\Gamma_{i\to j}$。

于是三能级固定流形中的三个核心矩阵统一为
\begin{equation}
\boxed{
H_N^{(3)}=P_N(H_S^{\rm RWA}+H_{\rm LS})P_N,
\qquad
K_N^{(3)}=P_N\left(\sum_qC_q^\dagger C_q\right)P_N,
\qquad
M_N^{(3)}=H_N^{(3)}-\frac{\ii\hbar}{2}K_N^{(3)}.}
\label{eq:phys-oqs-three-m}
\end{equation}
当 $d_N=3$ 时写 $M_N^{(3)}=(m_{ij}^{(N)})$。三个条件极点由
\begin{equation}
\boxed{
p_N(z)=\det(z\id(3)-M_N^{(3)})
=z^3-\tau_{1,N}z^2+\tau_{2,N}z-\tau_{3,N}=0,}
\label{eq:phys-oqs-three-cubic}
\end{equation}
其中
\[
\tau_{1,N}=\Tr M_N^{(3)},\qquad
\tau_{2,N}=\frac12\left[(\Tr M_N^{(3)})^2-\Tr((M_N^{(3)})^2)\right],\qquad
\tau_{3,N}=\det M_N^{(3)}.
\]
求出三个根以后直接使用\eqnrefs{eq:phys-oqs-spectral-propagator}；矩阵指数的证明已经在\secref{sec:phys-oqs-unified}完成，此处只负责把具体 $M_N^{(3)}$ 代入。

无条件动力学同样不需要新方法。令
\[
R_{NM}^{(3)}=(r_{ij}^{NM}),
\qquad
\mathcal S_{NM}^{(3)}=(s_{ij}^{NM}),
\]
把\eqnrefs{eq:phys-oqs-duhamel-block} 对应的微分方程逐矩阵元展开，先有
\[
(M_NR_{NM})_{ij}=\sum_{k=1}^{3}m_{ik}^{(N)}r_{kj}^{NM},
\qquad
(R_{NM}M_M^\dagger)_{ij}
=\sum_{k=1}^{3}r_{ik}^{NM}m_{jk}^{(M)*},
\]
而 recycle 源项是
\[
s_{ij}^{NM}
=\left[
\sum_qJ_q^{(N+s_q)}R_{N+s_q,M+s_q}^{(3)}J_q^{(M+s_q)\dagger}
\right]_{ij}.
\]
因此九个矩阵元共享同一个母式
\begin{equation}
\boxed{
\dot r_{ij}^{NM}
=-\frac{\ii}{\hbar}
\left[
\sum_{k=1}^3m_{ik}^{(N)}r_{kj}^{NM}
-\sum_{k=1}^3r_{ik}^{NM}m_{jk}^{(M)*}
\right]
+s_{ij}^{NM},
\qquad i,j=1,2,3.}
\label{eq:phys-oqs-three-block-elements}
\end{equation}
\eqnrefs{eq:phys-oqs-three-block-elements} 与\eqnrefs{eq:phys-oqs-three-atomic-jump}--\eqnrefs{eq:phys-oqs-three-cavity-jump-up} 已经构成量子单模三能级系统的完整无条件计算规则。

若耦合来自经典驱动，则先把显式时间相位吸收到合适的旋转框架中。

经典激光驱动同样从电偶极相互作用得到。更一般地，设外场含有若干频率分量
\begin{equation}
\mathbf E^{(+)}(t)
=\frac12\sum_r\boldsymbol{\mathcal E}_r
\ee^{-\ii\omega_rt},
\qquad
\mathbf E^{(-)}=(\mathbf E^{(+)})^\dagger,
\label{eq:phys-oqs-multitone-classical-field}
\end{equation}
并对每个场分量与每条跃迁分别定义共旋与反旋偶极投影
\begin{equation}
\boxed{
\Omega_{ij}^{(r)}
=\frac{\mathbf d_{ij}\cdot\boldsymbol{\mathcal E}_r}{\hbar},
\qquad
\Omega_{ij,{\rm cr}}^{(r)}
=\frac{\mathbf d_{ij}\cdot\boldsymbol{\mathcal E}_r^*}{\hbar}.}
\label{eq:phys-oqs-general-rabi-frequency}
\end{equation}
两者都是角频率。对线偏振并选择实偶极基时二者模长相同；对一般复偏振或有角动量选择规则的跃迁，不能把它们预先等同。把单个场分量 $r$ 与单条跃迁 $i\leftrightarrow j$ 单独取出，在相互作用绘景中，共旋项与反旋项分别带有频率
\[
\Delta_{ij}^{(r)}=\omega_{ij}-\omega_r,
\qquad
\Sigma_{ij}^{(r)}=\omega_{ij}+\omega_r.
\]
因此丢弃反旋转项的 optical RWA 由
\begin{equation}
\boxed{
\epsilon_{\rm optRWA}^{(ij,r)}
=\frac{|\Omega_{ij,{\rm cr}}^{(r)}|}{2|\Sigma_{ij}^{(r)}|}
\ll1}
\label{eq:phys-oqs-classical-optical-rwa-control}
\end{equation}
控制：反旋转分量给态矢造成 $O(\epsilon_{\rm optRWA}^{(ij,r)})$ 的快速微运动，并在二阶产生
\[
\delta\omega_{\rm BS}^{(ij,r)}
=O\!\left(\frac{|\Omega_{ij,{\rm cr}}^{(r)}|^2}{|\Sigma_{ij}^{(r)}|}\right)
\]
的 Bloch--Siegert 型频移。这个条件只允许删除反旋转项；它并不自动允许把一个仍然旋转、但远失谐的场分量从 Hamiltonian 中删除。若某个 $r$ 不收入后面的 $\mathcal R_{ij}$，还必须另行满足
\begin{equation}
\boxed{
\epsilon_{\rm off}^{(ij,r)}
=\frac{|\Omega_{ij}^{(r)}|}{2|\Delta_{ij}^{(r)}|}
\ll1,
\qquad \Delta_{ij}^{(r)}\ne0,}
\label{eq:phys-oqs-classical-offresonant-control}
\end{equation}
此时被消去的旋转项仍会留下
$O(|\Omega_{ij}^{(r)}|^2/|\Delta_{ij}^{(r)}|)$ 的 AC-Stark 位移；若这一二阶位移达到所需精度，就必须把它保留在有效 Hamiltonian 中。接近多光子共振、强驱动或 $|\Delta_{ij}^{(r)}|$ 与 $|\Omega_{ij}^{(r)}|$ 同阶时，\eqnrefs{eq:phys-oqs-classical-offresonant-control} 失效，不能通过“没有直接共振”删除该频率分量。

把 $-\mathbf d\cdot\mathbf E(t)$ 按照\eqnrefs{eq:phys-oqs-dipole-positive-negative} 所固定的正负频率约定分解；在\eqnrefs{eq:phys-oqs-classical-optical-rwa-control} 控制下先删除反旋转项，再按\eqnrefs{eq:phys-oqs-classical-offresonant-control} 删除确实可微扰消去的远失谐旋转项，便得到
\[
H_d^{\rm RWA}(t)
=-\frac{\hbar}{2}
\sum_{E_i>E_j}\sum_{r\in\mathcal R_{ij}}
\left[
\Omega_{ij}^{(r)}\ee^{-\ii\omega_rt}\sigma_{ij}
+\Omega_{ij}^{(r)*}\ee^{+\ii\omega_rt}\sigma_{ji}
\right],
\]
其中 $\mathcal R_{ij}$ 表示在所取光学 RWA 中保留、能够驱动该跃迁的频率分量集合。若每条允许边只保留一个主导激光，记该频率为 $\omega_{ij}^{d}$，并简写 $\Omega_{ij}^{(r)}\to\Omega_{ij}$。于是实验室系 Hamiltonian 才退化为后面三个例子使用的形式
\begin{equation}
\boxed{
H_{\rm lab}(t)
=\sum_iE_i\sigma_{ii}
-\frac{\hbar}{2}\sum_{E_i>E_j}
\left(
\Omega_{ij}\ee^{-\ii\omega_{ij}^{d}t}\sigma_{ij}
+\Omega_{ij}^*\ee^{+\ii\omega_{ij}^{d}t}\sigma_{ji}
\right).}
\label{eq:phys-oqs-three-drive-lab}
\end{equation}
因此\eqnrefs{eq:phys-oqs-three-drive-lab} 不是三能级专用假设，而是多频经典场在“每条边只保留一个近共振分量”后的记号简化。开放系统的相干部分实际是 $H_{\rm lab}(t)+H_{\rm LS}$；下面先变换 $H_{\rm lab}$，随后再说明 $H_{\rm LS}$ 与 dissipator 如何一起变换。

现在选择三个旋转频率 $\chi_i$，定义
\[
U_R(t)=\exp\!\left[-\ii t\sum_i\chi_i\sigma_{ii}\right],
\qquad
\rho_R=U_R^\dagger\rho U_R.
\]
由 $U_R^\dagger\sigma_{ij}U_R=
\ee^{\ii(\chi_i-\chi_j)t}\sigma_{ij}$ 以及
$-\ii\hbar U_R^\dagger\dot U_R=-\hbar\sum_i\chi_i\sigma_{ii}$，旋转框架 Hamiltonian 逐项变为
\begin{align}
H_R(t)
&=U_R^\dagger H_{\rm lab}(t)U_R
-\ii\hbar U_R^\dagger\dot U_R\\
&=\sum_i(E_i-\hbar\chi_i)\sigma_{ii}
-\frac{\hbar}{2}\sum_{E_i>E_j}
\left[
\Omega_{ij}
\ee^{-\ii[\omega_{ij}^{d}-(\chi_i-\chi_j)]t}\sigma_{ij}
+\mathrm{h.c.}
\right].
\label{eq:phys-oqs-three-drive-transform}
\end{align}
对一组可以同时静止化的驱动边，选择
\[
\chi_i-\chi_j=\omega_{ij}^{d}.
\]
再定义裸旋转框架失谐
\[
\delta_i^{(0)}=\frac{E_i}{\hbar}-\chi_i,
\qquad
\Delta_{ij}^{(0)}=\delta_i^{(0)}-\delta_j^{(0)}.
\]
给所有 $\delta_i^{(0)}$ 同时加同一个常数只给 Hamiltonian 增加单位阵倍数，不影响密度矩阵动力学，因此可以任取一个能级作为零点。在暂不计 Lamb shift 时，相干 Hamiltonian 为
\[
H_{\rm rot}^{(0)}
=\hbar\sum_i\delta_i^{(0)}\sigma_{ii}
-\frac{\hbar}{2}\sum_{E_i>E_j}
\left(\Omega_{ij}\sigma_{ij}+\Omega_{ij}^*\sigma_{ji}\right).
\]
若 $i\leftrightarrow j$ 本身被频率 $\omega_{ij}^{d}$ 的激光直接驱动，则
\[
\Delta_{ij}^{(0)}
=\frac{E_i-E_j}{\hbar}-\omega_{ij}^{d}
=\omega_{ij}-\omega_{ij}^{d},
\]
即通常的裸单光子失谐。若三条边都被独立驱动，则三个条件 $\chi_i-\chi_j=\omega_{ij}^{d}$ 只有在闭环频率满足
\[
\omega_{31}^{d}=\omega_{32}^{d}+\omega_{21}^{d}
\]
时才能同时成立；否则必有一条边保留多光子失谐对应的显式时间相位。这一闭环条件只由驱动相位决定，与后面是否加入 Lamb shift 无关。

旋转框架必须作用在整个主方程上，而不仅是 Hamiltonian。若实验室系写成
\[
\dot\rho
=-\frac{\ii}{\hbar}[H_{\rm lab}(t)+H_{\rm LS},\rho]
+\sum_q\mathcal D[C_q]\rho,
\]
由 $\rho_R=U_R^\dagger\rho U_R$ 对时间求导可得
\begin{equation}
\boxed{
\dot\rho_R
=-\frac{\ii}{\hbar}
[H_R(t)+U_R^\dagger H_{\rm LS}U_R,\rho_R]
+\sum_q\mathcal D[C_q^{R}(t)]\rho_R,
\qquad
C_q^{R}(t)=U_R^\dagger C_qU_R.}
\label{eq:phys-oqs-rotating-frame-master}
\end{equation}
若独立跃迁 jump 写成 $C_q=\sqrt{\gamma_q}\,\sigma_{ji}$，则
\[
C_q^{R}(t)
=\ee^{\ii(\chi_j-\chi_i)t}C_q.
\]
由于 $\mathcal D[\ee^{\ii\phi(t)}C_q]=\mathcal D[C_q]$，单个独立跃迁耗散子在这种对角旋转框架下形式不变；对角纯退相干算符甚至本身不变。

但对\eqnrefs{eq:phys-oqs-general-transition-kossakowski} 的交叉耗散，不能逐项丢掉相位。若
\[
F_\alpha^R(t)=\ee^{\ii\theta_\alpha(t)}F_\alpha,
\qquad
\theta_\alpha(t)=(\chi_{j_\alpha}-\chi_{i_\alpha})t,
\]
则交叉项的系数变为
\[
\Gamma_{\alpha\beta}^{(\omega)}
\ee^{\ii[\theta_\beta(t)-\theta_\alpha(t)]}.
\]
只有当同一耗散块内各跃迁在所选旋转框架中具有相同相位速度时，该 Kossakowski 矩阵仍保持时间无关。这一点在近简并 V 型和 $\Lambda$ 型的干涉耗散中尤其重要。

若 $H_{\rm LS}$ 在所选能量基中对角，记
\[
\Delta_i^{\rm LS}=\langle i|H_{\rm LS}|i\rangle,
\qquad
\delta_i=\delta_i^{(0)}+\frac{\Delta_i^{\rm LS}}{\hbar}.
\]
从这里开始，$\delta_i$ 专指已经包含对角 Lamb shift 的有效旋转框架频率，并定义
\begin{equation}
\boxed{\Delta_{ij}=\delta_i-\delta_j.}
\label{eq:phys-oqs-three-detuning}
\end{equation}
于是直接受驱跃迁的有效失谐为
\[
\Delta_{ij}
=\omega_{ij}-\omega_{ij}^{d}
+\frac{\Delta_i^{\rm LS}-\Delta_j^{\rm LS}}{\hbar}.
\]
为使后面的 Optical Bloch 方程只出现这一组有效参数，定义
\begin{equation}
\boxed{
H_{\rm rot}
=\hbar\sum_i\delta_i\sigma_{ii}
-\frac{\hbar}{2}\sum_{E_i>E_j}
\left(\Omega_{ij}\sigma_{ij}+\Omega_{ij}^*\sigma_{ji}\right).}
\label{eq:phys-oqs-three-hrot}
\end{equation}
并把 Rabi 频率延拓到所有有序指标：$\Omega_{ji}\equiv\Omega_{ij}^*$、$\Omega_{ii}=0$。若 $H_{\rm LS}$ 在简并或近简并子空间内含非对角元，则不能把它压缩进三个标量 $\delta_i$；这些矩阵元必须显式保留在旋转框架 Hamiltonian 中，此时下面的 commutator 母式仍成立，但需把 $H_{ij}$ 的相应非对角贡献一并加入。

对一般三能级密度矩阵，把 Hamiltonian 与耗散项逐个投影到矩阵元即可得到闭合方程。

先考虑独立的 population-transfer 通道
\[
C_{i\to j}=\sqrt{\Gamma_{i\to j}}\,\sigma_{ji},
\qquad i\ne j.
\]
对单个 $k\to l$ 通道，耗散子的 population 矩阵元是
\begin{align}
\langle i|\Diss[C_{k\to l}]\rho|i\rangle
&=\Gamma_{k\to l}\delta_{il}\rho_{kk}
-\frac{\Gamma_{k\to l}}2
\langle i|\{\sigma_{kk},\rho\}|i\rangle\\
&=\Gamma_{k\to l}\delta_{il}\rho_{kk}
-\Gamma_{k\to l}\delta_{ik}\rho_{ii}.
\end{align}
对所有通道求和，得到
\begin{equation}
\boxed{
\dot\rho_{ii}\big|_{\rm tr}
=\sum_{k\ne i}\Gamma_{k\to i}\rho_{kk}
-\Gamma_i^{\rm out}\rho_{ii},
\qquad
\Gamma_i^{\rm out}=\sum_{l\ne i}\Gamma_{i\to l}.}
\label{eq:phys-oqs-three-transfer-pop}
\end{equation}
对 $i\ne j$，同一个耗散子的非对角矩阵元为
\begin{align}
\langle i|\Diss[C_{k\to l}]\rho|j\rangle
&=-\frac{\Gamma_{k\to l}}2
\left(\delta_{ik}+\delta_{jk}\right)\rho_{ij}.
\end{align}
所以 population transfer 对 $\rho_{ij}$ 的总衰减是
\[
-\frac12(\Gamma_i^{\rm out}+\Gamma_j^{\rm out})\rho_{ij}.
\]
把\eqnrefs{eq:phys-oqs-general-pure-dephasing-rate} 的零频纯退相干与 population transfer 的寿命展宽相加，定义总 coherence 衰减率
\begin{equation}
\boxed{
\gamma_{ij}
=\frac12(\Gamma_i^{\rm out}+\Gamma_j^{\rm out})
+\gamma_{ij}^{\varphi}.}
\label{eq:phys-oqs-three-gamma}
\end{equation}

再计算相干项。由\eqnrefs{eq:phys-oqs-three-hrot}，对 population 有
\begin{align}
-\frac{\ii}{\hbar}\langle i|[H_{\rm rot},\rho]|i\rangle
&=-\frac{\ii}{\hbar}
\sum_{k\ne i}(H_{ik}\rho_{ki}-\rho_{ik}H_{ki})\\
&=\frac{\ii}{2}\sum_{k\ne i}
\left(\Omega_{ik}\rho_{ki}-\rho_{ik}\Omega_{ki}\right),
\end{align}
其中已经用 $H_{ik}=-\hbar\Omega_{ik}/2$。与\eqnrefs{eq:phys-oqs-three-transfer-pop} 相加得到一般 population 方程
\begin{equation}
\boxed{
\dot\rho_{ii}
=\frac{\ii}{2}\sum_{k\ne i}
\left(\Omega_{ik}\rho_{ki}-\rho_{ik}\Omega_{ki}\right)
+\sum_{k\ne i}\Gamma_{k\to i}\rho_{kk}
-\Gamma_i^{\rm out}\rho_{ii}.}
\label{eq:phys-oqs-three-obe-pop}
\end{equation}

对 coherence $\rho_{ij}$，设 $k$ 是除 $i,j$ 外的第三个能级。把 commutator 的三个中间态逐项写开，
\begin{align}
\langle i|[H_{\rm rot},\rho]|j\rangle
={}&\hbar(\delta_i-\delta_j)\rho_{ij}
-\frac{\hbar}{2}\Omega_{ij}\rho_{jj}
-\frac{\hbar}{2}\Omega_{ik}\rho_{kj}\\
&+\frac{\hbar}{2}\rho_{ii}\Omega_{ij}
+\frac{\hbar}{2}\rho_{ik}\Omega_{kj}.
\end{align}
乘以 $-\ii/\hbar$，再加入\eqnrefs{eq:phys-oqs-three-gamma} 的耗散，得到
\begin{equation}
\boxed{
\begin{aligned}
\dot\rho_{ij}
={}&-(\gamma_{ij}+\ii\Delta_{ij})\rho_{ij}
+\frac{\ii\Omega_{ij}}{2}(\rho_{jj}-\rho_{ii})\\
&+\frac{\ii\Omega_{ik}}{2}\rho_{kj}
-\frac{\ii}{2}\rho_{ik}\Omega_{kj},
\qquad i\ne j,\quad k\ne i,j.
\end{aligned}}
\label{eq:phys-oqs-three-obe-coh}
\end{equation}
\eqnrefs{eq:phys-oqs-three-obe-pop} 与\eqnrefs{eq:phys-oqs-three-obe-coh} 是经典驱动三能级系统在独立跃迁耗散下的普适 Optical Bloch 方程。下面把六个独立方程完整写出。利用 $\Omega_{ji}=\Omega_{ij}^*$ 与 $\rho_{ji}=\rho_{ij}^*$，三个 population 为
\begin{subequations}\label{eq:phys-oqs-three-obe-explicit}
\begin{align}
\dot\rho_{11}
={}&\sum_{k\ne1}\Gamma_{k\to1}\rho_{kk}-\Gamma_1^{\rm out}\rho_{11}
+\frac{\ii}{2}
(\Omega_{12}\rho_{21}-\rho_{12}\Omega_{21}
+\Omega_{13}\rho_{31}-\rho_{13}\Omega_{31}),\label{eq:phys-oqs-three-obe-explicit-pop1}\\
\dot\rho_{22}
={}&\sum_{k\ne2}\Gamma_{k\to2}\rho_{kk}-\Gamma_2^{\rm out}\rho_{22}
+\frac{\ii}{2}
(\Omega_{21}\rho_{12}-\rho_{21}\Omega_{12}
+\Omega_{23}\rho_{32}-\rho_{23}\Omega_{32}),\label{eq:phys-oqs-three-obe-explicit-pop2}\\
\dot\rho_{33}
={}&\sum_{k\ne3}\Gamma_{k\to3}\rho_{kk}-\Gamma_3^{\rm out}\rho_{33}
+\frac{\ii}{2}
(\Omega_{31}\rho_{13}-\rho_{31}\Omega_{13}
+\Omega_{32}\rho_{23}-\rho_{32}\Omega_{23}),\label{eq:phys-oqs-three-obe-explicit-pop3}
\end{align}
三个独立 coherence 为
\begin{align}
\dot\rho_{21}
={}&-(\gamma_{21}+\ii\Delta_{21})\rho_{21}
+\frac{\ii\Omega_{21}}2(\rho_{11}-\rho_{22})
+\frac{\ii\Omega_{23}}2\rho_{31}
-\frac{\ii}{2}\rho_{23}\Omega_{31},\label{eq:phys-oqs-three-obe-explicit-coh21}\\
\dot\rho_{31}
={}&-(\gamma_{31}+\ii\Delta_{31})\rho_{31}
+\frac{\ii\Omega_{31}}2(\rho_{11}-\rho_{33})
+\frac{\ii\Omega_{32}}2\rho_{21}
-\frac{\ii}{2}\rho_{32}\Omega_{21},\label{eq:phys-oqs-three-obe-explicit-coh31}\\
\dot\rho_{32}
={}&-(\gamma_{32}+\ii\Delta_{32})\rho_{32}
+\frac{\ii\Omega_{32}}2(\rho_{22}-\rho_{33})
+\frac{\ii\Omega_{31}}2\rho_{12}
-\frac{\ii}{2}\rho_{31}\Omega_{12}.\label{eq:phys-oqs-three-obe-explicit-coh32}
\end{align}
\end{subequations}
把前三式相加，所有 Hamiltonian 项按每一条边成对抵消；transfer 项也因每个 $k\to i$ 对一个方程是流出、对另一个方程是流入而抵消，因此
\[
\frac{\dd}{\dd t}(\rho_{11}+\rho_{22}+\rho_{33})=0.
\]
这给出了对显式方程的迹守恒检查。三种常见构形的差别现在只需在\eqnrefs{eq:phys-oqs-three-obe-explicit} 中置零不存在的 $\Omega_{ij}$、$\Gamma_{i\to j}$，再根据具体稳态决定零阶 population。

若某些近简并跃迁属于同一个 Bohr 频率块，独立跃迁形式要替换成\secref{sec:phys-oqs-gksl}的一般正半定矩阵 $\boldsymbol\Gamma=(\Gamma_{\mu\nu})$：
\begin{equation}
\boxed{
\mathcal D_{\Gamma}[\rho]
=\sum_{\mu\nu}\Gamma_{\mu\nu}
\left(F_\nu\rho F_\mu^\dagger
-\frac12\{F_\mu^\dagger F_\nu,\rho\}\right).}
\label{eq:phys-oqs-three-cross-damping}
\end{equation}
此时应直接从\eqnrefs{eq:phys-oqs-three-cross-damping} 取矩阵元；V 型的交叉阻尼例子将在下面显式计算。

最后给出统一的弱探针展开。用无量纲参数 $\eta$ 乘在所有弱探针 Rabi 频率上，写
\[
\Liouv(\eta)=\Liouv_0+\eta\Liouv_p+O(\eta^2),
\qquad
\rho_{\rm ss}=\rho^{(0)}+\eta\rho^{(1)}+O(\eta^2).
\]
代入 $\Liouv(\eta)|\rho_{\rm ss}\!\rangle\!\rangle=0$，零阶与一阶分别满足
\[
\Liouv_0|\rho^{(0)}\!\rangle\!\rangle=0,
\]
\begin{equation}
\boxed{
\Liouv_0|\rho^{(1)}\!\rangle\!\rangle
=-\Liouv_p|\rho^{(0)}\!\rangle\!\rangle,
\qquad
\Tr\rho^{(1)}=0.}
\label{eq:phys-oqs-three-linear-response}
\end{equation}
由于 $\Liouv_0$ 必然有稳态零模，它在整个 Liouville 空间上不可逆。若零阶稳态唯一，则在迹为零的子空间上可以使用约化逆或 Drazin 逆 $\Liouv_0^+$，写成
\begin{equation}
\boxed{
|\rho^{(1)}\!\rangle\!\rangle
=-\Liouv_0^+\Liouv_p|\rho^{(0)}\!\rangle\!\rangle.}
\label{eq:phys-oqs-linear-response-pseudoinverse}
\end{equation}
\eqnrefs{eq:phys-oqs-linear-response-pseudoinverse} 是弱探针线性响应的真正普适解；后面把两个 coherence 联立消元，只是这个 Liouville 空间方程在特定构形中的低维实现。线性响应的控制量也应在这里给出。对任意与所选算符范数相容的诱导范数，可以用
\begin{equation}
\boxed{
\epsilon_p
=|\eta|\,\|\Liouv_0^+\Liouv_p\|
\ll1}
\label{eq:phys-oqs-probe-small-parameter}
\end{equation}
作为充分的小参数诊断；它直接控制一阶修正相对于零阶态的尺度。若 $\Liouv_0$ 出现额外近零模，$\|\Liouv_0^+\|$ 会变大，即使探针 Rabi 频率本身很小也可能破坏线性响应。此时必须保留更高阶项或直接求完整非线性稳态。

下面三个例子都只是在\eqnrefs{eq:phys-oqs-three-obe-explicit}、\eqnrefs{eq:phys-oqs-general-transition-kossakowski} 与\eqnrefs{eq:phys-oqs-three-linear-response} 中代入不同的非零参数。

\begin{exbox}[$\Lambda$ 型]
$\Lambda$ 型有两个下能级 $|1\rangle,|2\rangle$ 与一个上能级 $|3\rangle$。单模量子场中可取
\[
(\nu_1,\nu_2,\nu_3)=(0,0,1),
\qquad
g_{31},g_{32}\ne0,
\qquad
g_{21}=0.
\]
由\eqnrefs{eq:phys-oqs-hn3-elements}，完整三维流形 $N\ge1$ 中只有 $1\leftrightarrow3$ 与 $2\leftrightarrow3$ 的非对角元：
\[
\left.H_{S,N}^{(3)}\right|_{\Lambda}=
\begin{pmatrix}
E_1+\hbar\omega_c(N+\tfrac12)&0&\hbar g_{31}^*\sqrt N\\
0&E_2+\hbar\omega_c(N+\tfrac12)&\hbar g_{32}^*\sqrt N\\
\hbar g_{31}\sqrt N&\hbar g_{32}\sqrt N&E_3+\hbar\omega_c(N-\tfrac12)
\end{pmatrix}.
\]
该矩阵由\eqnrefs{eq:phys-oqs-hn3-elements} 在上述 $\nu_i,g_{ij}$ 下逐项代入得到。若只考虑零温自发辐射，则非零原子 jump 为
\[
C_{3\to1}=\sqrt{\Gamma_{3\to1}}\sigma_{13},
\qquad
C_{3\to2}=\sqrt{\Gamma_{3\to2}}\sigma_{23},
\]
所以
\[
\Gamma_1^{\rm out}=\Gamma_2^{\rm out}=0,
\qquad
\Gamma_3^{\rm out}=\Gamma_{3\to1}+\Gamma_{3\to2},
\]
并由\eqnrefs{eq:phys-oqs-three-gamma}
\[
\gamma_{31}=\frac{\Gamma_{3\to1}+\Gamma_{3\to2}}2+\gamma_{31}^{\varphi},
\quad
\gamma_{32}=\frac{\Gamma_{3\to1}+\Gamma_{3\to2}}2+\gamma_{32}^{\varphi},
\quad
\gamma_{21}=\gamma_{21}^{\varphi}.
\]
特别地，在上述“两个自发辐射通道可由环境分辨、$\Gamma_{\alpha\beta}$ 已对角”的条件下，两个下能级的 coherence 不因上能级寿命直接衰减；这正是 $\Lambda$ 型能形成长寿命 Raman coherence 的原因。

若 $3\to1$ 与 $3\to2$ 近简并且跃迁偶极不正交，则必须回到\eqnrefs{eq:phys-oqs-general-transition-kossakowski}。取
\[
F_1=\sigma_{13},\qquad F_2=\sigma_{23},
\qquad \Gamma_{12}\ne0,
\]
由于 $F_1^\dagger F_2=F_2^\dagger F_1=0$，交叉项不增加上能级的额外 loss；但 recycle 项直接给
\begin{equation}
\boxed{
\dot\rho_{21}\big|_{\rm cross}
=\Gamma_{12}\rho_{33},
\qquad
\dot\rho_{12}\big|_{\rm cross}
=\Gamma_{12}^*\rho_{33}.}
\label{eq:phys-oqs-lambda-spontaneous-coherence}
\end{equation}
所以“长寿命下能级 coherence”与“自发产生下能级 coherence”是两个不同效应：前者来自 $\Gamma_1^{\rm out}=\Gamma_2^{\rm out}=0$，后者来自不可分辨跃迁的非对角 Kossakowski 元。下面的标准 $\Lambda$ 型 Optical Bloch 方程先采用独立通道，即 $\Gamma_{12}=0$。

现在考虑经典场，令
\[
\Omega_{31},\Omega_{32}\ne0,
\qquad
\Omega_{21}=0.
\]
把这些条件代入\eqnrefs{eq:phys-oqs-three-obe-pop}，三个 population 方程为
\begin{align}
\dot\rho_{33}
={}&-\Gamma_3^{\rm out}\rho_{33}
+\frac{\ii}{2}
(\Omega_{31}\rho_{13}-\rho_{31}\Omega_{13}
+\Omega_{32}\rho_{23}-\rho_{32}\Omega_{23}),\\
\dot\rho_{11}
={}&\Gamma_{3\to1}\rho_{33}
+\frac{\ii}{2}(\Omega_{13}\rho_{31}-\rho_{13}\Omega_{31}),\\
\dot\rho_{22}
={}&\Gamma_{3\to2}\rho_{33}
+\frac{\ii}{2}(\Omega_{23}\rho_{32}-\rho_{23}\Omega_{32}).
\end{align}
再从\eqnrefs{eq:phys-oqs-three-obe-coh} 分别取 $(i,j)=(3,1),(3,2),(2,1)$，得到
\begin{align}
\dot\rho_{31}
={}&-(\gamma_{31}+\ii\Delta_{31})\rho_{31}
+\frac{\ii\Omega_{31}}2(\rho_{11}-\rho_{33})
+\frac{\ii\Omega_{32}}2\rho_{21},\\
\dot\rho_{32}
={}&-(\gamma_{32}+\ii\Delta_{32})\rho_{32}
+\frac{\ii\Omega_{32}}2(\rho_{22}-\rho_{33})
+\frac{\ii\Omega_{31}}2\rho_{12},\\
\dot\rho_{21}
={}&-(\gamma_{21}+\ii\Delta_{21})\rho_{21}
+\frac{\ii\Omega_{32}^*}{2}\rho_{31}
-\frac{\ii}{2}\rho_{23}\Omega_{31}.
\end{align}
这里
\[
\Delta_{21}=\delta_2-\delta_1=\Delta_{31}-\Delta_{32},
\]
所以两光子共振就是 $\Delta_{31}=\Delta_{32}$。

若 $\Omega_{31}$ 是弱探针、$\Omega_{32}$ 是有限控制场，而且在探针加入前稳态主要为 $\rho_{11}^{(0)}=1$，其余 population 与 coherence 为零，则 $\rho_{31}^{(1)},\rho_{21}^{(1)}=O(\Omega_{31})$，而 $\rho_{23}^{(1)}\Omega_{31}=O(\Omega_{31}^2)$ 可在一阶忽略。上面两条 coherence 方程化为
\[
0=-(\gamma_{31}+\ii\Delta_{31})\rho_{31}^{(1)}
+\frac{\ii\Omega_{31}}2
+\frac{\ii\Omega_{32}}2\rho_{21}^{(1)},
\]
\[
0=-(\gamma_{21}+\ii\Delta_{21})\rho_{21}^{(1)}
+\frac{\ii\Omega_{32}^*}{2}\rho_{31}^{(1)}.
\]
第二式给
\[
\rho_{21}^{(1)}
=\frac{\ii\Omega_{32}^*/2}{\gamma_{21}+\ii\Delta_{21}}\rho_{31}^{(1)},
\]
代回第一式得到
\begin{equation}
\boxed{
\rho_{31}^{(1)}
=\frac{(\ii\Omega_{31}/2)(\gamma_{21}+\ii\Delta_{21})}
{(\gamma_{31}+\ii\Delta_{31})(\gamma_{21}+\ii\Delta_{21})
+|\Omega_{32}|^2/4}.}
\label{eq:phys-oqs-lambda-response}
\end{equation}
因此吸收谱不是从某个额外公式得到，而是一般 coherence 方程的二元线性方程组直接求解。

在两光子共振且下能级近简并时，去掉共同对角能量后相干耦合部分为
\[
H_{\rm int}^{(\Lambda)}
=-\frac{\hbar}{2}
(\Omega_{31}\sigma_{31}+\Omega_{32}\sigma_{32}+\mathrm{h.c.}).
\]
作用在
\[
|D\rangle
=\frac{\Omega_{32}|1\rangle-\Omega_{31}|2\rangle}
{\sqrt{|\Omega_{31}|^2+|\Omega_{32}|^2}}
\]
上时，$|3\rangle$ 的系数正比于
$\Omega_{31}\Omega_{32}-\Omega_{32}\Omega_{31}=0$，故 $H_{\rm int}^{(\Lambda)}|D\rangle$ 不含激发态分量。这就是暗态；\eqnrefs{eq:phys-oqs-lambda-response} 中两光子共振处的吸收抑制是同一干涉结构在密度矩阵线性响应中的表现。
\end{exbox}

\clearpage
\begin{exbox}[V 型]
V 型有一个下能级 $|1\rangle$ 与两个上能级 $|2\rangle,|3\rangle$。单模量子场中取
\[
(\nu_1,\nu_2,\nu_3)=(0,1,1),
\qquad
g_{21},g_{31}\ne0,
\qquad
g_{32}=0.
\]
代入\eqnrefs{eq:phys-oqs-hn3-elements}，$N\ge1$ 时
\[
\left.H_{S,N}^{(3)}\right|_{V}=
\begin{pmatrix}
E_1+\hbar\omega_c(N+\tfrac12)&\hbar g_{21}^*\sqrt N&\hbar g_{31}^*\sqrt N\\
\hbar g_{21}\sqrt N&E_2+\hbar\omega_c(N-\tfrac12)&0\\
\hbar g_{31}\sqrt N&0&E_3+\hbar\omega_c(N-\tfrac12)
\end{pmatrix}.
\]
若两个下降通道 $2\to1$ 与 $3\to1$ 在环境看来可分辨，则只需在\eqnrefs{eq:phys-oqs-three-obe-pop}--\eqnrefs{eq:phys-oqs-three-obe-coh} 中取
\[
\Gamma_{2\to1},\Gamma_{3\to1}\ne0,
\qquad
\Omega_{21},\Omega_{31}\ne0,
\qquad
\Omega_{32}=0.
\]
于是
\[
\Gamma_2^{\rm out}=\Gamma_{2\to1},
\qquad
\Gamma_3^{\rm out}=\Gamma_{3\to1},
\qquad
\Gamma_1^{\rm out}=0,
\]
并且上能级 coherence 的自然寿命贡献是
\[
\gamma_{23}\big|_{\rm tr}
=\frac12(\Gamma_{2\to1}+\Gamma_{3\to1}).
\]
把 $\Omega_{32}=0$ 和上述跃迁率代入\eqnrefs{eq:phys-oqs-three-obe-pop}，得到
\begin{align}
\dot\rho_{11}
={}&\Gamma_{2\to1}\rho_{22}+\Gamma_{3\to1}\rho_{33}
+\frac{\ii}{2}
(\Omega_{12}\rho_{21}-\rho_{12}\Omega_{21}
+\Omega_{13}\rho_{31}-\rho_{13}\Omega_{31}),\\
\dot\rho_{22}
={}&-\Gamma_{2\to1}\rho_{22}
+\frac{\ii}{2}(\Omega_{21}\rho_{12}-\rho_{21}\Omega_{12}),\\
\dot\rho_{33}
={}&-\Gamma_{3\to1}\rho_{33}
+\frac{\ii}{2}(\Omega_{31}\rho_{13}-\rho_{31}\Omega_{13}).
\end{align}
相应地，由\eqnrefs{eq:phys-oqs-three-obe-coh}
\begin{align}
\dot\rho_{21}
={}&-(\gamma_{21}+\ii\Delta_{21})\rho_{21}
+\frac{\ii\Omega_{21}}2(\rho_{11}-\rho_{22})
-\frac{\ii}{2}\rho_{23}\Omega_{31},\\
\dot\rho_{31}
={}&-(\gamma_{31}+\ii\Delta_{31})\rho_{31}
+\frac{\ii\Omega_{31}}2(\rho_{11}-\rho_{33})
-\frac{\ii}{2}\rho_{32}\Omega_{21},\\
\dot\rho_{23}
={}&-(\gamma_{23}+\ii\Delta_{23})\rho_{23}
+\frac{\ii\Omega_{21}}2\rho_{13}
-\frac{\ii}{2}\rho_{21}\Omega_{13}.
\end{align}
因此 V 型中两条光学 coherence $\rho_{21},\rho_{31}$ 通过上能级 coherence $\rho_{23}$ 相互联系。若其中一条驱动作为强场，弱探针展开的零阶态应先由保留该强场的 $\Liouv_0\rho^{(0)}=0$ 求出；一般不能预先把 $\rho_{11}^{(0)}$ 固定为 1，因为强场本身会重新分配 $|1\rangle$ 与相应上能级的 population。

V 型的另一重要情形是两个上能级近简并、两条偶极矩共享同一环境模式。此时不能使用相互独立的两个标量耗散子，而应从\eqnrefs{eq:phys-oqs-three-cross-damping} 出发，取
\[
F_2=\sigma_{12},\qquad F_3=\sigma_{13},
\qquad
\boldsymbol\Gamma=
\begin{pmatrix}
\Gamma_{2\to1}&\Gamma_{23}\\
\Gamma_{23}^*&\Gamma_{3\to1}
\end{pmatrix}\succeq0.
\]
现在直接取矩阵元。jump 项只回注到 $|1\rangle$，所以
\begin{align}
\dot\rho_{11}\big|_{\Gamma}
={}&\Gamma_{2\to1}\rho_{22}+\Gamma_{3\to1}\rho_{33}
+\Gamma_{23}\rho_{32}+\Gamma_{23}^*\rho_{23}.
\end{align}
反对易子作用在上能级子空间，给
\begin{align}
\dot\rho_{22}\big|_{\Gamma}
={}&-\Gamma_{2\to1}\rho_{22}
-\frac12(\Gamma_{23}\rho_{32}+\Gamma_{23}^*\rho_{23}),\\
\dot\rho_{33}\big|_{\Gamma}
={}&-\Gamma_{3\to1}\rho_{33}
-\frac12(\Gamma_{23}\rho_{32}+\Gamma_{23}^*\rho_{23}),\\
\dot\rho_{23}\big|_{\Gamma}
={}&-\frac12(\Gamma_{2\to1}+\Gamma_{3\to1})\rho_{23}
-\frac{\Gamma_{23}}2(\rho_{22}+\rho_{33}),\\
\dot\rho_{21}\big|_{\Gamma}
={}&-\frac{\Gamma_{2\to1}}2\rho_{21}
-\frac{\Gamma_{23}}2\rho_{31},\\
\dot\rho_{31}\big|_{\Gamma}
={}&-\frac{\Gamma_{3\to1}}2\rho_{31}
-\frac{\Gamma_{23}^*}{2}\rho_{21}.
\end{align}
这些交叉项均由同一个 $\Gamma_{\mu\nu}$ 矩阵的非对角元逐矩阵元展开得到。正半定性给出
\[
|\Gamma_{23}|^2\le
\Gamma_{2\to1}\Gamma_{3\to1}.
\]
耗散矩阵的两个本征值为
\[
\gamma_{\pm}
=\frac12\left[
\Gamma_{2\to1}+\Gamma_{3\to1}
\pm
\sqrt{(\Gamma_{2\to1}-\Gamma_{3\to1})^2+4|\Gamma_{23}|^2}
\right].
\]
当等号 $|\Gamma_{23}|^2=\Gamma_{2\to1}\Gamma_{3\to1}$ 成立时，$\gamma_-=0$。若写
$\Gamma_{23}=\ee^{\ii\phi}\sqrt{\Gamma_{2\to1}\Gamma_{3\to1}}$，并假设
$\Gamma_{2\to1}+\Gamma_{3\to1}>0$，则归一化的零衰减方向为
\[
|D_{\Gamma}\rangle
=\frac{
\sqrt{\Gamma_{3\to1}}|2\rangle
-\ee^{-\ii\phi}\sqrt{\Gamma_{2\to1}}|3\rangle}
{\sqrt{\Gamma_{2\to1}+\Gamma_{3\to1}}}.
\]
它是否同时成为完整动力学的暗态，还要继续把它代入\eqnrefs{eq:phys-oqs-three-hrot} 检查相干驱动是否把它耦合到 $|1\rangle$；因此“耗散暗态”和“Hamiltonian 暗态”是两个需要同时验证的条件。
\end{exbox}

\clearpage
\begin{exbox}[$\Xi$ 型]
$\Xi$ 型是级联结构 $|1\rangle\leftrightarrow|2\rangle\leftrightarrow|3\rangle$。单模量子场中取
\[
(\nu_1,\nu_2,\nu_3)=(0,1,2),
\qquad
g_{21},g_{32}\ne0,
\qquad
g_{31}=0.
\]
代入\eqnrefs{eq:phys-oqs-hn3-elements}，对 $N\ge2$ 的完整三维流形有
\[
\left.H_{S,N}^{(3)}\right|_{\Xi}=
\begin{pmatrix}
E_1+\hbar\omega_c(N+\tfrac12)&\hbar g_{21}^*\sqrt N&0\\
\hbar g_{21}\sqrt N&E_2+\hbar\omega_c(N-\tfrac12)&\hbar g_{32}^*\sqrt{N-1}\\
0&\hbar g_{32}\sqrt{N-1}&E_3+\hbar\omega_c(N-\tfrac32)
\end{pmatrix}.
\]
这里第二条边的系数是 $\sqrt{N-1}$，直接来自\eqnrefs{eq:phys-oqs-hn3-elements} 中 $\sqrt{N-\nu_2}$；这也是为什么级联结构不能简单照抄二能级的 $\sqrt N$。

在零温级联衰减中取
\[
C_{3\to2}=\sqrt{\Gamma_{3\to2}}\sigma_{23},
\qquad
C_{2\to1}=\sqrt{\Gamma_{2\to1}}\sigma_{12}.
\]
于是
\[
\Gamma_3^{\rm out}=\Gamma_{3\to2},
\qquad
\Gamma_2^{\rm out}=\Gamma_{2\to1},
\qquad
\Gamma_1^{\rm out}=0,
\]
从\eqnrefs{eq:phys-oqs-three-gamma} 得
\[
\gamma_{21}=\frac{\Gamma_{2\to1}}2+\gamma_{21}^{\varphi},
\quad
\gamma_{32}=\frac{\Gamma_{3\to2}+\Gamma_{2\to1}}2+\gamma_{32}^{\varphi},
\quad
\gamma_{31}=\frac{\Gamma_{3\to2}}2+\gamma_{31}^{\varphi}.
\]
注意 $\rho_{31}$ 虽然对应一个没有直接偶极驱动的两光子 coherence，但它的寿命仍由 $|3\rangle$ 与 $|1\rangle$ 的总出射率按\eqnrefs{eq:phys-oqs-three-gamma} 决定。

经典驱动时取
\[
\Omega_{21},\Omega_{32}\ne0,
\qquad
\Omega_{31}=0.
\]
由\eqnrefs{eq:phys-oqs-three-obe-pop} 得
\begin{align}
\dot\rho_{11}
={}&\Gamma_{2\to1}\rho_{22}
+\frac{\ii}{2}(\Omega_{12}\rho_{21}-\rho_{12}\Omega_{21}),\\
\dot\rho_{22}
={}&\Gamma_{3\to2}\rho_{33}-\Gamma_{2\to1}\rho_{22}
+\frac{\ii}{2}
(\Omega_{21}\rho_{12}-\rho_{21}\Omega_{12}
+\Omega_{23}\rho_{32}-\rho_{23}\Omega_{32}),\\
\dot\rho_{33}
={}&-\Gamma_{3\to2}\rho_{33}
+\frac{\ii}{2}(\Omega_{32}\rho_{23}-\rho_{32}\Omega_{23}).
\end{align}
三个独立 coherence 方程由\eqnrefs{eq:phys-oqs-three-obe-coh} 直接得到
\begin{align}
\dot\rho_{21}
={}&-(\gamma_{21}+\ii\Delta_{21})\rho_{21}
+\frac{\ii\Omega_{21}}2(\rho_{11}-\rho_{22})
+\frac{\ii\Omega_{32}^*}{2}\rho_{31},\\
\dot\rho_{32}
={}&-(\gamma_{32}+\ii\Delta_{32})\rho_{32}
+\frac{\ii\Omega_{32}}2(\rho_{22}-\rho_{33})
-\frac{\ii}{2}\rho_{31}\Omega_{21}^*,\\
\dot\rho_{31}
={}&-(\gamma_{31}+\ii\Delta_{31})\rho_{31}
+\frac{\ii\Omega_{32}}2\rho_{21}
-\frac{\ii}{2}\rho_{32}\Omega_{21}.
\end{align}
由\eqnrefs{eq:phys-oqs-three-detuning} 自动有
\[
\Delta_{31}=\delta_3-\delta_1
=(\delta_3-\delta_2)+(\delta_2-\delta_1)
=\Delta_{32}+\Delta_{21},
\]
所以这里的两光子失谐不需要再引入新的符号。

令 $\Omega_{21}$ 为弱探针，$\Omega_{32}$ 为有限控制场，并取无探针稳态 $\rho_{11}^{(0)}=1$。一阶中 $\rho_{21}^{(1)},\rho_{31}^{(1)}=O(\Omega_{21})$，而 $\rho_{32}^{(1)}\Omega_{21}=O(\Omega_{21}^2)$，故
\[
0=-(\gamma_{21}+\ii\Delta_{21})\rho_{21}^{(1)}
+\frac{\ii\Omega_{21}}2
+\frac{\ii\Omega_{32}^*}{2}\rho_{31}^{(1)},
\]
\[
0=-(\gamma_{31}+\ii\Delta_{31})\rho_{31}^{(1)}
+\frac{\ii\Omega_{32}}2\rho_{21}^{(1)}.
\]
第二式给
\[
\rho_{31}^{(1)}
=\frac{\ii\Omega_{32}/2}{\gamma_{31}+\ii\Delta_{31}}\rho_{21}^{(1)},
\]
代回第一式，得到
\begin{equation}
\boxed{
\rho_{21}^{(1)}
=\frac{(\ii\Omega_{21}/2)(\gamma_{31}+\ii\Delta_{31})}
{(\gamma_{21}+\ii\Delta_{21})(\gamma_{31}+\ii\Delta_{31})
+|\Omega_{32}|^2/4}.}
\label{eq:phys-oqs-xi-response}
\end{equation}
控制场通过两光子 coherence $\rho_{31}$ 改写探针极化 $\rho_{21}$ 的分母。响应表现为干涉透明还是可分辨的 dressed-state 双峰，由 $|\Omega_{32}|$ 与 $\gamma_{21},\gamma_{31}$ 的相对尺度决定，不能仅凭极点分裂的代数形式将两种机制等同。
\end{exbox}

接下来考察应用近似与数值截断的控制。

具体腔与多能级模型还引入系统 RWA、经典场 optical RWA、远失谐频率筛选、裸/缀饰耗散、弱探针和 Fock 截断；它们彼此独立，也独立于系统--环境 Born、Markov 与 secular 近似。系统 RWA 由\eqnrefs{eq:phys-oqs-rwa-small-parameter} 控制，失效时回到 $H_S^{\rm full}$ 并保留 Bloch--Siegert 位移、虚光子成分与激发数不守恒；经典场的反旋项与远失谐旋转项分别由\eqnrefs{eq:phys-oqs-classical-optical-rwa-control}、\eqnrefs{eq:phys-oqs-classical-offresonant-control} 控制，前者失效时需保留 counter-rotating dynamics，后者失效时需保留对应驱动及其 AC-Stark 修正。裸耗散只有在\eqnrefs{eq:phys-oqs-local-dissipator-condition} 的环境谱不可分辨与粗粒化不可完全分辨条件同时成立时才受控，否则必须用实际耦合 Hamiltonian 的 dressed Bohr 算符重建耗散块。

多频驱动还要求旋转框架一致：只有沿任意闭环的驱动相位和为零，才存在一个对角 $U_R(t)$ 同时消去全部显式相位；否则剩余拍频使 Liouvillian 成为周期或准周期时间依赖，周期情形可用 Floquet--Liouville 方法。弱探针则由\eqnrefs{eq:phys-oqs-probe-small-parameter} 中的 $\|\Liouv_0^+\|$ 直接控制；靠近额外慢模或耗散临界点时必须检查实际的一阶态修正，而不能只比较 Rabi 频率与某个裸衰减率。

Fock 截断只是数值截断。定义顶层占据
\begin{equation}
P_{\rm edge}
=\sum_i\sum_{n\in\text{最高若干层}}
\langle i,n|\rho|i,n\rangle,
\label{eq:phys-oqs-fock-edge-population}
\end{equation}
应令 $P_{\rm edge}$ 远小于目标数值精度，并递增 $n_{\max}$ 检查目标可观测量收敛；高温、强驱动或接近激光阈值时尤其需要这一检查。自由空间有限能级偶极模型的 Lamb-shift 主值积分还可能具有高频敏感性；绝对 Lamb shift 应由足以描述高频结构的原子/QED 模型给出，并把可测跃迁频率理解为重整化后的参数。

\begin{center}
\small
\begin{tabularx}{0.96\textwidth}{>{\raggedright\arraybackslash\bfseries}p{2.6cm} X X}
\toprule
条件 & 接近失效时的信号 & 更完整描述\\
\midrule
系统 RWA & Bloch--Siegert 位移、激发数不守恒 & $H_S^{\rm full}$、Rabi/Floquet 模型\\
经典场 optical RWA & counter-rotating 微运动、Bloch--Siegert 位移 & 保留反旋驱动项\\
远失谐频率筛选 & AC-Stark 位移或远失谐激发不可忽略 & 保留对应旋转驱动项\\
Born 截断 & 高阶累积量、bath back-action 不可忽略 & 高阶 NZ/TCL、扩大系统\\
Markov & 回波、长尾相关、非指数记忆 & 有限记忆 NZ/TCL、扩大系统\\
Complete secular & 近简并跃迁干涉 & finite coarse graining / partial secular\\
裸耗散 & bath 可分辨 dressed splitting & dressed $S_\mu(\omega)$\\
弱探针 & 饱和、population 大幅改变 & 完整非线性 Liouvillian\\
Fock 截断 & 顶层 population 不可忽略 & 增大 $n_{\max}$ 并做收敛检查\\
\bottomrule
\end{tabularx}
\end{center}

\begin{derivation}[量子 Zeno 效应与脉冲控制中的有效 Hamiltonian]
量子 Zeno 效应描述频繁测量对量子演化的抑制，其数学结构是脉冲控制与有效 Hamiltonian 的极限。分四步建立。

\emph{第一步：Zeno 效应的精确表述。}对初始态 \(|\psi_0\rangle\) 在投影 \(P=|\psi_0\rangle\langle\psi_0|\) 上，\(N\) 次等间隔投影测量（间隔 \(\tau=T/N\)）后的存活概率
\begin{equation*}
  P_{\rm surv}(T,N)=\left|\langle\psi_0|(P\ee^{-\ii H\tau/\hbar}P)^N|\psi_0\rangle\right|^2=\left|\langle\psi_0|\ee^{-\ii H\tau/\hbar}|\psi_0\rangle\right|^{2N}.
\end{equation*}
短时展开 \(\langle\psi_0|\ee^{-\ii H\tau/\hbar}|\psi_0\rangle=1-\frac{\ii\tau}{\hbar}\langle H\rangle-\frac{\tau^2}{2\hbar^2}\langle H^2\rangle+\cdots\)，故
\begin{equation*}
  |\langle\psi_0|\ee^{-\ii H\tau/\hbar}|\psi_0\rangle|^2=1-\frac{\tau^2}{\hbar^2}\Delta H^2+O(\tau^3),
\end{equation*}
其中 \(\Delta H^2=\langle H^2\rangle-\langle H\rangle^2\)。因此
\begin{equation*}
  P_{\rm surv}(T,N)=\left(1-\frac{T^2\Delta H^2}{N^2\hbar^2}+\cdots\right)^N\to1\quad(N\to\infty),
\end{equation*}
即无穷频繁测量完全冻结演化（Zeno 效应）。

\emph{第二步：Zeno 有效 Hamiltonian。}更精确地，\(P\ee^{-\ii H\tau/\hbar}P\) 在 \(P\) 的支撑上可写为 \(\ee^{-\ii H_Z\tau/\hbar}\)，其中 Zeno 有效 Hamiltonian
\begin{equation*}
  H_Z=PHP+\frac{1}{2}\sum_{\alpha}\frac{PHQ_\alpha Q_\alpha HP}{\Delta E_\alpha}+O(\tau),
\end{equation*}
\(Q_\alpha=|\phi_\alpha\rangle\langle\phi_\alpha|\) 是 \(H\) 在 \(P\) 补空间中的本征投影，\(\Delta E_\alpha\) 是能级差。首项 \(PHP\) 是投影后的裸 Hamiltonian（Zeno 子空间动力学），次项是虚跃迁修正（类似二阶微扰）。当 \(PHP\) 在 \(P\) 支撑上非简并时，Zeno 效应完全冻结；当 \(PHP\) 有非平凡谱时，Zeno 子空间内仍有动力学，由 \(PHP\) 生成。

\emph{第三步：反 Zeno 效应。}当测量间隔 \(\tau\) 不趋于零而取有限值时，存活概率可能小于无测量值——频繁测量反而加速衰变（反 Zeno）。对不稳定态，衰变率 \(\Gamma(\tau)=\frac{1-\langle\psi_0|\ee^{-\ii H\tau/\hbar}|\psi_0\rangle|^2}{\tau}\) 在 \(\tau\to0\) 时 \(\Gamma\to0\)（Zeno），但在有限 \(\tau^*\) 处 \(\Gamma(\tau^*)\) 可能超过自然衰变率 \(\Gamma_{\rm nat}\)。\(\tau^*\) 由能谱密度 \(g(E)\) 决定：若 \(g(E)\) 在 \(E_0=\langle H\rangle\) 处有峰且宽度 \(\sigma_E\)，则 \(\tau^*\sim\hbar/\sigma_E\)，此时 \(\Gamma(\tau^*)\sim\sigma_E/\hbar\) 可远大于 \(\Gamma_{\rm nat}\)。这把 Zeno/反 Zeno 的分界与能谱结构精确联系。

\emph{第四步：脉冲控制与动力学解耦。}Zeno 效应是脉冲控制的连续极限。周期施加投影 \(P\) 等价于脉冲序列 \(U(T)=\prod_k P\ee^{-\ii H\tau/\hbar}\)，其有效 Hamiltonian \(H_{\rm eff}=\frac{\ii\hbar}{T}\log U(T)\) 由 Magnus 展开给出。对自旋 echo（施加 \(\pi\)-脉冲翻转），\(H_{\rm eff}\) 消去奇数阶环境耦合（\([H_{\rm env},\pi_x]=0\) 的项保留，其余消去），是动力学解耦的原理。Uhrig 动力学解耦（UDD）优化脉冲时间 \(\{t_k\}\) 使 \(H_{\rm eff}\) 到 \(O(T^{N+1})\) 恒为零，脉冲间隔 \(\delta t_k=T\sin^2\frac{k\pi}{2(N+1)}\)（Uhrig 2007）。这把 Zeno 效应从"投影测量"推广到"酉脉冲控制"，是量子纠错和噪声抑制的核心技术。

\emph{物理意义。}Zeno 效应的数学本质是投影 \(P\) 对演化的截断：\(P\ee^{-\ii H\tau}P\approx\ee^{-\ii PHP\tau}\) 在 \(\tau\to0\) 极限精确，有限 \(\tau\) 时有二阶修正。反 Zeno 效应则源于有限 \(\tau\) 时衰变率 \(\Gamma(\tau)\) 的非单调性——能谱结构决定 \(\Gamma\) 在何处取极大。在量子计算中，UDD 脉冲序列利用这一结构实现高阶噪声抑制，其脉冲间隔的 \(\sin^2\) 分布是数论优化（Chebyshev 多式）的结果。

\emph{量子纠错中的 Zeno 子空间。}在主动量子纠错中，频繁的纠错投影使系统动力学限制在纠错码子空间（Zeno 子空间），由 \(PHP\) 生成有效 Hamiltonian。当纠错码距离 \(d\) 足于无穷时，Zeno 极限使逻辑量子比特的退相干率趋于零——这是容错量子计算的物理基础。UDD 脉冲与纠错投影的组合（"动力学纠错"）在实验上已用于超导量子比特和离子阱系统，实现 \(10^2\)--\(10^3\) 倍的退相干抑制，验证了 Zeno 有效 Hamiltonian 理论的定量预言。











\emph{[Lindblad 方程的完全正定性与 GKLS 表示]} Lindblad 主方程的完全正定性保证物理密度矩阵在演化下保持正定，其代数结构由 Gorini--Kossakowski--Sudarshan (GKS) 表示给出。分四步建立。

\emph{第一步：完全正映射的定义。}量子信道 \(\Phi\) 是完全正的（CP）：对任意辅助空间 \(\mathcal H_A\)，\(\Phi\otimes\mathrm{Id}_A\) 正。CP 条件强于正性——存在正映射非 CP（转置映射 \(\Theta\) 正但 \(\Theta\otimes\mathrm{Id}\) 给部分转置，可负）。CP 映射的 Kraus 表示：\(\Phi(\rho)=\sum_k V_k\rho V_k^\dagger\)（\(V_k\) 任意算符），是 CP 性的代数特征。

\emph{第二步：GKLS 表示。}保迹 CP 映射的生成元（Lindblad 方程）由 GKS 定理给出：
\begin{equation}
  \frac{\dd\rho}{\dd t}=-\ii[H,\rho]+\sum_{\alpha,\beta}c_{\alpha\beta}\left(F_\alpha\rho F_\beta^\dagger-\frac{1}{2}\{F_\beta^\dagger F_\alpha,\rho\}\right),
  \label{eq:phys-oqs-gkls}
\end{equation}
其中 \(H=H^\dagger\) 是有效 Hamiltonian，\(\{F_\alpha\}\) 是算符基，\(C=(c_{\alpha\beta})\) 是半正定矩阵（Kossakowski 矩阵）。\(C\ge0\) 保证\eqnrefs{eq:phys-oqs-gkls} 是完全正的——这是 CP 性的代数判据。对 \(C\) 做 Cholesky 分解 \(C=WW^\dagger\)，定义 \(L_k=\sum_\alpha W_{k\alpha}F_\alpha\)，得到标准 Lindblad 形式 \(\sum_k(L_k\rho L_k^\dagger-\frac{1}{2}\{L_k^\dagger L_k,\rho\})\)。

\emph{第三步：保迹性与熵单调性。}Lindblad 方程自动保迹：\(\operatorname{tr}(\dot\rho)=0\)（由 \(\operatorname{tr}[H,\rho]=0\) 和 \(\operatorname{tr}(L\rho L^\dagger-\frac{1}{2}\{L^\dagger L,\rho\})=0\)）。更深刻地，Lindblad 演化满足量子相对熵单调性：\(D(\Phi_t(\rho)\|\Phi_t(\sigma))\le D(\rho\|\sigma)\)（\(D(\rho\|\sigma)=\operatorname{tr}(\rho\log\rho-\rho\log\sigma)\)）。这是 CP 映射的 data processing inequality，保证 Lindblad 演化不可逆（信息只减不增）。

\emph{第四步：稳态与耗散结构。}Lindblad 方程的稳态 \(\rho_{\rm ss}\)（\(\dot\rho_{\rm ss}=0\)）由 Hamiltonian 和耗散算符 \(\{L_k\}\) 的代数关系决定。若 \(\{L_k\}\) 生成整个算符代数（遍历条件），稳态唯一。有限维系统中稳态总是存在（Brouwer 不动点），但唯一性需要遍历性。稳态的纯度由耗散-相干竞争决定：纯耗散（\(H=0\)）稳态通常是混合态（最大混合 \(\rho\propto I\) 当 \(L_k\) 遍历），相干驱动（\(H\neq0\)）可维持纯态稳态（如驱动-耗散方程的相干态稳态）。这一代数结构把开放量子系统的渐近行为完全编码在 \(\{H,L_k\}\) 的代数中。
\end{derivation}
